% !TeX program = lualatex

% Generated by the object-oriented LaTeX project merger.

\documentclass{book}

\usepackage{titlesec}
\usepackage[table]{xcolor}
\usepackage{array}
\usepackage[colorinlistoftodos]{todonotes}
\usepackage{amsmath}
\usepackage{tikz}
\usepackage{graphicx}
\usepackage{svg}
\usepackage{pdfpages}
\usepackage{dirtree}
\usepackage{booktabs}
\usepackage{adjustbox}
\usepackage[edges]{forest}
\usepackage{float}
\usepackage{caption}
\usepackage{minted}
\usepackage{amssymb}
\usepackage{multirow}
\usepackage{xstring}
\usepackage[most]{tcolorbox}
\usepackage{xparse}
\usepackage{natbib}
\usepackage[a4paper,margin=2.5cm]{geometry}
\usepackage{tabularx}
\usepackage{longtable}
\usepackage{enumitem}
\usepackage{algorithm}
\usepackage{algpseudocode}
\usepackage[unicode,colorlinks=true,linkcolor=blue,urlcolor=blue]{hyperref}
\usepackage{bookmark}

\usetikzlibrary{graphs,graphdrawing,arrows.meta,mindmap,fit,positioning,calc,shapes.geometric,backgrounds,patterns,decorations.pathreplacing}

\usegdlibrary{layered,force}

\newcommand{\projname}{robotic_learning_project}

\newcommand{\figspath}{/home/donkarlo/Dropbox/repo/\projname/figs/}

% Convenience macro

\newcommand{\filepath}[1]{\href{file://#1}{\path{#1}}}

%inline python with minted

\newcommand{\pyl}[1]{\mintinline{python}{#1}}

% graphics

\graphicspath{{/home/donkarlo/Dropbox/repo/\projname/figs/}}

\captionsetup{singlelinecheck=false}

% Minted (needs -shell-escape)

\usemintedstyle{friendly}

\setminted{
    fontsize=\small,
    linenos,
    breaklines,
    frame=single,
    tabsize=4,
    autogobble,
}

\setminted[python]{fontsize=\footnotesize}

% Colors and tikz styles (can stay here)

\definecolor{mainfill}{HTML}{FFF2CC}

\definecolor{mainborder}{HTML}{D6B656}

\tikzset{
    mainnode/.style={
        draw=mainborder,
        fill=mainfill,
        rectangle,
        align=center,
        inner sep=5pt,
        rounded corners,
        line width=0.8pt,
        text=black,
    }
}

% Depths

\setcounter{tocdepth}{4}

\setcounter{secnumdepth}{4}

% ---- Load hyperref LAST (then bookmark, optional) ----

\hypersetup{
    colorlinks=true,
    linkcolor=black,
    citecolor=black,
    urlcolor=black,
    pdfstartview=FitH,
    bookmarksopen=true,
    bookmarksnumbered=true
}
% bookmarks

\definecolor{TagBackground}{HTML}{DCEEFF}

\definecolor{TagText}{HTML}{1E5FAF}

\ExplSyntaxOn

\cs_new_protected:Npn \semantic_tags_print:n #1
    {
    \clist_map_inline:nn {#1}
    {
        \SemanticTagBox{\strut\tl_trim_spaces:n {##1}}
        \hspace{0.35em}
    }
}

\NewDocumentCommand{\semtag}{m}
{
    \semantic_tags_print:n {#1}
}

\ExplSyntaxOff

\newcommand{\seehere}[1]{\href{#1}{see here}}

\hypersetup{
    colorlinks=true,
    linkcolor=blue,
    urlcolor=blue,
    citecolor=blue,
    pdftitle={Surprise in Active Inference: Causes and Modulators}
}

\titleformat{\section}{\large\bfseries}{\thesection}{0.6em}{}

\titleformat{\subsection}{\normalsize\bfseries}{\thesubsection}{0.6em}{}

\setlength{\parindent}{0pt}

\setlength{\parskip}{0.6em}

\renewcommand{\arraystretch}{1.25}

% Compatibility commands added by the merger.
\providecommand{\ProjectDeepHeading}[2]{%
    \par\medskip\noindent\textbf{#2}\par\smallskip
}
\providecommand{\ProjectMissingAsset}[1]{%
    \par\medskip\noindent\fbox{\parbox{0.9\linewidth}{Missing asset: \texttt{#1}}}\par\medskip
}
\providecommand{\ProjectSafeIncludeGraphics}[2][]{\includegraphics[#1]{#2}}
\providecommand{\ProjectSafeIncludePdf}[2][]{\includepdf[#1]{#2}}
\providecommand{\ProjectSafeIncludeSvg}[2][]{\includesvg[#1]{#2}}
\providecommand{\seehere}[1]{\href{#1}{see here}}
\providecommand{\seeherefooter}[1]{\footnote{\url{#1}}}
\providecommand{\clickurlfooter}[1]{\footnote{\url{#1}}}
\providecommand{\seefootnotehere}[1]{\footnote{\url{#1}}}
\providecommand{\semtag}[1]{\texttt{#1}}
\providecommand{\srcpath}[1]{#1}
\providecommand{\figspath}{}
\providecommand{\assetspath}{}
\providecommand{\figspathP}{}
\providecommand{\State}[1]{\ensuremath{#1}}


\begin{document}

% Source: nd_robotic_ai_learning.tex

%table of contents
\tableofcontents
\cite{busoniu-2008-a-comprehensive-survey-of-multi-agent-reinforcement-learning}












% List todos
\listoftodos
% Bibliography

 % Reference the BibTeX file here

\chapter{Robot}

\section{Kind}

% Source: robot/kind/robot_kinds.tex

\subsubsection{Robots Types}
\paragraph{Quadcopters}
    \subparagraph{Propplers}
        \ProjectDeepHeading{2}{Pitch}
            the distance the proppler moves the air forward in one full rotation.
            The \ensuremath{\theta} is the rotation dependent on the distance to the center of the proppler
            \[
            \text{Pitch}(r) = 2\ensuremath{\pi} r \times \tan\big(\ensuremath{\theta}(r)\big)
            \]
    \subparagraph{Rotors}
        \ProjectSafeIncludePdf[pages=-]{../assets/rotor.pdf}
        \ProjectDeepHeading{1}{Brushless rotors}
            \ProjectDeepHeading{2}{ECG (Electronic Speed Controller)}
                provides the engine with patterns for it three phases as it has three magnetofields to connect voltage to each at once
            \ProjectDeepHeading{2}{PWM (Pulse with Modulation)}
                is to determine how longe the voltage is established in microseconds

\chapter{Math}

% Source: math/math.tex

\subsection{Math}

\subsubsection{Geometry}
    \paragraph{Euclidean}
        \paragraph{Coordinate systems}
        
    
    
    \subsubsection{Linear algebra}
\subsubsection{Vector spaces}
    \paragraph{Li algebra}
        \subparagraph{SO3}
        \subparagraph{SE3}

% Source: state_estimation.tex

\section{State estimation}

\subsection{Odometry}
    is to estimate robots position ONLY from \textbf{motion sensors}

% Source: controllers.tex

\section{Controller}
    \ProjectSafeIncludePdf[pages=-]{../assets/controllers.pdf}
    With regard to the distance the robot has with its mission, the controller generates commands for actuators to achieve that mission. Remember mission is a desired state
    
    
    \subsection{Model Predictive Control (MPC)}
    
    
    \subsection{Proportional–Integral–Derivative (PID)}

% Source: /home/donkarlo/Dropbox/repo/nd_language_ai_project/src/nd_language_ai/formal/computer/programming/tex/latex/code_clips/external_graphviz.tex

\newcommand{\includegraphvizrel}[2][0.9\textwidth]{%
    \IfEndWith{#2}{.dot}{%
        \StrGobbleRight{#2}{4}[\graphrelativebase]%
    }{%
        \edef\graphrelativebase{#2}%
    }%
    \directlua{
		local function get_current_directory()
			local pipe = io.popen("pwd")
			if pipe == nil then
				return nil
			end
			local directory = pipe:read("*l")
			pipe:close()
			return directory
		end

		local function get_parent_directory(path)
			return path:match("^(.*)/[^/]+$")
		end

		local function get_basename(path)
			return path:match("([^/]+)$")
		end

		local function find_src_directory(start_path)
			local current_path = start_path
			while current_path ~= nil do
				if get_basename(current_path) == "src" then
					return current_path
				end
				current_path = get_parent_directory(current_path)
			end
			return nil
		end

		local name = "\luaescapestring{\graphrelativebase}"
    local current_directory = get_current_directory()
    
    if current_directory == nil then
    tex.error("Could not determine current working directory")
    end
    
    local src_directory = find_src_directory(current_directory)
    if src_directory == nil then
    tex.error("Could not find a parent directory named src")
    end
    
    local project_root = get_parent_directory(src_directory)
    local output_directory = project_root .. "/out"
    local output_pdf_path = output_directory .. "/" .. name .. ".pdf"
    local output_pdf_directory = output_pdf_path:match("^(.*)/[^/]+$")
    
    if output_pdf_directory then
    os.execute('mkdir -p "' .. output_pdf_directory .. '"')
    end
    
    local build_command = 'dot -Tpdf "' .. name .. '.dot" -o "' .. output_pdf_path .. '"'
    local build_result = os.execute(build_command)
    
    if not (build_result == true or build_result == 0) then
    tex.error("Graphviz build failed", {"Command: " .. build_command})
    end
    
    tex.sprint("\\gdef\\graphoutputpath{" .. output_pdf_path .. "}")
}%
\begin{center}
    \ProjectMissingAsset{/graphoutputpath}
\end{center}
}

% Source: /home/donkarlo/Dropbox/repo/nd_language_ai_project/src/nd_language_ai/formal/computer/programming/tex/latex/code_clips/seehere.tex

\RequirePackage{xurl}

\makeatletter
\@ifpackageloaded{hyperref}{}{
    \RequirePackage{hyperref}
}
\makeatother

\AtBeginDocument{
    \hypersetup{
        colorlinks=true,
        linkcolor=blue,
        urlcolor=blue,
        citecolor=blue
    }
}

\providecommand{\seehere}[1]{\href{\#1}{see here}}
\providecommand{\clickurl}[1]{\href{\#1}{url}}
\providecommand{\seeherefooter}[1]{\href{\#1}{see here}\footnote{\tiny\url{\#1}}}
\providecommand{\clickurlfooter}[1]{\href{\#1}{here}\footnote{\tiny\url{\#1}}}

\section{Part}

\subsection{Mind}

\subsubsection{Cognition}

\paragraph{Part}

\subparagraph{Object level}

\ProjectDeepHeading{1}{Process}

\ProjectDeepHeading{2}{Parts}

\ProjectDeepHeading{3}{Thinking}

\ProjectDeepHeading{4}{Planning}

% Source: robot/part/mind/cognition/part/object_level/process/parts/thinking/planning/planning.tex

\todo{Planning is not decision making, planning uses decion making to choose the best action at each step}
    \todo{Biologically , it is called action planning}
    
    Planning is the process of thinking regarding the activities required to achieve a desired goal. Planning is based on foresight, the fundamental capacity for mental time travel. \seeherefooter{https://en.wikipedia.org/wiki/Planning}
    
    
    \ProjectDeepHeading{7}{From neorology perpective}
        Planning is one of the executive functions of the brain, encompassing the neurological processes involved in the formulation, evaluation and selection of a sequence of thoughts and actions to achieve a desired goal
        \seeherefooter{https://en.wikipedia.org/wiki/Planning\#Neurology}
    
    
    \ProjectDeepHeading{7}{Cognitive planning}
        
        Cognitive planning is one of the executive functions. It encompasses the neurological processes involved in the formulation, evaluation and selection of a sequence of thoughts and actions to achieve a desired goal.
        \seeherefooter{https://en.wikipedia.org/wiki/Planning_(cognitive)}
    
    
    \ProjectDeepHeading{7}{Automated planning and scheduling}
        is a branch of artificial intelligence that concerns the realization of strategies or action sequences, typically for execution by intelligent agents \seehere{https://en.wikipedia.org/wiki/Automated_planning_and_scheduling}
        
        Planning is the process of thinking regarding the activities required to achieve a desired goal. \seehere{https://en.wikipedia.org/wiki/Planning}
    
    
    \ProjectDeepHeading{7}{Motor planning}
        \clickurlfooter{https://en.wikipedia.org/wiki/Motor_planning} defines motor planning is a set of processes related to the preparation of a movement that occurs during the reaction time (the time between the presentation of a stimulus to a person and that person's initiation of a motor response).
    
    
    \ProjectDeepHeading{7}{Motion planning}
        \clickurlfooter{https://en.wikipedia.org/wiki/Motion_planning} defines motion planning, also path planning (also known as the navigation problem or the piano mover's problem) as a computational problem to find a sequence of valid configurations that moves the object from the source to destination.

\chapter{Ros}

% Source: ros/ros.tex

\subsection{ROS}
\subsubsection{Comparision with human body}
    \begin{table}[h!]
        \centering
        \resizebox{\textwidth}{!}{%
            \begin{tabular}{|>{\raggedright\arraybackslash}p{3cm}|>{\raggedright\arraybackslash}p{4cm}|>{\raggedright\arraybackslash}p{8.5cm}|}
                \hline
                \textbf{ROS Component}              & \textbf{Role}                          & \textbf{Human Analog}                                                                                               \\
                \hline
                \textbf{Nodes}                      & Independent modules                    & \textbf{Organs or neural nuclei} (e.g., basal ganglia, cerebellum, liver)                                           \\
                \hline
                \textbf{Topics (Pub/Sub)}           & Broadcast data streams                 & \textbf{Axons or nerve bundles}; broadcast of neural signals between neurons                                        \\
                \hline
                \textbf{Messages}                   & Data format / schema                   & \textbf{Patterns of action potentials / neurotransmitters} (information encoding)                                   \\
                \hline
                \textbf{Services (Req/Res)}         & Short request–response interaction     & \textbf{Spinal reflexes / point-to-point neural arcs}                                                               \\
                \hline
                \textbf{Actions}                    & Long-running task with feedback        & \textbf{Goal-directed behaviors} with continuous feedback (e.g., walking toward a goal; closed-loop motor behavior) \\
                \hline
                \textbf{Parameters}                 & Runtime configuration                  & \textbf{Hormones / homeostatic regulators} + “genes/learning” as pre-settings                                       \\
                \hline
                \textbf{tf / tf2}                   & Reference frames and transformations   & \textbf{Skeletal structure and proprioception (joint position sensing)}                                             \\
                \hline
                \textbf{Sensor drivers}             & Sensor input interfaces                & \textbf{Sensory organs:} eyes (camera), inner ear / vestibular system (IMU), skin (pressure/touch)                  \\
                \hline
                \textbf{Actuators / cmd\_vel}       & Output to actuators                    & \textbf{Muscles + motor neurons}                                                                                    \\
                \hline
                \textbf{ros\_control / controllers} & Stable feedback controllers            & \textbf{Cerebellum + spinal cord} (fine motor control, PID, balance)                                                \\
                \hline
                \textbf{Navigation / Nav2}          & Mapping and path planning              & \textbf{Hippocampus + parietal/occipital cortex} (cognitive map, place and direction cells)                         \\
                \hline
                \textbf{MoveIt}                     & Arm motion planning                    & \textbf{Motor/premotor cortex + cerebellum}                                                                         \\
                \hline
                \textbf{ROS Master (ROS1)}          & Node naming / discovery service        & \textbf{Thalamus} (central coordinator)                                                                             \\
                \hline
                \textbf{DDS (ROS2)}                 & Distributed communication layer        & \textbf{White matter / myelin} (communication network with QoS)                                                     \\
                \hline
                \textbf{QoS}                        & Delay and reliability guarantees       & \textbf{Axon diameter / degree of myelination} (speed and reliability of conduction)                                \\
                \hline
                \textbf{Clock}                      & Time reference                         & \textbf{Cardiac pacemaker (SA node)} + \textbf{neural rhythms}                                                      \\
                \hline
                \textbf{rosbag}                     & Data recording and playback            & \textbf{Hippocampal memory} (replay / recall during sleep)                                                          \\
                \hline
                \textbf{diagnostics}                & Health monitoring                      & \textbf{Pain / nociceptors + immune system} (fault and damage detection)                                            \\
                \hline
                \textbf{logging}                    & Event logging                          & \textbf{Cortical narrative traces / episodic memory records}                                                        \\
                \hline
                \textbf{launch}                     & Scenario or system startup             & \textbf{Daily routines / reflexive initiation behaviors}                                                            \\
                \hline
                \textbf{lifecycle manager}          & start / stop / configure state control & \textbf{Brainstem / autonomic nervous system (ANS)} – wakefulness, sleep, arousal                                   \\
                \hline
                \textbf{security}                   & Authentication and access restriction  & \textbf{Blood–brain barrier + immune system}                                                                        \\
                \hline
                \textbf{namespaces}                 & Separation of functional spaces        & \textbf{Functional brain regions / organ modularity}                                                                \\
                \hline
            \end{tabular}
        }
        \caption{Mapping between ROS components and their biological analogs in the human body.}
        \label{tab:ros_human_mapping}
    \end{table}
\subsubsection{Concepts}
\url{https://wiki.ros.org/ROS/Concepts}
    \paragraph{Filesystem level}
        \begin{itemize}
            \item packages
            \item meta packages
            \item package manifests
            \item repositories
            \item message/msg types
            \item service/srv types
        \end{itemize}
    \paragraph{Computation graph level concepts}
        \begin{itemize}
            \item Node \url{}
            \item Master \url{}
            \item Parameter server \url{}
            \item messages \url{}
            \item topics \url{}
            \item services \url{}
            \item bags \url{}
        \end{itemize}
    \paragraph{Names}
        bbc
        

\subsubsection{Setup}
    \paragraph{TF}
        \url{https://wiki.ros.org/navigation/Tutorials/RobotSetup/TF}
        \begin{figure}
            \centering
            \ProjectSafeIncludeGraphics{../assets/ros/tf.png}
            \caption{}
            \label{fig:tf}
        \end{figure}
\subsubsection{Names}
    see \url{https://wiki.ros.org/Names}
    \begin{itemize}
        \item global
        \item relative
        \item private
        \item base
    \end{itemize}

\subsubsection{TF}
    A package responsible for managing and publishing coordinate frames and transformation between them.
    
    TF keeps all the frames in a directed tree called ''Transform tree''
    
    
    map is coming from GPS which is from world frame. base\_link is robots body and lidar\_link as well
    map → odom → base\_link → lidar\_link
    
    \ProjectDeepHeading{1}{time\_stamped}
        Each transform is published with timestamp
        
\subsubsection{ROSPlan}
    The ROSPlan framework provides a collection of tools for AI Planning in a ROS system. ROSPlan has a variety of nodes which encapsulate plan, problem generation, and plan execution.
    \url{https://kcl-plan.github.io/ROSPlan/documentation/}
\subsubsection{Messages}
    \url{https://wiki.ros.org/msg}
    Message 
    \paragraph{geometry\_msgs}
    \paragraph{nav\_msgs}
        \ProjectDeepHeading{1}{Odometry}
            has Pose of 7 dimensions (3 for position and 4 quaternions) and 6 as twist or velocity (3 linear velocity and 3 angular)
    \paragraph{sensor\_msgs}
        \ProjectDeepHeading{1}{LaserScan}
    \paragraph{diagnostic\_msgs}
    \paragraph{actionlib\_msgs}
\subsubsection{Topics}
    Is a communication channel to which nodes can publish messages or subscribe for messages

    \ProjectDeepHeading{1}{Notes}
        They are ros runtime proccesses
        \begin{itemize}
            \item each topic can have as many publishers and subscribers that you need but the message types must be always the same
        \end{itemize}
\subsubsection{Node}
    Node a piece of runable code in ROS. It is either a c++/python code or a Gazebo plugin

    to list the running node:
    \begin{minted}{bash}
        rosnode list
    \end{minted}
\subsubsection{RQT Graph}
    Is a visual tool that shows nodes as boxes and topics as ovals and the direction of the edges show publisher and subscriber nodes
\subsubsection{RQT console}
    To visualize logs
\subsubsection{Bags}
\subsubsection{Frames}
     startes by world frame and then main robot's frame.
    
    Frames can have children. for example uav2/gps\_origin is the parent of uav2/fcu in which fcu is the the flight controller unit.
\subsubsection{Parameters}
    \url{https://wiki.ros.org/Parameter%20Server}
    \paragraph{Parameter server}
\subsubsection{Services}
    request-response communication form. Such as calling gazebo/get\_model\_state and getting models location
\subsubsection{Action}
    \url{https://wiki.ros.org/actionlib}
    \ProjectSafeIncludePdf[pages=-]{../assets/ros/action.pdf}
    \paragraph{Goal}

Is a combination of services and topics for a longterm mission with feedback and result and it could be terminated in the middle.
    \ProjectDeepHeading{1}{components}
        \ProjectDeepHeading{2}{Action client}
            Within application client
        \ProjectDeepHeading{2}{Action server}
            Within application server
    
    \ProjectDeepHeading{1}{Messages}
        \ProjectDeepHeading{2}{Goal}
        \ProjectDeepHeading{2}{Feedback}
        \ProjectDeepHeading{2}{Result}

\subsubsection{Navigation}
    \url{https://wiki.ros.org/navigation?distro=noetic}

\subsubsection{Planning}
\subsubsection{rviz}
    For visualizing data from sensors
\subsubsection{MAVROS}
\subsubsection{MAVLINK}
\subsubsection{QGroundControl}
\subsubsection{Flight controller}
    It is the hardware board and it usually contains
    \begin{itemize}
        \item IMU
        \item Barometer (for height)
        \item GNSS
        \item Microcontroller/cpu
        \item I/O ports
    \end{itemize}

    examples:
    \begin{itemize}
        \item Pixhawk
        \item Cubeorange
        \item Holybro
    \end{itemize}

\subsubsection{Autopilot}
    \paragraph{PX4}
        \paragraph{PXmini}
    \paragraph{Ardupilot}

\chapter{Control theory}

% Source: control_theory/control.tex

\subsubsection{Cognitive control}
        \cite{badre-2025-cognitive-control} argues that cognitive control is the system that enables flexible, goal-directed behavior by selecting actions according to context and by monitoring behavior to adjust it when needed. A central claim of the paper is that control depends strongly on how tasks are represented, especially whether they are encoded compositionally or as integrated conjunctive representations. Compositional representations support generalization and rapid learning of new tasks, but they are more vulnerable to interference and multitasking costs. Conjunctive representations reduce interference and support efficient action selection, but they are less suited for flexible transfer to new situations. The paper concludes that cognitive control emerges from the interaction between representational structure and feedback mechanisms such as monitoring errors, conflict, context, and outcomes

% Source: control_theory/control_input.tex

%
    %

% Source: control_theory/control_sgnal.tex

The numerical value given as a control input
    %
    %

% Source: control_theory/loop.tex

A control loop is the fundamental building block of control systems in general and industrial control systems in particular. It consists of the process sensor, the controller function, and the final control element (FCE) which controls the process necessary to automatically adjust the value of a measured process variable (PV) to equal the value of a desired set-point (SP).
    \url{https://en.wikipedia.org/wiki/Control_loop}
    % 
    %

% Source: control_theory/process_variable.tex

a process variable (PV; also process value or process parameter) is the current measured value of a particular part of a process which is being monitored or controlled. \seeherefooter{https://en.wikipedia.org/wiki/Process_variable}
    %
    %

% Source: control_theory/set_point.tex

a setpoint (SP;[1] also set point) is the desired or target value for an essential variable, or process value (PV) of a control system, \seeherefooter{https://en.wikipedia.org/wiki/Setpoint_(control_system)}
    %
    %

% Source: control_theory/single_input_and_single_output.tex

\url{https://en.wikipedia.org/wiki/MIMO}
    %
    %

% Source: control_theory/state_space_representation.tex

\clickurlfooter{https://en.wikipedia.org/wiki/State-space_representation} defines a state-space representation is a mathematical model of a physical system that uses state variables to track how inputs shape system behavior over time through first-order differential equations or difference equations. 
    %
    %

\section{Controller}

% Source: control_theory/controller/controller.tex

\[
        J = \sum_{k=0}^{N-1}
        \left( x_k - x_k^{\mathrm{ref}} \right)^T
        Q
        \left( x_k - x_k^{\mathrm{ref}} \right)
        +
        u_k^T R u_k
    \]
    
    %
    %

% Source: control_theory/controller/closed_loop.tex

A closed-loop controller or feedback controller is a control loop which incorporates feedback, in contrast to an open-loop controller or non-feedback controller.  \seeherefooter{https://en.wikipedia.org/wiki/Closed-loop_controller}
    %
    %

% Source: control_theory/controller/open_loop_controller.tex

In control theory, an open-loop controller, also called a non-feedback controller, is a control loop part of a control system in which the control action ("input" to the system[1]) is independent of the "process output", which is the process variable that is being controlled. \seeherefooter{https://en.wikipedia.org/wiki/Open-loop_controller}
    %
    %

\section{Feedback}

% Source: control_theory/feedback/feedback.tex

\clickurlfooter{https://en.wikipedia.org/wiki/Feedback} says feedback occurs when outputs of a system are routed back as inputs as part of a chain of cause and effect that forms a circuit or loop
    \begin{figure}[H]
        \centering
        \ProjectSafeIncludeGraphics[width=0.9\textwidth]{control_theory/feedback/feedback.png}
    \end{figure}
    %
    %

\subsection{Kinds}

% Source: control_theory/feedback/kinds/negative.tex

%
    %

% Source: control_theory/feedback/kinds/positive.tex

%
    %

\section{Kinds}

\subsection{Model predictive control}

% Source: control_theory/kinds/model_predictive_control/model_predictive_control.tex

MPC was first developed from early receding-horizon control ideas
    \citep{propoi-1963-use-of-linear-programming-methods-for-synthesizing-sampled-data-automatic-systems} and later became prominent in industrial process control through Model Predictive
    Heuristic Control and Dynamic Matrix Control
    \citep{richalet-1978-model-predictive-heuristic-control-applications-to-industrial-processes,
        cutler-1980-dynamic-matrix-control-a-computer-control-algorithm}.
    
    is an advanced method of process control that is used to control a process while satisfying a set of constraints. \seeherefooter{https://en.wikipedia.org/wiki/Model\_predictive\_control}
    
    
    \subparagraph{Cost function}
        
        \[
            J =
            \sum_{k=0}^{N-1}
            \left[
                \left( x_k - x_k^{\mathrm{ref}} \right)^{T}
                Q
                \left( x_k - x_k^{\mathrm{ref}} \right)
                +
                u_k^{T} R u_k
            \right]
        \]
        
        \begin{itemize}
            \item[$J$] is the total cost function to be minimized.
            \item[$k$] is the discrete time-step index.
            \item[$N$] is the prediction horizon length.
            \item[$x_k$] is the system state at time step $k$.
            \item[$x_k^{\mathrm{ref}}$] is the reference state at time step $k$.
            \item[$x_k - x_k^{\mathrm{ref}}$] is the state tracking error at time step $k$.
            \item[$(\cdot)^T$] denotes the transpose operator.
            \item[$Q$] is the weight of the state error. It determines how important the deviation of the system from the desired path or reference value is. For example, if the state includes position, velocity, and orientation, a larger weight can be assigned to position so that the controller tries harder to keep the position close to its reference value.
            \item[$u_k$] is the control input at time step $k$.
            \item[$R$] is the weight of the control input. It determines how costly it is to use the control command. If R is chosen to be large, the controller produces less aggressive control commands. As a result, the motion usually becomes smoother and calmer, but the trajectory-tracking error may become larger.
            \item[$u_k^T R u_k$] is the control-effort penalty.
        \end{itemize}
    %
    %

\subsection{Proportional integral derivative controller}

% Source: control_theory/kinds/proportional_integral_derivative_controller/proportional_integral_derivative_controller.tex

%
    %

\section{Process}

% Source: control_theory/process/process.tex

process control is a system used in modern manufacturing which uses the principles of control theory and physical industrial control systems to monitor, control and optimize continuous industrial production processes using control algorithms. \seeherefooter{https://en.wikipedia.org/wiki/Industrial_process_control}
    %
    %

\section{System}

% Source: control_theory/system/system.tex

%
    %

\chapter{Frameworks}

\section{Gazebo}

\subsection{Sdf}

% Source: frameworks/gazebo/sdf/sdf.tex

XML format
    
    Can read URDF and convert it to SDF

\subsubsection{Example}

% Source: frameworks/gazebo/sdf/example/quadcopter_and_tree_in_day_light.tex

\begin{minted}{xml}
        <?xml version="1.0" ?>
        <sdf version="1.6">
          <world name="day_world">
            <gravity>0 0 -9.81</gravity>
        
            <!-- Day light -->
            <scene>
              <ambient>0.6 0.6 0.6 1</ambient>
              <background>0.7 0.9 1.0 1</background>
              <shadows>true</shadows>
            </scene>
        
            <light name="sun" type="directional">
              <cast_shadows>true</cast_shadows>
              <pose>0 0 10 0 0 0</pose>
              <diffuse>1 1 1 1</diffuse>
              <specular>0.2 0.2 0.2 1</specular>
              <direction>-0.3 0.2 -1</direction>
            </light>
        
            <!-- Ground -->
            <model name="ground_plane">
              <static>true</static>
              <link name="link">
                <collision name="collision">
                  <geometry>
                    <plane>
                      <normal>0 0 1</normal>
                      <size>200 200</size>
                    </plane>
                  </geometry>
                </collision>
                <visual name="visual">
                  <geometry>
                    <plane>
                      <normal>0 0 1</normal>
                      <size>200 200</size>
                    </plane>
                  </geometry>
                  <material>
                    <ambient>0.2 0.6 0.2 1</ambient>
                    <diffuse>0.2 0.6 0.2 1</diffuse>
                  </material>
                </visual>
              </link>
            </model>
        
            <!-- A simple tree -->
            <model name="tree_1">
              <static>true</static>
              <pose>8 3 0 0 0 0</pose>
        
              <link name="trunk">
                <collision name="trunk_collision">
                  <pose>0 0 1 0 0 0</pose>
                  <geometry>
                    <cylinder>
                      <radius>0.25</radius>
                      <length>2.0</length>
                    </cylinder>
                  </geometry>
                </collision>
                <visual name="trunk_visual">
                  <pose>0 0 1 0 0 0</pose>
                  <geometry>
                    <cylinder>
                      <radius>0.25</radius>
                      <length>2.0</length>
                    </cylinder>
                  </geometry>
                  <material>
                    <ambient>0.35 0.2 0.1 1</ambient>
                    <diffuse>0.35 0.2 0.1 1</diffuse>
                  </material>
                </visual>
              </link>
        
              <link name="crown">
                <pose>0 0 2.4 0 0 0</pose>
                <collision name="crown_collision">
                  <geometry>
                    <sphere>
                      <radius>1.0</radius>
                    </sphere>
                  </geometry>
                </collision>
                <visual name="crown_visual">
                  <geometry>
                    <sphere>
                      <radius>1.0</radius>
                    </sphere>
                  </geometry>
                  <material>
                    <ambient>0.1 0.5 0.1 1</ambient>
                    <diffuse>0.1 0.5 0.1 1</diffuse>
                  </material>
                </visual>
              </link>
            </model>
        
            <!-- Quadcopter (simple rigid body + sensors) -->
            <model name="quadcopter">
              <static>false</static>
              <pose>0 0 1.0 0 0 0</pose>
        
              <link name="base_link">
                <inertial>
                  <mass>1.5</mass>
                  <inertia>
                    <ixx>0.03</ixx>
                    <iyy>0.03</iyy>
                    <izz>0.05</izz>
                    <ixy>0</ixy>
                    <ixz>0</ixz>
                    <iyz>0</iyz>
                  </inertia>
                </inertial>
        
                <collision name="body_collision">
                  <geometry>
                    <box>
                      <size>0.35 0.35 0.12</size>
                    </box>
                  </geometry>
                </collision>
        
                <visual name="body_visual">
                  <geometry>
                    <box>
                      <size>0.35 0.35 0.12</size>
                    </box>
                  </geometry>
                  <material>
                    <ambient>0.2 0.2 0.25 1</ambient>
                    <diffuse>0.2 0.2 0.25 1</diffuse>
                  </material>
                </visual>
        
                <!-- IMU sensor -->
                <sensor name="imu_sensor" type="imu">
                  <always_on>true</always_on>
                  <update_rate>200</update_rate>
                  <imu>
                    <angular_velocity>
                      <x><noise type="gaussian"><mean>0</mean><stddev>0.0002</stddev></noise></x>
                      <y><noise type="gaussian"><mean>0</mean><stddev>0.0002</stddev></noise></y>
                      <z><noise type="gaussian"><mean>0</mean><stddev>0.0002</stddev></noise></z>
                    </angular_velocity>
                    <linear_acceleration>
                      <x><noise type="gaussian"><mean>0</mean><stddev>0.01</stddev></noise></x>
                      <y><noise type="gaussian"><mean>0</mean><stddev>0.01</stddev></noise></y>
                      <z><noise type="gaussian"><mean>0</mean><stddev>0.01</stddev></noise></z>
                    </linear_acceleration>
                  </imu>
        
                  <!-- If you use gazebo_ros_pkgs, this publishes a ROS topic -->
                  <plugin name="imu_plugin" filename="libgazebo_ros_imu_sensor.so">
                    <robotNamespace>/quadcopter</robotNamespace>
                    <bodyName>base_link</bodyName>
                    <topicName>imu</topicName>
                    <frameName>base_link</frameName>
                    <updateRate>200.0</updateRate>
                  </plugin>
                </sensor>
        
                <!-- GPS sensor (Gazebo "navsat") -->
                <sensor name="gps_sensor" type="navsat">
                  <always_on>true</always_on>
                  <update_rate>10</update_rate>
                  <navsat>
                    <position_sensing>
                      <horizontal>
                        <noise type="gaussian"><mean>0</mean><stddev>0.5</stddev></noise>
                      </horizontal>
                      <vertical>
                        <noise type="gaussian"><mean>0</mean><stddev>1.0</stddev></noise>
                      </vertical>
                    </position_sensing>
                    <velocity_sensing>
                      <horizontal>
                        <noise type="gaussian"><mean>0</mean><stddev>0.1</stddev></noise>
                      </horizontal>
                      <vertical>
                        <noise type="gaussian"><mean>0</mean><stddev>0.2</stddev></noise>
                      </vertical>
                    </velocity_sensing>
                  </navsat>
        
                  <!-- If you use gazebo_ros_pkgs, this publishes a ROS topic -->
                  <plugin name="gps_plugin" filename="libgazebo_ros_gps.so">
                    <robotNamespace>/quadcopter</robotNamespace>
                    <bodyName>base_link</bodyName>
                    <topicName>gps</topicName>
                    <frameName>base_link</frameName>
                    <updateRate>10.0</updateRate>
                  </plugin>
                </sensor>
        
                <!-- 2D LiDAR (ray) -->
                <sensor name="lidar_2d" type="ray">
                  <pose>0.15 0 0.05 0 0 0</pose>
                  <always_on>true</always_on>
                  <update_rate>10</update_rate>
                  <ray>
                    <scan>
                      <horizontal>
                        <samples>720</samples>
                        <resolution>1</resolution>
                        <min_angle>-3.141592653589793</min_angle>
                        <max_angle>3.141592653589793</max_angle>
                      </horizontal>
                    </scan>
                    <range>
                      <min>0.15</min>
                      <max>12.0</max>
                      <resolution>0.01</resolution>
                    </range>
                    <noise type="gaussian">
                      <mean>0</mean>
                      <stddev>0.01</stddev>
                    </noise>
                  </ray>
        
                  <!-- If you use gazebo_ros_pkgs, this publishes sensor_msgs/LaserScan -->
                  <plugin name="lidar_plugin" filename="libgazebo_ros_laser.so">
                    <robotNamespace>/quadcopter</robotNamespace>
                    <topicName>scan</topicName>
                    <frameName>lidar_link</frameName>
                    <updateRate>10.0</updateRate>
                  </plugin>
                </sensor>
              </link>
        
              <!-- An explicit frame for the lidar (optional but helpful) -->
              <link name="lidar_link">
                <pose>0.15 0 1.05 0 0 0</pose>
                <inertial>
                  <mass>0.01</mass>
                  <inertia>
                    <ixx>1e-6</ixx><iyy>1e-6</iyy><izz>1e-6</izz>
                    <ixy>0</ixy><ixz>0</ixz><iyz>0</iyz>
                  </inertia>
                </inertial>
              </link>
        
              <joint name="lidar_joint" type="fixed">
                <parent>base_link</parent>
                <child>lidar_link</child>
              </joint>
        
              <!-- Aerodynamics / flight control is intentionally omitted (this is a sensor-ready SDF example). -->
            </model>
          </world>
        </sdf>
    \end{minted}

\subsubsection{Ros plugin}

% Source: frameworks/gazebo/sdf/ros_plugin/ros_plugin.tex

With codes such as this you can send sensor data via a topic to ROS
    \begin{minted}{xml}
        <sensor name="imu" type="imu">
          <always_on>true</always_on>
          <update_rate>200</update_rate>
        
          <plugin name="imu_plugin" filename="libgazebo_ros_imu_sensor.so">
            <robotNamespace>/simple_robot</robotNamespace>
            <topicName>imu</topicName>
            <frameName>base_link</frameName>
          </plugin>
        </sensor>
    \end{minted}

\section{Ros}

\subsection{Xarco}

% Source: frameworks/ros/xarco/xarco.tex

\subparagraph{Example}
        \begin{minted}{xml}
            <?xml version="1.0"?>
                <robot xmlns:xacro="\url{http://www.ros.org/wiki/xacro"} name="simple_quadcopter">
                
                  <!-- --------------------
                       Parameters
                  --------------------- -->
                  <xacro:property name="base_mass" value="1.5"/>
                  <xacro:property name="base_size_x" value="0.35"/>
                  <xacro:property name="base_size_y" value="0.35"/>
                  <xacro:property name="base_size_z" value="0.12"/>
                
                  <!-- --------------------
                       Helper macros
                  --------------------- -->
                  <xacro:macro name="box_inertial" params="mass x y z">
                    <inertial>
                      <mass value="${mass}"/>
                      <!-- Simple box inertia (approx) -->
                      <inertia
                        ixx="${(1.0/12.0) * mass * (y*y + z*z)}"
                        iyy="${(1.0/12.0) * mass * (x*x + z*z)}"
                        izz="${(1.0/12.0) * mass * (x*x + y*y)}"
                        ixy="0" ixz="0" iyz="0"/>
                    </inertial>
                  </xacro:macro>
                
                  <!-- --------------------
                       Base link (quadcopter body)
                  --------------------- -->
                  <link name="base_link">
                    <xacro:box_inertial mass="${base_mass}" x="${base_size_x}" y="${base_size_y}" z="${base_size_z}"/>
                
                    <collision name="base_collision">
                      <geometry>
                        <box size="${base_size_x} ${base_size_y} ${base_size_z}"/>
                      </geometry>
                    </collision>
                
                    <visual name="base_visual">
                      <geometry>
                        <box size="${base_size_x} ${base_size_y} ${base_size_z}"/>
                      </geometry>
                    </visual>
                  </link>
                
                  <!-- --------------------
                       GPS link + joint
                  --------------------- -->
                  <link name="gps_link">
                    <inertial>
                      <mass value="0.05"/>
                      <inertia ixx="1e-5" iyy="1e-5" izz="1e-5" ixy="0" ixz="0" iyz="0"/>
                    </inertial>
                
                    <visual name="gps_visual">
                      <geometry>
                        <box size="0.04 0.04 0.02"/>
                      </geometry>
                    </visual>
                  </link>
                
                  <joint name="gps_joint" type="fixed">
                    <parent link="base_link"/>
                    <child link="gps_link"/>
                    <!-- mounted on top -->
                    <origin xyz="0 0 ${base_size_z/2.0 + 0.02}" rpy="0 0 0"/>
                  </joint>
                
                  <!-- --------------------
                       LiDAR link + joint
                  --------------------- -->
                  <link name="lidar_link">
                    <inertial>
                      <mass value="0.12"/>
                      <inertia ixx="3e-5" iyy="3e-5" izz="3e-5" ixy="0" ixz="0" iyz="0"/>
                    </inertial>
                
                    <visual name="lidar_visual">
                      <geometry>
                        <cylinder radius="0.03" length="0.05"/>
                      </geometry>
                    </visual>
                  </link>
                
                  <joint name="lidar_joint" type="fixed">
                    <parent link="base_link"/>
                    <child link="lidar_link"/>
                    <!-- mounted in front-top -->
                    <origin xyz="0.12 0 ${base_size_z/2.0 + 0.03}" rpy="0 0 0"/>
                  </joint>
                
                  <!-- --------------------
                       (Optional) Gazebo tags: materials + sensor placeholders
                       This is only a placeholder; actual Gazebo sensor plugins depend on your setup.
                  --------------------- -->
                  <gazebo reference="base_link">
                    <material>Gazebo/Grey</material>
                  </gazebo>
                
                  <gazebo reference="gps_link">
                    <material>Gazebo/Black</material>
                    <!-- Placeholder: real GPS in Gazebo usually uses a plugin -->
                  </gazebo>
                
                  <gazebo reference="lidar_link">
                    <material>Gazebo/Black</material>
                    <!-- Placeholder: real LiDAR uses a ray sensor + plugin -->
                  </gazebo>
                
                </robot>
        \end{minted}

\subsubsection{Urdf}

% Source: frameworks/ros/xarco/urdf/urdf.tex

Structure
    base\_link
         ├── joint1
         │    └── link1
         │         └── joint2
         │              └── link2
    
    Each sub tree is a limb




    \begin{minted}{xml}
        <sdf version="1.6">
          <model name="robot">
            <link name="base_link">
              ...
            </link>
          </model>
        </sdf>
        
    \end{minted}

\chapter{From phd review}

% Source: from_phd_review/phd_thesis_review.tex

%table of contents




The review part is to gather all existing knowledge

















% List todos

% Bibliography

 % Reference the BibTeX file here

\subsubsection{Plan}

% Source: robot/part/mind/plan/plan.tex

A plan is typically any list of steps, with details of timing and resources, used to achieve an objective. \seehere{https://en.wikipedia.org/wiki/Plan}
    %
    %

\section{Robotics}

\subsection{Robot}

\subsubsection{Parts}

\paragraph{Mind}

\subparagraph{From data to wisdom}

% Source: from_phd_review/robotics/robot/parts/mind/from_data_to_wisdom/from-data-to-information-theory.tex

In this chapter we draw the boundries between data , sensory data, sensor and data fusion, information , etc
        
        
        \ProjectDeepHeading{3}{Data}
        
        
        \ProjectDeepHeading{3}{Information} Information is processed, structured, or organized data that has been given meaning and context, making it understandable and useful. In the following table you can find the difference between information and perception
            \begin{tabular}{|l|l|l|}
                \hline
                \textbf{Aspect}             & \textbf{Perception}                                       & \textbf{Information}                                              \\ \hline
                \textbf{Definition}         & The process of interpreting sensory input.                & Processed data that provides meaning and context.                 \\ \hline
                \textbf{Process vs. Output} & A dynamic process involving sensors and interpretation. & A static or dynamic product resulting from processing. \\ \hline
                \textbf{Role}               & A mechanism to generate understanding of the environment. & Provides insights or knowledge for decision-making.               \\ \hline
                \textbf{Dependency}         & Requires sensory input or stimuli.                        & Can exist independently (e.g., in a database or memory).          \\ \hline
                \textbf{Examples in Humans} & Seeing a tree, hearing a sound.                           & Recognizing the tree as an oak or identifying the sound as music. \\ \hline
                \textbf{Examples in Robots} & Using cameras and algorithms to recognize objects.        & Storing recognized objects as labeled data.                       \\ \hline
            \end{tabular}
        
        
        \ProjectDeepHeading{3}{Knowledge}
            Refers to the collection of information, facts, concepts, or skills that an individual has acquired through learning, experience, or education. An example of knowledge in robotics or any other intelligent agent, a database of facts or pre-programmed rules (e.g., rules of chess) can be considered as knowledge. In humans, knowing that water boils at 100 degree is a knowledge. Knowledge is the static outcome of learning and experience; it is what we know.
            \begin{itemize}
                \item Is on aspect of knowledge predicting future states from current and past states?
                \item What is the difference between knowledge and information?
            \end{itemize}
        
        
        \ProjectDeepHeading{3}{Information theory}
        
        
        \ProjectDeepHeading{3}{Signal processing}
        
        
        \ProjectDeepHeading{3}{Questions}

% Source: from_phd_review/robotics/robot/parts/mind/learning_in_robots.tex

\ProjectDeepHeading{1}{Learning in Robots}
Learning refers to the process of acquiring or updating knowledge, skills, or behavior based on data, experience, or instruction. It involves adapting internal models to better represent the external world or to improve mission performance. Robots use learning algorithms like supervised, unsupervised, or reinforcement learning to build capabilities such as object recognition or path plan. For example, a robot might learn to classify objects in its environment by training on labeled images. Learning provides the foundational knowledge needed for reasoning and decision-making.

\ProjectDeepHeading{2}{Learning Types in Robotics}

\ProjectDeepHeading{3}{Supervised Learning}
Supervised learning is a method where a model learns to map inputs to outputs using labeled datasets. It is extensively used for missions requiring high accuracy, such as object detection and signal classification. In drones, supervised learning can be employed to classify objects in the environment using camera data. For instance, a drone can be trained to recognize trees, vehicles, and humans by using a dataset of labeled aerial images. This capability is crucial for autonomous navigation and obstacle avoidance in complex terrains.

\ProjectDeepHeading{3}{Unsupervised Learning}
Unsupervised learning discovers hidden patterns or intrinsic structures in unlabeled data, making it useful for exploring unknown environments. In drones, unsupervised learning can process data from LiDAR and IMU sensors to cluster 3D point clouds into meaningful groups, such as open paths, vegetation, or obstacles. This allows the drone to map and understand uncharted areas without explicit labeling, enabling autonomous navigation in new environments. 

\ProjectDeepHeading{3}{Semi-Supervised Learning}
Semi-supervised learning combines labeled and unlabeled data to improve model efficiency in data-scarce scenarios. For a drone navigating a forest, a small amount of labeled LiDAR data can help infer the characteristics of a larger, unlabeled dataset, enabling it to identify safe paths or landing zones. This approach bridges the gap between fully supervised and unsupervised learning in dynamic environments. 

\ProjectDeepHeading{3}{Multi Modal learning}


\ProjectDeepHeading{3}{Cross modal learning}

\ProjectDeepHeading{3}{Transfer Learning}
Transfer learning reuses knowledge from a pre-trained model to improve performance on a related mission, reducing the need for extensive training. In drones, a model pre-trained to detect vehicles in urban environments using camera data can be fine-tuned to recognize animals in rural areas. This adaptability allows drones to operate effectively across diverse missions with minimal retraining.

\ProjectDeepHeading{3}{Self-Supervised Learning}
Self-supervised learning generates labels from data itself, enabling the model to learn meaningful representations without manual annotations. A drone can use self-supervised learning to predict future LiDAR readings based on its current IMU and GPS data, allowing it to understand motion dynamics and adapt to changing environmental conditions autonomously. 

\ProjectDeepHeading{3}{Meta-Learning (Learning to Learn)}
Meta-learning, or "learning to learn," equips models with the ability to quickly adapt to new missions using minimal data. For example, a drone trained to navigate deserts can adapt to forested environments with only a few examples by leveraging its prior experience with obstacle avoidance and pathfinding. This rapid adaptability is essential for multi-environment missions.

\ProjectDeepHeading{3}{Continual Learning}
Continual learning allows drones to acquire new skills incrementally while retaining previously learned knowledge, addressing the problem of catastrophic forgetting. For instance, a drone initially trained to navigate open fields can incrementally learn to operate in dense urban areas without forgetting its prior navigation skills. This capability is vital for long-term autonomous operation in dynamic settings. 

\ProjectDeepHeading{3}{Reinforcement Learning}
Reinforcement learning trains agents to make sequential decisions by maximizing cumulative rewards. In drones, reinforcement learning can optimize flight paths by balancing energy consumption and coverage. For instance, a drone might learn to map a large area efficiently by receiving rewards for increased coverage and penalties for excessive battery usage, refining its navigation strategy through trial and error. 

\textbf{Questions}
\\
Are rewards also learnable? For example, if I have a drone that sends random control inputs to its rotators in a jungle, can it understand control inputs are good that help it increase it power source or balance?

\ProjectDeepHeading{4}{Q-learning}
Q-learning is a reinforcement learning algorithm that enables an agent to find optimal actions in various states without needing a model of the environment. A drone navigating a maze of obstacles can use Q-learning to identify the best path to its destination by updating a Q-table based on LiDAR-based proximity sensors and rewards for collision-free navigation.

\ProjectDeepHeading{3}{Active Learning}
Active learning minimizes the effort required for labeling data by focusing on the most informative samples. For example, a drone might identify ambiguous objects in its camera feed, such as unusual terrain features, and request human input for labeling. This targeted approach improves model accuracy while reducing annotation costs, making it ideal for exploration missions.

\ProjectDeepHeading{3}{Imitation Learning}
Imitation learning enables robots to mimic expert behaviors by observing demonstrations. For instance, a drone can learn to navigate through narrow corridors by analyzing the control inputs and camera footage of an experienced pilot. This approach reduces the need for explicit programming in scenarios requiring complex decision-making. 

\ProjectDeepHeading{3}{Online Learning}
Online learning allows models to update incrementally as new data becomes available, ensuring adaptability to dynamic environments. A delivery drone, for example, might continuously refine its flight model based on real-time wind data from its IMU and GPS sensors, maintaining stability and improving efficiency in varying conditions. 

\ProjectDeepHeading{3}{Evolutionary Learning}
Evolutionary learning optimizes models through algorithms inspired by natural selection, such as mutation and crossover. A swarm of drones might use evolutionary learning to evolve collaborative strategies for mapping large areas, where each generation improves performance by combining the strengths of previous attempts. 

\ProjectDeepHeading{3}{Adversarial Learning}
Adversarial learning involves training models in adversarial setups to improve robustness and generalization. For example, a drone can learn to detect and reject spoofed GPS signals by training with adversarially generated GPS noise, ensuring reliable operation even in challenging environments. 

\ProjectDeepHeading{3}{Federated Learning}
Federated learning enables collaborative model training across multiple drones without sharing raw data, preserving privacy and security. Each drone trains a local model using its camera and LiDAR data, and the aggregated updates improve a global model shared across the fleet. This approach is particularly useful in sensitive applications such as surveillance or disaster response. 

\ProjectDeepHeading{3}{Curriculum Learning}
Curriculum learning organizes training by progressing from simpler to more complex missions, mimicking human learning. For example, a drone may first learn to hover, then navigate open spaces, and finally handle complex, obstacle-rich environments like dense forests. This structured approach improves learning efficiency and robustness.

\ProjectDeepHeading{3}{Multi-Mission Learning}
Multi-mission learning enables models to perform several missions simultaneously by leveraging shared representations. For instance, a drone might simultaneously detect obstacles using LiDAR and estimate wind conditions with IMU data, optimizing its operations for both safety and efficiency.

\ProjectDeepHeading{3}{Intrinsic Motivation Learning}
Intrinsic motivation learning encourages agents to explore autonomously by leveraging curiosity or novelty. A drone might prioritize exploring areas with high variability in terrain, using LiDAR and GPS data to discover new landmarks, driven by its intrinsic desire to gather information about its environment. 

\ProjectDeepHeading{3}{Zero-Shot Learning}
Zero-shot learning enables models to generalize to unseen classes or missions using semantic relationships or prior knowledge. For example, a drone trained to identify common tree species can recognize a new species based on its similarity to known categories, using visual and structural features from its camera and LiDAR data.

\ProjectDeepHeading{3}{Few-Shot Learning}
Few-shot learning allows models to learn missions with minimal labeled examples, leveraging prior knowledge and generalization. A search-and-rescue drone could quickly identify specific types of debris or hazards in disaster areas using only a few labeled images, enhancing its response capabilities in emergencies.

\ProjectDeepHeading{3}{sensory motor learning}

\ProjectDeepHeading{2}{Reasoning}
Reasoning is the cognitive process of deriving conclusions, solving problems, or making decisions based on existing knowledge or logic. Unlike learning, reasoning often operates on pre-existing knowledge or models to produce actionable insights. It encompasses techniques like deductive reasoning (general to specific), inductive reasoning (specific to general), and abductive reasoning (best plausible explanation). In robotics, reasoning allows systems to decide the best course of action, such as selecting a path to avoid an obstacle. Reasoning leverages learned information to achieve missions efficiently.

\ProjectDeepHeading{2}{Inference}
Inference is the act of drawing specific conclusions or making deductions based on evidence or observed data. It can involve probabilistic reasoning, such as Bayesian inference, where a robot estimates the likelihood of events based on prior knowledge and current observations. For example, a drone might infer that a detected object is a tree based on sensory data from its camera and LiDAR. Inference often operates as a step within broader reasoning processes, making it a more targeted, evidence-driven activity. 

Inference techniques in robotics enable robots to derive conclusions, make decisions, or interpret data based on sensory inputs, learned models, and pre-existing knowledge. These techniques handle uncertainties, optimize actions, and allow robots to adapt to dynamic environments. Below, we discuss key inference methods used in robotics.



\ProjectDeepHeading{3}{Comparison Table}
\begin{table}[h!]
    \centering
    \begin{tabular}{|l|l|l|l|}
        \hline
        \textbf{Aspect} & \textbf{Learning} & \textbf{Reasoning} & \textbf{Inference} \\ \hline
        \textbf{Definition} & Acquiring knowledge or skills through experience or data. & Deriving conclusions or solving problems based on knowledge or logic. & Drawing specific conclusions from data or observations. \\ \hline
        \textbf{Input} & External data or experiences. & Pre-existing knowledge or models. & Evidence or current observations. \\ \hline
        \textbf{Output} & Updated knowledge or internal model. & Decisions, plans, or broader insights. & Specific, evidence-based conclusions. \\ \hline
        \textbf{Scope} & Broad, involves building general understanding. & Broader problem-solving or decision-making process. & Narrow, focused on individual conclusions. \\ \hline
        \textbf{Examples in Robots} & Training a model to recognize objects in the environment. & Deciding the best route to avoid obstacles based on a map. & Inferring that a detected object is a tree using sensory data. \\ \hline
        \textbf{Relation} & Provides the foundation for reasoning. & Uses knowledge from learning to guide inference and decision-making. & Supports reasoning by providing specific conclusions. \\ \hline
    \end{tabular}
    \caption{Comparison of Learning, Reasoning, and Inference}
\end{table}
\ProjectDeepHeading{3}{Inference techniques in robotics}

    \ProjectDeepHeading{4}{Bayesian Inference}
    Bayesian inference is a probabilistic technique that updates the likelihood of a hypothesis based on prior knowledge and observed data. In robotics, it is widely used for sensor fusion, where data from multiple sensors like LiDAR, cameras, and IMUs are combined to estimate the robot's position or environment. This method handles uncertainties effectively, making it ideal for localization missions such as Monte Carlo Localization.

    \ProjectDeepHeading{4}{Maximum Likelihood Estimation (MLE)}
    Maximum Likelihood Estimation determines the parameters of a model that maximize the likelihood of observed data. In robotics, MLE is often used for model fitting, such as estimating the motion dynamics of a robot or calibrating sensor parameters. For example, a robot can use MLE to optimize its camera parameters to improve visual recognition accuracy. 

    \ProjectDeepHeading{4}{Monte Carlo Methods}
    Monte Carlo methods are stochastic approaches that use random sampling to estimate outcomes and probabilities. In robotics, Monte Carlo Localization is a prominent application, where particle filters estimate a robot’s pose using sensory data and motion models. These methods are particularly effective in high-dimensional and nonlinear systems, such as plan collision-free paths for drones.

    \ProjectDeepHeading{4}{Deductive Inference}
    Deductive inference involves drawing specific conclusions from general rules or knowledge. Robots employ this approach for rule-based decision-making, where predefined if-then rules dictate actions based on observed states. For example, a robot might deduce that an obstacle detected ahead requires it to stop and plan an alternate path. This method is intuitive and interpretable. 

    \ProjectDeepHeading{4}{Neural Network-Based Inference}
    Neural network-based inference uses artificial neural networks to predict outcomes based on learned patterns in data. In robotics, this approach is widely applied in missions like object detection using convolutional neural networks (CNNs) or motion prediction with recurrent neural networks (RNNs). For instance, a drone might infer the trajectory of a moving object to avoid collisions. Neural inference excels in handling complex, high-dimensional data.

    \ProjectDeepHeading{4}{Fuzzy Inference}
    Fuzzy inference applies fuzzy logic to handle uncertainties and vagueness in data. Robots use this method for decision-making in imprecise or noisy environments. For example, a robotic vacuum might decide how much to adjust its path based on partially observable obstacles. Fuzzy inference provides intuitive and human-like reasoning capabilities, making it suitable for navigation and human-robot interaction.

% Source: from_phd_review/robotics/robot/parts/mind/machine-learning-models.tex

\ProjectDeepHeading{1}{Models}
\ProjectDeepHeading{2}{Different models}
\ProjectDeepHeading{3}{Classification Models}
Classification models are designed to categorize input data into predefined classes by learning decision boundaries in the feature space. These models work by analyzing patterns in the training data and assigning new inputs to one of the learned categories. Common algorithms include decision trees, support vector machines, and neural networks. They are widely used in missions like object recognition in computer vision, where an image is classified as "cat" or "dog," and anomaly detection, where unusual patterns in data are flagged. Their performance heavily depends on the quality and diversity of the training data.
    \ProjectDeepHeading{4}{Discriminative Models}
Discriminative models focus on learning the decision boundary between classes by modeling the conditional probability of the output given the input. Unlike generative models, they do not attempt to model the entire data distribution, making them efficient for classification missions. Examples include logistic regression, random forests, and discriminative neural networks. These models are used in applications like spam email detection, where the mission is to classify emails as "spam" or "not spam," and medical diagnosis, where symptoms are mapped to diseases. They are favored for their accuracy and efficiency in prediction missions.


\ProjectDeepHeading{3}{Regression Models}
Regression models predict continuous numerical values by capturing the relationships between input features and outputs. These models fit a function to the data points, minimizing the error between predicted and actual values. For instance, linear regression models establish linear relationships, while more complex techniques like polynomial regression or neural networks capture non-linear relationships. Applications include predicting house prices based on features like size and location, and estimating battery life in drones using flight parameters. Regression models are foundational in missions requiring precise numerical estimations.

\ProjectDeepHeading{3}{Predictive Models}
Predictive models forecast future outcomes by analyzing patterns in historical and present data. They use statistical methods or machine learning algorithms to make informed predictions. For example, a time series model might predict future stock prices by examining past trends, or trajectory prediction models estimate the future path of a drone based on current velocity and environmental data. Predictive models are critical in applications like weather forecasting, financial analysis, and autonomous navigation, where anticipating future states is essential for decision-making. 

\ProjectDeepHeading{3}{Generative Models}
Generative models learn the underlying distribution of training data to generate new, realistic data points. These models aim to capture the essence of the data, enabling missions like image synthesis, text generation, and data augmentation. Popular examples include Generative Adversarial Networks (GANs), Variational Autoencoders (VAEs), and diffusion models. For instance, GANs can generate high-resolution images of artificial objects, while VAEs can create realistic reconstructions of training data. Generative models are vital in applications like creating synthetic datasets for training or enhancing artistic creativity.



\ProjectDeepHeading{3}{Hybrid Models}
Hybrid models combine the strengths of generative and discriminative approaches to achieve improved accuracy and robustness. These models leverage generative methods to learn the underlying data distribution and discriminative techniques to optimize decision boundaries. For instance, a hybrid model might use a generative network to preprocess data and a discriminative network for classification. Applications include semi-supervised learning, where labeled and unlabeled data are combined, and robotics, where hybrid models can improve both perception and action plan. Hybrid models are particularly effective in missions with limited labeled data.

\ProjectDeepHeading{3}{Probabilistic Models}
Probabilistic models use principles of probability theory to represent uncertainty and make predictions. These models calculate the likelihood of outcomes given the input data, enabling robust decision-making in uncertain environments. Examples include Bayesian networks, which model probabilistic relationships among variables, and Hidden Markov Models (HMMs), used for sequential data analysis. In robotics, probabilistic models are essential for missions like localization, where the robot estimates its position in an environment with uncertain sensor readings. These models excel in scenarios requiring explicit uncertainty quantification.

\ProjectDeepHeading{3}{Neural Models}
Neural models are advanced machine learning architectures inspired by the human brain, designed for specific missions through layers of artificial neurons. Convolutional Neural Networks (CNNs) excel at processing visual data, such as recognizing objects in drone imagery. Transformers, on the other hand, are specialized for sequence data and dominate missions like natural language understanding and time-series analysis. These models are highly flexible, allowing applications in diverse fields, including robotics, where drones use CNNs for obstacle detection and transformers for mission plan based on sequential data. Neural models continue to evolve, shaping modern AI advancements.






\ProjectDeepHeading{2}{Pattern, pattern recognition models and anomalies}
    \ProjectDeepHeading{3}{Pattern} A pattern in pattern recognition refers to a set of features or data points that represent some meaningful structure or regularity. It can be thought of as a recognizable arrangement or characteristic that distinguishes one set of data from another. Patterns are typically identified within raw data through various algorithms and methods. Two ways of representing patterns are:
    \begin{itemize}
        \item Feature vectors: Numerical representations of attributes (e.g., size, color, or frequency).
        \item Graphs or sets: Representing relationships or groupings (e.g., network graphs).
    \end{itemize}
    \ProjectDeepHeading{3}{Model}
    A model in pattern recognition is a mathematical framework that learns to identify or classify patterns from data. The model is built using algorithms trained on sample data from patterns and is then used to recognize patterns in new data.
    \begin{itemize}
        \item First pattern should be detected by a model so that anomaly becomes meaningful? What about statistical models that calculate the distance between two set of points?
        \item Can a pattern be determined by one time experience?
    \end{itemize}
    \begin{itemize}
        \item What is the difference between pattern recognition in sensory data and anomaly detection? Anomalies are realized from a model that distinguishes a pattern from a non-pattern
        \item Should we expect any continuous anomaly over continuous time interval a model for pattern recognition?
    \end{itemize}

\subparagraph{Intelligence}

% Source: from_phd_review/robotics/robot/parts/mind/intelligence/integrated-robotic-intelligence.tex

\ProjectDeepHeading{2}{Unifying Frameworks for Perception, Action, and Learning in Robotics or Integrated Robotic Intelligence}
This chapter explores unifying frameworks that integrate perception, action, and learning in robotics. These frameworks range from foundational approaches like reinforcement learning to advanced methodologies such as neuro-symbolic integration and embodied intelligence. Each framework contributes uniquely to the development of robots capable of adaptive and autonomous behavior.

\ProjectDeepHeading{3}{Overview of Frameworks}
    \ProjectDeepHeading{4}{Reinforcement Learning and Hierarchical Extensions}
    Reinforcement Learning (RL) provides a framework for robots to learn through interaction with their environment, optimizing cumulative rewards. Hierarchical Reinforcement Learning (HRL) extends RL by decomposing missions into manageable submissions, simplifying the learning process. For example, a humanoid robot missioned with cleaning a room can use HRL to break the mission into locating objects, picking them up, and organizing them. Drones applying RL optimize flight paths to maximize coverage while minimizing energy usage.

    \ProjectDeepHeading{4}{Sensorimotor and Embodied Intelligence Frameworks}
    Sensorimotor learning focuses on coupling sensory input and motor output, enabling robots to refine actions based on feedback. Embodied intelligence expands this by emphasizing the robot's physical interaction with the environment. For instance, a soft robot navigating uneven terrain adjusts its movements based on tactile feedback, demonstrating the integration of physical and sensory adaptation. These frameworks are essential for missions like grasping, balancing, and navigating complex terrains.


    \ProjectDeepHeading{4}{Bayesian Brain and Active Inference}
    The Bayesian brain framework and active inference approach integrate perception, action, and learning through generative models that minimize prediction errors. Robots continuously update their beliefs and act to reduce discrepancies between predictions and observations. For example, a self-navigating robot adjusts its trajectory based on predicted and observed movements of nearby objects. This framework provides a robust basis for autonomous adaptation. 

    \ProjectDeepHeading{4}{Damasio’s Emotional and Embodied Cognition}
    Antonio Damasio's work highlights the role of emotions and bodily states in guiding reasoning and decision-making. The somatic marker hypothesis suggests that emotions derived from bodily signals can inform intelligent behavior in uncertain scenarios. Robots inspired by this framework simulate emotional reasoning to enhance human-robot interaction and decision-making, such as adapting responses based on perceived human emotions. 

    \ProjectDeepHeading{4}{Haykin’s Cognitive Dynamic Systems}
    Simon Haykin’s cognitive dynamic systems integrate neural networks, signal processing, and real-time adaptation to address uncertainty and noise in dynamic environments. Applications include multi-agent communication and adaptive control, where robots dynamically adjust behavior based on environmental conditions. For example, drones utilize Haykin’s principles for decentralized communication and navigation in complex terrains. 

    \ProjectDeepHeading{4}{End-to-End Deep Learning Frameworks}
    End-to-end deep learning frameworks map raw sensory inputs directly to motor outputs using deep neural networks. This approach eliminates the need for intermediate modules, enabling missions like autonomous driving and robotic navigation. For instance, drones process video feeds in real time to navigate complex environments.

    \ProjectDeepHeading{4}{Neuro-Symbolic Integration Frameworks}
    Neuro-symbolic frameworks combine neural networks for perception with symbolic reasoning for decision-making. Robots using these frameworks recognize objects and follow logical rules to complete missions. For instance, a robot assembling a mechanical part uses neural vision for identification and symbolic reasoning for correct assembly.

    \ProjectDeepHeading{4}{Predictive Coding Frameworks}
    Predictive coding minimizes discrepancies between predicted and observed sensory inputs by continuously updating internal models. This framework enables robots to act proactively and adapt dynamically. For example, a balancing robot anticipates shifts in its center of gravity and adjusts its posture to maintain stability. 

    \ProjectDeepHeading{4}{Multi-Agent Reinforcement Learning (MARL)}
    Multi-Agent Reinforcement Learning (MARL) extends RL to scenarios involving multiple robots collaborating or competing to achieve missions. For example, a swarm of drones collaboratively mapping a disaster area uses MARL to share data and optimize coverage while avoiding overlaps or collisions. This framework integrates shared observations, policy optimization, and coordinated actions.

    \ProjectDeepHeading{4}{Integration of Frameworks}
    While each framework contributes uniquely, their integration can lead to more robust and versatile robotic systems. For instance, predictive coding can enhance neuro-symbolic reasoning by providing anticipatory models, while embodied intelligence can ground deep learning models in real-world physical interactions. Exploring synergies among these frameworks opens new avenues for research in unified robotic intelligence.

% Source: from_phd_review/robotics/from-awareness-to-consciousness.tex

\subsubsection{From Awareness to Consciousness}
\paragraph{(Sensory) Data} Data refers to raw, unprocessed facts or measurements that lack context or interpretation. It is the basic building block for deriving meaning but does not convey any significance on its own.

\paragraph{Awareness}
Awareness refers to the ability to perceive, sense, or be conscious of one’s surroundings, internal states, or external stimuli. It is a fundamental cognitive capability that allows an entity (human, animal, or even machine) to respond to changes in its environment or internal state.

Does awareness need sensors? For example a robot with all disconnected sensors by only reviewing its memory patterns is aware? Does the literature support internal and external awareness?




\paragraph{Intelligence}  is the ability to learn, reason, adapt, and solve problems. It includes applying knowledge to achieve missions. What is the most basic form of intelligence?



\paragraph{Cognition}
refers to the processes of acquiring, processing, storing, and using information. It includes perception, memory, reasoning, decision-making, problem-solving, and learning. Cognition in robots is artificial and computational. It involves algorithms and models that enable perception, learning, and decision-making. A cognitive robot integrates sensors, learning systems, reasoning mechanisms, and actuators to perceive, understand, adapt, and interact effectively with its environment. The cognitive architecture ties these components together, allowing the robot to exhibit intelligent and adaptive behavior akin to human cognition but without subjective awareness. Here are the components of a cognitive robot:
    \begin{tabular}{|l|l|l|}
        \hline
        \textbf{Component} & \textbf{Purpose} & \textbf{Examples} \\ \hline
        \textbf{Sensors} & Collect data from the environment. & Cameras, LiDAR, microphones, touch sensors. \\ \hline
        \textbf{Perception System} & Process and interpret sensor data. & Object detection, SLAM, speech recognition. \\ \hline
        \textbf{Memory and Knowledge} & Store past experiences and learned data. & Semantic graphs, long-term memory. \\ \hline
        \textbf{Learning System} & Learn and adapt from experience. & Reinforcement learning, supervised learning. \\ \hline
        \textbf{Decision-Making} & Make optimal choices for actions. & Path plan, probabilistic reasoning. \\ \hline
        \textbf{Actuators and Motion} & Execute physical actions. & Motors, robotic arms, motion planners. \\ \hline
        \textbf{Communication} & Interact with humans or systems. & NLP, gesture recognition, HRI frameworks. \\ \hline
        \textbf{Feedback and Adaptation} & Improve behavior using feedback. & Error correction, real-time adaptation. \\ \hline
        \textbf{Cognitive Architectures} & Integrate all components cohesively. & SOAR, ACT-R, ROS, BICA. \\ \hline
        \textbf{Ethics and Safety Systems} & Ensure safe and ethical operation. & Collision avoidance, failsafe mechanisms. \\ \hline
    \end{tabular}



    \subparagraph{Cognitive model} \url{https://en.wikipedia.org/wiki/Cognitive_model\#Dynamical_systems}{url}
        Cognitive model is a computational or theoretical framework that describes and simulates cognitive processes (e.g., perception, attention, memory, reasoning, and decision-making).
        \begin{itemize}
            \item mental model: A mental model is a cognitive representation or internal simulation that individuals form to understand, predict, and interact with the world around them. It allows people to reason about how things work and to solve problems. A mental model is a subset of a broader cognitive model. Mental models describe an individual's internal understanding, while cognitive models explain the processes that lead to this understanding.
        \end{itemize}
    \subparagraph{Dynamic cognitive systems and robots}

    \subparagraph{Cognitive architectures}
        \textbf{Cognitive Architectures:} Cognitive architectures are computational frameworks that model and simulate the structures and processes underlying human cognition. They provide a unified framework for integrating components like memory, reasoning, perception, and learning to mimic human-like intelligence. By organizing cognitive processes into modular structures, these architectures enable the study of general cognitive abilities and their application in artificial systems. Many cognitive architectures also emphasize the interaction between symbolic and sub-symbolic processes, ensuring adaptability in dynamic environments. Their development bridges theoretical cognitive science and practical artificial intelligence research. 

        \ProjectDeepHeading{1}{Review Papers on Cognitive Architectures:}

            \textbf{"A Review of 40 Years of Cognitive Architecture Research: Core Cognitive Abilities and Practical Applications"} by Iuliia Kotseruba and John K. Tsotsos provides an extensive overview of cognitive architecture research, discussing core cognitive abilities and various applications. 

            \textbf{"Progress and Challenges in Research on Cognitive Architectures"} by Pat Langley reviews the notion of cognitive architectures and recurring themes in their study, highlighting both advancements and ongoing challenges in the field. 

            \textbf{"Cognitive Architectures and Autonomy: A Comparative Review"} by Kristinn R. Thórisson and Thórhallur E. Rögnvaldsson examines systems and architectures designed to address integrated skills, discussing principles and features that contribute to autonomy in intelligent systems. 



        \ProjectDeepHeading{1}{Most famous cognitive architectures}
            \textbf{ACT-R:} ACT-R (Adaptive Control of Thought-Rational) is a cognitive architecture designed to simulate human cognitive processes. It integrates declarative and procedural memory, enabling systems to retrieve relevant information based on pattern-matching mechanisms. For instance, in robotics, ACT-R is applied to model human-like reasoning and memory retrieval during missions such as human-robot interactions. This architecture emphasizes modularity, allowing cognitive processes like memory recall and learning to be closely aligned with human cognition.

            \textbf{Soar:} Soar is a cognitive architecture that focuses on learning, problem-solving, and decision-making by using long-term memory retrieval. It enables agents to recall stored plans or strategies when encountering missions similar to past experiences. In robotics, Soar facilitates navigation and adaptive behaviors by retrieving relevant memory traces during dynamic problem-solving. Its ability to learn from past outcomes makes it a robust framework for modeling cognitive flexibility in artificial systems.

            \textbf{CLARION:} CLARION (Connectionist Learning with Adaptive Rule Induction ON-line) is a hybrid cognitive architecture combining symbolic reasoning with sub-symbolic neural processes. This architecture supports memory retrieval by associating high-level rules with underlying distributed representations. In robotic systems, CLARION is particularly useful for integrating reasoning and action, such as adapting to new environments based on learned associations and generalizations. Its dual-process approach ensures flexibility and robustness in memory-dependent missions.

            \textbf{LIDA:} LIDA (Learning Intelligent Decision Agent) is a biologically inspired cognitive architecture that emphasizes cognitive cycles and attention mechanisms. It models human-like perception, decision-making, and learning through its multiple memory systems, including sensory, episodic, and procedural memory. LIDA has been applied in multi-robot systems to enable distributed decision-making and collaborative learning, making it well-suited for dynamic and adaptive environments. 

            \textbf{OpenCog:} OpenCog is an open-source cognitive architecture aimed at achieving artificial general intelligence (AGI). It uses knowledge graphs and probabilistic reasoning to represent and process complex knowledge structures. In multi-robot systems, OpenCog enables collaborative reasoning and shared decision-making, allowing robots to collectively solve missions by leveraging distributed knowledge. \url{https://wiki.opencog.org/w/The_Open_Cognition_Project}

            \textbf{SPAUN:} SPAUN (Semantic Pointer Architecture Unified Network) is a large-scale brain-inspired model that simulates cognitive processes such as visual recognition, decision-making, and learning. It uses spiking neural networks to perform various missions, making it particularly useful for simulating human-like cognitive behaviors in robotic systems.

            \textbf{NARS:} NARS (Non-Axiomatic Reasoning System) is a cognitive architecture designed for reasoning and decision-making under uncertainty. It adapts to changing environments by continuously learning from incomplete and dynamic knowledge. NARS has been employed in real-time robotics, where adaptive and flexible reasoning is critical. 

            \textbf{BECCA:} BECCA (Brain-Emulating Cognition and Control Architecture) focuses on unsupervised learning and adaptive control by mimicking biological learning processes. It enables robots to autonomously adapt to new missions and environments, making it suitable for applications requiring continuous learning and behavioral flexibility.

            \begin{table}[h!]
                \centering
                \begin{tabular}{|l|l|l|l|}
                    \hline
                    \textbf{Cognitive Architecture} & \textbf{Type} & \textbf{Key Feature} & \textbf{Application in Robotics} \\ \hline
                    ACT-R & Symbolic & Human-like reasoning and memory retrieval & Human-robot interaction  \\ \hline
                    Soar & Symbolic & Problem-solving and decision-making & Adaptive navigation  \\ \hline
                    CLARION & Hybrid & Combines symbolic reasoning with subsymbolic learning & Dynamic adaptation  \\ \hline
                    LIDA & Hybrid & Cognitive cycles and distributed learning & Distributed decision-making  \\ \hline
                    OpenCog & Subsymbolic & Knowledge graphs and probabilistic reasoning & Collaborative reasoning  \\ \hline
                    SPAUN & Neuromorphic & Brain-inspired spiking neural network & Cognitive mission simulation  \\ \hline
                    NARS & Subsymbolic & Adaptive reasoning under uncertainty & Real-time decision-making  \\ \hline
                    BECCA & Neuromorphic & Unsupervised learning and adaptive control & Adaptive control in robotics  \\ \hline
                \end{tabular}
                \caption{Summary of Cognitive Architectures for Robotics}
            \end{table}




    \subparagraph{Embodied Cognition}
        \todo{Embodied agent}

    \subparagraph{Attention}
        % I checked attention is a sesnsory of cognition
        \todo{Write about attention as a cognition ability}



\paragraph{Consciousness} is the state of being aware of one’s existence, thoughts, surroundings, and self. It involves subjective experience and introspection. It is a level higher than awareness and in most literature is claimed to be not reachable by robots. Consequently, self-consciousness is also a level higher than self-awareness. A robot can never experience what a burn feel likes.
% \todo{check this paper:  \url{https://www.researchgate.net/publication/360659133_Robots_and_Machine_Consciousness}}

% Source: from_phd_review/robotics/environment.tex

\subsubsection{Environment}

\paragraph{Describing formats}
    \ProjectDeepHeading{2}{Occipancy grid}
        Occupied - Free  - unknown
    
    
    \ProjectDeepHeading{2}{GEOTIFF}
        Pixel + 3D GPS
        
    \ProjectDeepHeading{2}{DEM - Digital Elevation Model}
        in shape file format (.shp) formed of GeoTIFF or ASCII grid
        
    \ProjectDeepHeading{2}{GeoJason}
        vector format for web to speed up
        
    \ProjectDeepHeading{2}{KML/KMZ for Google maps}
        For google maps


\paragraph{From robot perspective}
    \subparagraph{SLAM}
        Simultaneous Localization And Mapping
        \ProjectDeepHeading{1}{LIDAR based}
            \ProjectDeepHeading{1}{Hector SLAM}
        \ProjectDeepHeading{1}{Output}
            \ProjectDeepHeading{2}{Grid}

\paragraph{From real point of view for testing}
    This is just for estimating the fault. For examzple we have the real distance of a robot from a wall and we have the estimation of the robot from the wall. Now we can estimate the distance between what the robot thinks of the environment and what it really is.


\paragraph{Gazebo}
Simulation of physical world for robots

% Source: from_phd_review/robotics/control-theory.tex

\subsubsection{Control theory}
\textbf{Feedback theory}
Feedback theory is a foundational concept in control theory that explains how systems regulate their behavior through feedback loops. It provides a mathematical framework for monitoring outputs, comparing them to desired targets, and making adjustments to minimize errors. This theory ensures stability, robustness, and adaptability in dynamic systems, allowing them to maintain equilibrium and respond to disturbances effectively. Feedback theory is widely applied in robotics for missions such as stabilization and trajectory control, where sensors provide real-time data to refine motor commands. Its principles are crucial for designing both biological and artificial systems that integrate perception, action, and learning.


\paragraph{Overview of Control Systems}

\subparagraph{Open-Loop Control Systems}
Open-loop control systems operate without feedback, relying on predefined commands to control the robot's actions. These systems are straightforward to implement and suitable for predictable missions where system dynamics and disturbances are negligible. For example, in robotics, an open-loop system might control a conveyor belt or move a robotic arm to a fixed position without considering external influences. While cost-effective, these systems lack the ability to correct errors caused by changes in the environment or system inaccuracies, making them unsuitable for dynamic missions.

\subparagraph{Closed-Loop (Feedback) Control Systems}
Closed-loop control systems use feedback from sensors to adjust the robot's actions in real-time. This feedback loop ensures that the desired outcomes are achieved even in the presence of disturbances or uncertainties. For instance, a balancing robot like a Segway uses feedback from gyroscopes to maintain its balance, or a drone adjusts its altitude based on barometric pressure readings. These systems are accurate and robust but require sensors and additional computational power, increasing system complexity and cost.

\subparagraph{Proportional-Integral-Derivative (PID) Control}
PID control is a fundamental closed-loop control algorithm that adjusts system outputs based on the proportional, integral, and derivative of the error signal. It is widely used for missions requiring precise control, such as maintaining the position or speed of robotic arms or ensuring drones hover steadily in midair. While PID controllers are easy to implement and tune for linear systems, they can struggle with highly dynamic or nonlinear systems where advanced control techniques may be required.

\subparagraph{Adaptive Control}
Adaptive control dynamically adjusts its parameters to maintain optimal performance under changing conditions. This type of control is ideal for robots operating in unpredictable environments, such as drones dealing with fluctuating wind conditions. For example, an adaptive controller can modify a drone's flight parameters in real-time to maintain stability and performance as the payload or environmental factors change. While powerful, designing and implementing adaptive control systems can be complex and computationally demanding.

\subparagraph{Optimal Control}
Optimal control focuses on minimizing or maximizing a specific performance criterion, such as energy consumption, time, or trajectory accuracy, over a given system trajectory. In robotics, optimal control might be used to plan efficient paths for robotic arms or drones to minimize travel time and energy use. For example, drones missioned with area surveillance can optimize their flight paths to maximize coverage with minimal battery consumption. While achieving high performance, these systems require accurate models and significant computational resources.

\subparagraph{Model Predictive Control (MPC)}
Model Predictive Control predicts future states of the system using a model and optimizes control inputs over a finite time horizon to meet constraints and objectives. This approach is particularly useful for complex and multivariable robotic systems. For instance, autonomous drones can use MPC to plan collision-free trajectories in dynamic environments, taking into account sensor data from LiDAR and cameras. Although MPC provides excellent performance in constrained systems, its computational requirements can be a limitation for fast-moving robots.

\subparagraph{Robust Control}
Robust control ensures system stability and performance even in the presence of uncertainties and disturbances. This type of control is commonly used in industrial robots or drones operating under varying environmental conditions, such as strong winds. For example, a robust controller for a drone would maintain stability and precise operation despite model inaccuracies or sensor noise. While highly reliable, robust control systems often involve conservative designs and require detailed modeling efforts.

\subparagraph{Nonlinear Control}
Nonlinear control addresses systems where linear assumptions do not hold, such as robotic arms with flexible joints or drones performing aerobatic maneuvers. Nonlinear controllers use specialized techniques to handle complex dynamics, enabling robots to operate effectively in challenging conditions. For instance, a nonlinear controller might stabilize a humanoid robot walking on uneven terrain or a drone performing tight turns at high speeds. Designing nonlinear control systems is mathematically challenging but essential for advanced robotics.

\subparagraph{Learning-Based Control}
Learning-based control integrates machine learning algorithms to optimize control strategies based on data or experience. Robots, such as drones, can learn to navigate dynamic environments by adapting their flight controllers using data from sensors like IMU, GPS, and cameras. For example, a drone might learn to avoid obstacles in a cluttered environment by combining reinforcement learning with traditional control techniques. This approach is highly effective for complex and evolving missions but requires extensive training data and computational power.

\subparagraph{Hybrid Control Systems}
Hybrid control systems combine continuous control methods with discrete decision-making processes, such as finite state machines. These systems are ideal for robots with mixed dynamics, such as drones switching between hover and travel modes or robotic arms transitioning between different operational states. For example, a hybrid controller could manage a drone's transition from autonomous navigation to manual control in emergency scenarios. While versatile, hybrid systems require careful integration of discrete and continuous control elements.

\subparagraph{H-Infinity Control}
H-Infinity control is a robust control approach that minimizes the worst-case gain of disturbances in the system. This method ensures reliable operation under extreme uncertainties, making it ideal for high-stakes applications such as aerospace robotics. For instance, an H-Infinity controller might stabilize a space robot in the presence of unpredictable gravitational forces or disturbances. While effective in maintaining stability, this approach requires advanced mathematical design and analysis.

\subparagraph{Fuzzy Logic Control}
Fuzzy logic control uses degrees of truth rather than binary logic to handle uncertainty and imprecision in control systems. This method is well-suited for applications where precise models are unavailable. For example, fuzzy logic can control a robotic vacuum cleaner navigating cluttered rooms or a drone making smooth altitude adjustments. While intuitive and effective in many scenarios, fuzzy logic systems can sometimes lack precision compared to traditional control methods.

\paragraph{Questions}

Are there neural control theories?

\section{Mrs}

% Source: from_phd_review/mrs/mrs.tex

\subsubsection{Collectiveness}


\paragraph{The difference between an MRS and SRS}
A system can be classified as a \textit{Multi-Robot System (MRS)} if it meets a set of necessary and sufficient conditions. These conditions ensure that the system involves multiple autonomous agents interact and collaborating in a decentralized manner.

\subparagraph{Necessary and Sufficient Conditions}
Combining the necessary and sufficient conditions, a system qualifies as an MRS if it meets the following criteria:

\begin{enumerate}
    \item The system must consist of two or more robots capable of independent decision-making and act on their own.
    \item The system must lack a centralized decision-maker, with control distributed across the robots. 
    \item The robots must interact through communication or observation to achieve shared missions.
    \item The system must distribute missions across multiple robots, with each robot contributing to the mission.
    \item The robots must be able to learn independently and adapt to their environment and each other. 
    \item There must be mechanisms for coordinating missions, negotiating roles, or sharing information.
\end{enumerate}


\paragraph{Different multi robot systems}

\textbf{Heterogeneous MRS}

\textbf{Homogeneous MRS}

\textbf{Swarms} \todo{How to choose who are your neighbors?}

\textbf{Swarmnoids}

\todo{Are there more categories?}


\paragraph{Different forms of relationship between multiple robots}

\textbf{Cooperation} refers to collaborative interaction between robots to achieve a common mission that is difficult or impossible to achieve individually. In cooperative scenarios, robots share information, resources, and missions to improve efficiency. For instance, robots in search and rescue missions can cover larger areas by cooperating, thus increasing the chances of finding victims faster. Cooperation requires mission allocation, seamless communication, and synchronization of actions.

\textbf{Coordination} focuses on synchronizing the actions of multiple robots to avoid conflicts and optimize resource usage. Coordination ensures that robots operating in the same environment do not interfere with each other’s missions. For example, in warehouse automation, robots coordinate their movements to avoid collisions while picking and delivering items. Effective coordination mechanisms are essential to prevent resource conflicts and ensure scalability in large robot fleets.

\textbf{Competition} arises when robots compete for limited resources or missions. This adversarial interaction is common in game-theoretic approaches where each robot tries to maximize its individual performance. For instance, in robot soccer, teams of robots compete to score missions while blocking opponents. Market-based mission allocation systems also incorporate competition, where robots bid for missions, and the most efficient robot wins the mission.

\textbf{Collaboration} is a higher-level form of cooperation where robots adapt their behavior dynamically to work together toward a shared mission. Collaboration requires cognitive capabilities, such as reasoning and decision-making, to adjust strategies based on the environment. For instance, robots collaborating in construction missions adjust their roles in real-time based on progress and obstacles. Collaborative robots can also work alongside humans in manufacturing plants by adapting to human actions.

\textbf{Swarm Behavior} refers to emergent collective behavior in multi-robot systems that is inspired by biological systems such as ant colonies and bee swarms. In swarm robotics, each robot follows simple local rules, and complex global behavior emerges without central control. This approach is particularly useful in missions like environmental monitoring and foraging, where robots spread out to cover large areas efficiently. Swarm behavior is robust and scalable but requires careful design of local interaction rules.

\textbf{Negotiation} occurs when robots interact to resolve conflicts or agree on mission assignments. In negotiation, robots exchange proposals and counterproposals until they reach a mutually acceptable agreement. This type of interaction is crucial in resource-constrained environments, where robots need to negotiate access to limited charging stations or mission priorities. Effective negotiation protocols ensure fairness and optimize mission allocation in multi-robot teams.

\textbf{Role Assignment} focuses on assigning specific roles or missions to different robots based on their capabilities and the mission requirements. Role assignment ensures efficiency and mission specialization in complex multi-robot systems. For example, in search and rescue missions, one robot might focus on mapping the area while another searches for victims. Dynamic role assignment allows robots to switch roles based on changing conditions and mission demands.

\textbf{Social Interaction} involves robots engaging in meaningful communication with humans or other robots in a socially appropriate way. This type of interaction is essential for human-robot collaboration in areas like healthcare and customer service. Social robots are designed to recognize social norms and respond to human emotions appropriately. They use verbal and non-verbal cues to enhance interaction quality and build trust with users. 


\paragraph{Collective learning techniques in robots}
\textbf{Centralized learning}

\textbf{Decentralized learning}

\textbf{Hybrid learning}

\textbf{Cooperative Imitation Learning/Multi-Agent Imitation Learning}


\paragraph{Communication}


\paragraph{Scalability}


\paragraph{Anomaly detection in MRS}


\paragraph{Question}
\todo{Maybe the difference between mrs and srs is cognition of the boundaries of effect}
\todo{What is Coalition making?}
\todo{Consensus building}
\todo{Flocking}
\todo{Distributed cognition}
\todo{What is swarm intelligence and how does it apply to robots?}
\todo{Are there relationships between sociology and MRS?}
\todo{What are different forms of consensus making in an MRS?}
\todo{Language grounding in MRS}

% Source: from_phd_review/mrs/collective-self.tex

\subsubsection{Collective-self}
Imagine the biggest population of robots in cosmus. Collective is the ability of a robot to understand the boundry of the largest number of robots that their actions are either bounded to it or its actions are bounded to them. This ability already exist in human sociaety.


\paragraph{Social-awareness}
Social awareness refers to an robot's ability to recognize the presence, actions, and roles of other agents in its environment and adjust its behavior accordingly. It involves understanding that other agents are distinct entities with their own intentions and capabilities. In multi-robot systems (MRS), social awareness enables robots to coordinate actions, avoid conflicts, and engage in collaborative missions. It is essential for human-robot interaction and team-based robotics, as it allows robots to interact meaningfully with humans and other robots by adapting to their actions and roles.

Self-awareness is a prerequisite for social awareness because an agent must first understand itself before it can recognize and interact with others. Self-awareness involves identifying one’s own body, actions, and internal states, which is essential for distinguishing self-generated changes from external changes caused by other agents or the environment. Without this distinction, an agent would struggle to recognize which actions it is responsible for and which actions are caused by others, leading to confusion in social interactions. For a robot to adapt its behavior in response to another robot or human, it must first know its own capabilities and limitations. Therefore, social awareness builds upon self-awareness by extending the robot’s understanding from itself to its interactions with others, enabling effective coordination and cooperation in multi-agent systems.



\paragraph{Collective awareness}
Collective awareness refers to an robot's ability to understand that it is part of a group working toward a shared mission and to adjust its behavior to optimize the group's overall performance. It goes beyond recognizing individual agents to encompass an understanding of group dynamics, shared resources, and collective missions. In MRS, collective awareness is crucial for missions like swarm robotics, where robots need to coordinate without central control to achieve emergent behaviors that maximize group efficiency. This form of awareness allows robots to distribute missions, avoid redundancies, and adapt to changes in the group’s composition.

Social awareness is a prerequisite for collective awareness because recognizing individual agents and their roles is essential before understanding the group as a whole. Social awareness enables an agent to identify other agents, perceive their actions, and understand their missions and intentions. Without this ability to recognize and respond to individual interactions, an agent would be unable to coordinate missions or participate in shared group missions. Collective awareness builds upon social awareness by requiring an agent to integrate its knowledge of individual agents into an understanding of group dynamics and shared objectives. Therefore, without first understanding who is part of the group and what each agent is doing, a robot cannot optimize group performance or contribute to collective success.
\begin{table}[h!]
    \centering
    \begin{tabular}{|l|l|l|}
        \hline
        \textbf{Aspect}            & \textbf{Social Awareness}                                   & \textbf{Collective Awareness}                            \\ \hline
        \textbf{Focus}             & Individual interactions                                     & Group-level understanding                                \\ \hline
        \textbf{Main Mission}         & Recognizing other agents and their roles                    & Understanding shared missions and collective missions          \\ \hline
        \textbf{Adaptation}        & Adapting to the actions of individual agents                & Adapting to the needs of the group as a whole            \\ \hline
        \textbf{Example}           & Waiting for another robot to finish using a shared tool     & Adjusting search patterns based on the group’s coverage  \\ \hline
        \textbf{Type of Awareness} & Interpersonal awareness                                     & Group-level awareness                                    \\ \hline
        \textbf{Relevance in MRS}  & Essential for avoiding conflicts and improving coordination     & Essential for achieving group efficiency and cooperation         \\ \hline
    \end{tabular}
    \caption{Comparison Between Social Awareness and Collective Awareness}
\end{table}


\paragraph{Social cognition}
Social cognition in multi-robot systems (MRS) refers to the ability of robots to understand and respond to individual agents’ actions, roles, and intentions. This involves recognizing other robots as distinct entities, predicting their behaviors, and adapting accordingly. Social cognition enables role recognition, conflict resolution, and collaboration in shared environments. For example, in a warehouse, robots use social cognition to avoid collisions by predicting each other’s paths and adjusting their trajectories. Inspired by human social cognition, this concept incorporates elements like theory of mind, which allows robots to infer the missions and intentions of others. Social learning is another critical aspect, enabling robots to learn new missions by observing the actions of peers. Effective social cognition is essential for missions that require close interaction between robots or between robots and humans. It also facilitates human-robot collaboration by allowing robots to adapt to human commands or gestures. However, implementing social cognition in robotics faces challenges like intention prediction and communication limitations.


\paragraph{Collective cognition}
Collective cognition refers to the ability of a group of robots to process information collectively, make joint decisions, and adapt as a team to achieve shared missions. It goes beyond individual interactions to focus on the emergent group behavior that arises from decentralized coordination. In MRS, collective cognition is exemplified by swarm robotics, where robots collaboratively explore an environment or distribute missions. Shared memory systems and distributed decision-making are fundamental elements of collective cognition, allowing robots to maintain a unified understanding of their environment. For instance, drones in a disaster zone use collective cognition to collectively map an area, ensuring coverage without redundancy. Collective cognition is driven by local interactions among robots, leading to global behaviors like efficient mission allocation or cooperative transport. This concept is inspired by biological systems like ant colonies and bird flocks, where individuals follow simple rules to achieve complex group objectives. While collective cognition enables scalable and robust multi-robot systems, challenges like communication overhead and consensus mechanisms remain significant.


\paragraph{The role of anomaly detection in collective awareness, social awareness , collective cognition and }


\paragraph{self, others and self}
\textbf{Questions}
\begin{itemize}
    \item Is the notion of being multiple true?
    \item Does collective self-awareness exist?
    \item How can two agents agree that they are talking about the same concept?
    \item How does the presence of others help with self-awareness?
    \item Is there really a difference between personal self-awareness and collective self-awareness?
    \begin{itemize}
        \item What is the relation between pattern recognition and anomaly detection
    \end{itemize}
\end{itemize}


\paragraph{Questions}
\todo{How do multiple robots know that they are working together?}
\todo{How does an MRS robots undetstand that they are different than another single or mrs?}

% Source: from_phd_review/prediction.tex

\subsection{Prediction}

% Source: from_phd_review/robotics/robot.tex

\subsubsection{Robots}
A robot is a programmable, multifunctional machine designed to perform missions by sense its environment, processing information, and taking action either autonomously or semi-autonomously.



\paragraph{Sensors}
    \subparagraph{IMU}
        IMUs come with 6 axis or 9 axis. 6 axis includes accelerometer and gyroscope and 9 axis also includes magnometer
        
        \ProjectDeepHeading{1}{Accelometer}
            To measure the x,y,z acceleration
        
        \ProjectDeepHeading{1}{Giroscope}
            to measure rotational rate, eitrher (roll, yaw, pitch) or quaternion
        
        \ProjectDeepHeading{1}{Magnetometer}
            It measures earth magnetic field. the output is a threed number
        
        \ProjectDeepHeading{1}{GPS}
        
        \ProjectDeepHeading{1}{LIDAR}
        
\paragraph{Sense Plan Action loop}
    It is also called Perception-Planning-Control
    \begin{figure}[H]
        \centering
        \ProjectSafeIncludeGraphics[width=0.8\linewidth]{../assets/sense_plan_act.png}
        \label{fig:my_png}
    \end{figure}
    
    \subparagraph{Sense}
        Perception has four layers
        \ProjectDeepHeading{1}{Sensing}
            Raw sensory data
        \ProjectDeepHeading{1}{Preprocessing}
        \ProjectDeepHeading{1}{Odometry and then State estimation}
            Odometry is the process of estimating a robot’s change in position (translation) and orientation (rotation) over time, using motion sensors.
        \ProjectDeepHeading{1}{Localication and mapping}
            \ProjectDeepHeading{2}{SLAM}
                stands for simultanious localization and mapping. This is very related to self-awareness in a robot.
                    \ProjectDeepHeading{3}{Output}
                    SLAM output is  something like an ocuupancy grid
    
    \subparagraph{Plan}
        \ProjectDeepHeading{2}{Mission} Mission is an input to plan
    \subparagraph{Action}
        Executes motor commands to follow the planned trajectory
            

\paragraph{Action-Mission}
    A robot perform a set of actions to achieve a (set of) missions
    \ProjectDeepHeading{2}{Action}
        Action is an assesable long running missions such as goto point (x,y,z). Mavlink defines actions.
    \ProjectDeepHeading{2}{Mission}
        A pose in which a robot should try to achieve
    \ProjectDeepHeading{2}{Commands}
        instructions sent to actuators such as increase rotation speed of rotor 3 in a quadcopter

\paragraph{The main components of a robot}:
The main components of a robot can be divided into five essential categories:

    \subparagraph{Physical components}
        \textbf{Sensors}: Allow robots to perceive both their environment and internal states, using sensors like cameras and LiDAR or IMUs.

        \textbf{Actuators}: Enable the robot to move or perform missions by converting control signals into physical motion.
        Examples include motors, servos, robotic arms, and wheels.

        \textbf{Power Supply}: Provides the energy needed to operate the sensors, actuators, and processing unit.
        Examples include batteries, solar cells, and fuel cells.

        \textbf{Mechanical Structure}: The physical frame or body of the robot that supports and integrates the other components.
        Examples include wheels, legs, manipulators, and drone frames.

    \subparagraph{Virtual components:}
        \textbf{Controller (Processing Unit)}: The "brain" of the robot that processes sensor inputs, makes decisions, and sends commands to actuators.
        Examples include microcontrollers, processors, and embedded systems.

        \textbf{Missions} A mission in robotics refers to a desired outcome or end state that a robot aims to achieve through its actions. Unlike missions, missions are typically long-term and result-oriented, requiring the robot to plan, adapt, and choose appropriate missions to reach the desired objective. Missions can be static or dynamic, depending on environmental changes or new information. In autonomous systems, missions often drive decision-making by determining which missions to prioritize and how to allocate resources. In cognitive architectures, mission management is critical for enabling mission-directed behavior and learning from experience.
        




\paragraph{Different robots}

    \subparagraph{Non-Intelligent Robots:} These robots operate based on pre-programmed instructions with limited or no adaptability to changing environments. Examples include automated assembly-line robots or simple vacuum robots. 

    \subparagraph{Intelligent Robots:} These robots are capable of perceiving their environment, reasoning, and making decisions autonomously, often using AI techniques like machine learning. Examples include humanoid robots, autonomous vehicles, or drones with adaptive navigation systems. 

        \textbf{Cognitive Robots:} A subset of intelligent robots, cognitive robots are designed to mimic human-like cognitive processes, such as reasoning, problem-solving, memory, and understanding intentions. They focus on emulating human intelligence and are often used for missions requiring higher-order reasoning and natural interaction, such as social robots like Pepper.

        \textbf{Cognitive dynamic robots}

        \textbf{Concinnous robots}
        Probably never achievable. Pain must reduce computation resources in robot.They probably can never answer the question "What burning feels like" from recognized patterns in their in sensory data in their memory. Another aspect of pain might be messiness of words or meanings as sensory patterns in their memory.

\paragraph{Feedback system}

\paragraph{Control theory}

\paragraph{Sensors}
    \begin{itemize}
        \item proprioceptive and exteroceptive sensors?
        \item sensory data
        \item Are sensory data only for state estimation ?
    \end{itemize}

\paragraph{State estimation}
    \subparagraph{Questions}
        \begin{itemize}
            \item Is relating current sensor data with previous control inputs to solve problem a kind of state estimation?
        \end{itemize}
    
\paragraph{Mission}
    \todo{This part must go under sense-plan-action cycle}
    \ProjectDeepHeading{2}{Homestatis and survival}
        Somehow must prove that without having a surviavl mission self aware ness is not possible
        
        The survival final mission is the only effort that turns a robots attention to itself.
        
        Obeying the human ruler is one surviavl skill because the human ruler may kill the robot


        
\paragraph{Simulation}
    \subparagraph{ROS ecosystem}
    
    
    \subparagraph{ROS}
    
    
    \subparagraph{MAVROS}
    A layer over ROS
    
    
    \subparagraph{MAVLink}
    The language to establish communication between
    \begin{itemize}
        \item QGroundControl and (MAVROS-MAVLink)
        \item (MAVROS-MAVLink) and (Autopilot)
    \end{itemize}
    
    
    \subparagraph{QGround Control}
    
    
    \subparagraph{Flight controller}
    The hardware board that usually containes
    \begin{itemize}
        \item IMU
        \item barometer
        \item GPS
        \item microcontroller/CPU
        \item I/O ports
    \end{itemize}
    
    Here are sime examples
    \begin{itemize}
        \item Pixhawk
        \item Cube orange
        \item Holybro
    \end{itemize}
    
    
    \subparagraph{Autopilot}
    The software
    
    \ProjectDeepHeading{1}{PX4}
    
    \ProjectDeepHeading{1}{PX4mini}
    
    \ProjectDeepHeading{1}{Ardupilot}
    
\paragraph{Quadcopter}
    \subparagraph{Homeostatis}
    This could be considered as the final mission of a robot. Even a robot which follows its human master, it is better to be regarded as something that follows its master orders in order not to be destroyed by it
    This will be very intresting. A flying quadcopter is very prone to crashing. That is why it is a very good subject for self-awarness studies

% Source: from_phd_review/robotics/self.tex

\subsubsection{Self}
\todo{Markov blanket}


\ProjectDeepHeading{2}{Why is it important that a robot knows what it is?}
    Why is it important that a robot knows what it is and how it can know what its?

\paragraph{Self}
In robotics, self is defined as the robot's ability to model, predict, and distinguish its own internal state, physical body, and actions from external influences. This "self" is achieved through mathematical models, sensor data, and computational processes.


\paragraph{Distinguishing Between Self and Non-Self in Robots}

\subparagraph{Internal Models and Forward Models}
Robots use \textbf{internal models} to predict the outcomes of their actions. By comparing predicted and observed outcomes, they can distinguish between self-generated actions and external influences.

\textbf{Mathematical Formulation:}
\begin{equation}
    \hat{x}_{t+1} = f(x_t, u_t)
\end{equation}
where:
\begin{itemize}
    \item \( \hat{x}_{t+1} \) is the predicted next state.
    \item \( x_t \) is the current state.
    \item \( u_t \) is the control input.
    \item \( f \) is the internal model function.
\end{itemize}
The error between the predicted and observed states helps identify external disturbances:
\begin{equation}
    \text{Error} = \| x_{\text{observed}} - \hat{x}_{t+1} \|
\end{equation}

\subparagraph{Artificial Immune System (AIS) Models}
Inspired by biological immune systems, AIS models distinguish self from non-self through anomaly detection.

\textbf{Mathematical Approach:}
\begin{itemize}
    \item Define a \textit{self-space} of normal behaviors.
    \item Use the \textbf{Negative Selection Algorithm} to generate detectors for anomalies.
\end{itemize}

The affinity function measures the deviation of an observation from the self-space:
\begin{equation}
    \text{Affinity} = \| x_{\text{observed}} - x_{\text{self}} \|
\end{equation}

\subparagraph{Dynamical Systems and Control Theory}
Robots model their own dynamics and detect external forces (non-self) by analyzing deviations.

\textbf{Dynamics Equation:}
\begin{equation}
    \dot{x} = Ax + Bu + d
\end{equation}
where:
\begin{itemize}
    \item \( x \) is the state variable.
    \item \( A, B \) are system matrices.
    \item \( u \) is the control input.
    \item \( d \) represents external disturbances.
\end{itemize}
Non-zero \( d \) indicates the presence of non-self influences.

\subparagraph{Sensory-Motor Contingencies (SMCs)}
SMCs describe the predictable relationship between motor actions and sensory outcomes. Self-generated actions produce expected sensory consequences.

\textbf{Mathematical Representation:}
\begin{equation}
    \hat{s}_{t+1} = g(s_t, a_t)
\end{equation}
where:
\begin{itemize}
    \item \( \hat{s}_{t+1} \) is the predicted sensory outcome.
    \item \( s_t \) is the current sensory state.
    \item \( a_t \) is the motor action.
\end{itemize}
Unpredictable sensory feedback indicates non-self behavior.

\subparagraph{Machine Learning-Based Models}
Machine learning models, such as autoencoders, detect deviations from learned self-behavior.

\textbf{Autoencoder Reconstruction Error:}
\begin{equation}
    \text{Reconstruction Error} = \| x_{\text{input}} - x_{\text{reconstructed}} \|
\end{equation}
A high reconstruction error signals non-self behavior.

\subparagraph{Bayesian Inference Models}
Bayesian models use probabilistic reasoning to classify observations as self or non-self.

\textbf{Posterior Probability:}
\begin{equation}
    P(\text{self} | \text{observation}) \propto P(\text{observation} | \text{self}) P(\text{self})
\end{equation}
Low posterior probability indicates non-self influences.

\ProjectDeepHeading{1}{Summary Table of Models}

\begin{tabular}{|l|p{6cm}|p{5cm}|}
    \hline
    \textbf{Model}              & \textbf{Core Idea}                              & \textbf{Use Case}                             \\ \hline
    Internal Models             & Predict outcomes and compare with observations. & Collision detection, motion verification.     \\ \hline
    Artificial Immune Systems   & Detect deviations from self-space.              & Fault detection, anomaly recognition.         \\ \hline
    Dynamical Systems           & Analyze external forces in system dynamics.     & Detecting external disturbances.              \\ \hline
    Sensory-Motor Contingencies & Predict sensory outcomes of motor actions.      & Body ownership, tool use.                     \\ \hline
    Machine Learning            & Learn self-behavior and detect anomalies.       & Anomaly detection, fault isolation.           \\ \hline
    Bayesian Inference          & Probabilistically classify self vs. non-self.   & Uncertainty estimation, disturbance analysis. \\ \hline
\end{tabular}

\subparagraph{Questions}
Is there a determined boundary between self and none-self?

what is the definition of Self-modeling?

if there are no external forces, then the robot can not distinguish between self and non-self


\paragraph{Self-awareness}
In robotics, self-awareness refers to a robot's ability to recognize, understand, and respond to its internal states and its role within its environment or a system. It involves functional awareness rather than subjective experience, enabling robots to monitor their own status, actions, and interactions with the external world. Self-awareness is a higher-order form of awareness. It refers to the ability to recognize oneself as an individual separate from the environment and others. Self-awareness involves not only perceiving external and internal states but also reflecting on them as being "one's own".
\begin{tabular}{|l|l|l|}
    \hline
    \textbf{Aspect}              & \textbf{Awareness}                              & \textbf{Self-Awareness}                                               \\ \hline
    \textbf{Definition}          & The ability to perceive and respond to stimuli. & The ability to reflect on oneself as distinct from the environment.   \\ \hline
    \textbf{Focus}               & External environment and internal states.       & Internal reflection and recognition of self.                          \\ \hline
    \textbf{Complexity}          & Basic and reactive.                             & Higher-order cognitive process.                                       \\ \hline
    \textbf{Recognition of Self} & Absent.                                         & Present (e.g., recognizing self in a mirror).                         \\ \hline
    \textbf{Example (Humans)}    & Hearing a sound and turning toward it.          & Realizing “I am nervous because of an upcoming test.”                 \\ \hline
    \textbf{Example (Robots)}    & Detecting an obstacle and avoiding it.          & Recognizing that a movement or failure was caused by its own actions. \\ \hline
    \textbf{Level of Cognition}  & Low-level cognition or sensory processing.      & High-level cognition involving reflection and self-modeling.          \\ \hline
\end{tabular}
\begin{itemize}
    \item Is it an agent's knowledge about its existence?

    \begin{itemize}
        \item This boundary classifies between what and what?
        \begin{itemize}
            \item sensory readings and possible states
            \item Are there a trainable and improvable models to classify between self and non-self?
        \end{itemize}
    \end{itemize}
\end{itemize}

\subparagraph{Different levels of self-awareness}
\begin{tabular}{|l|l|l|l|}
    \hline
    \textbf{Level}               & \textbf{Description}                 & \textbf{Examples in Humans}    & \textbf{Examples in AI/Robots} \\ \hline
    Zero-Level                   & No self-perception.                  & None.                          & Reflex-based systems.          \\ \hline
    Physical Self-Awareness      & Recognizing one’s body and position. & Proprioception.                & Collision avoidance in robots. \\ \hline
    Situational Self-Awareness   & Awareness of actions in context.     & Adjusting behavior in traffic. & Path plan in environments. \\ \hline
    Mirror Self-Awareness        & Recognizing oneself in a reflection. & Seeing oneself in a mirror. & Basic visual recognition. \\ \hline
    Extended Self-Awareness      & Awareness of self across time.       & Reflecting on events.          & Tracking operational history.  \\ \hline
    Introspective Self-Awareness & Reflecting on thoughts and emotions. & Understanding desires. & Not achievable. \\ \hline
    Social Self-Awareness        & Awareness in social contexts.        & Adapting to norms.             & Sentiment-aware chatbots.      \\ \hline
    Meta-Self-Awareness          & Awareness of awareness itself.       & Pondering existence.           & Not achievable.                \\ \hline
\end{tabular}

\subparagraph{Aspects of self-awareness}
\begin{itemize}
    \item Is self-awareness an ability or an illusion?
\end{itemize}
\begin{itemize}
    \item What is the most primitive or primary aspect of self-awareness? (link to staring at a picture)
    \begin{itemize}
        \item Is it the ability to draw a boundary between self and non-self? Is this possible at all? Is there a clear definite boundary between self and non-self?
    \end{itemize}
    \begin{itemize}
        \item Does self-awareness exist without time? For example, if a robot stares at a painting and it only has an RGB camera, can it understand that it is alive?
        \item remembering
        \begin{itemize}
            \item What is the definition of remembering?
            \item How remembering is different with reviewing?
            \item Is one aspect of self awareness is to detect more patterns or sub patterns when reviewing an experience?
        \end{itemize}
        \begin{itemize}
            \item Is remembering alone enough to say that a robot is self-aware?
            \item Can an agent be self-aware without remembering?
            \item What is the definition of remembering?
            \item Is it only passing a set of sensory data through attention?
            \item Is there a relation between language and remembering? For example do we remember better with a language?
        \end{itemize}
        \item Is attention an aspect of self-awareness? What is the definition of attention?
    \end{itemize}
\end{itemize}
\begin{itemize}
    \item Is self-awareness an ability? (what is the definition of an ability)
    \item Is self-awareness improvable?
    \item Is self-awareness a continuous feeling or it can be provoked whenever the agent wills to?
    \item Agent's knowledge being able to do things further than its programmed missions is an aspect of self-awareness?
    \item What is illusion? (Definition over sensory data and time)
    \item Causality? Its relation with language
\end{itemize}

\subparagraph{Applications of self-awareness in robotics}
\begin{itemize}
    \item Does it make a difference wether a robot is self-aware? Maybe such meta knowledge is useless
\end{itemize}


\paragraph{Self-modeling}


\paragraph{self-cognition}
% self-cognition is about the ability of a robot to think about itself, self-self_awareness is about the ability of an agent to perceive itself as an entity separate from the environment and recognize its own existence. for example distinguishing between robotic are and what it is holding or recognising itself in mirror. Self conciouness is the ability to reflect on oneself as an entity, including one’s thoughts, emotions, and existence over time.
Self-cognition in robots refers to a robot's ability to reason about its own cognitive processes, knowledge, and limitations. It involves introspection, self-monitoring, and understanding how its decisions and actions are related to its internal models or knowledge.

\begin{tabular}{|l|l|l|}
    \hline
    \textbf{Aspect}       & \textbf{Cognition}                                                     & \textbf{Self-Cognition}                                                           \\ \hline
    \textbf{Definition}   & The ability to perceive, process, learn, and act on external inputs. & The ability to reason about its own cognitive processes, states, and limitations. \\ \hline
    \textbf{Focus}        & Solving missions, making decisions, and interact with the environment. & Reflecting on its performance, knowledge, and cognitive processes. \\ \hline
    \textbf{Complexity}   & Involves perception, learning, and reasoning.                          & Involves self-monitoring, introspection, and meta-reasoning.                      \\ \hline
    \textbf{Key Features} & Decision-making, learning, and act autonomously.                    & Error detection, self-modeling, confidence estimation, and self-correction.       \\ \hline
    \textbf{Examples}     & Identifying and picking an object.                                     & Recognizing it failed to pick an object and correcting the behavior.              \\ \hline
    \textbf{Role}         & Basic autonomous operation and mission-solving.                           & Enhancing adaptability, reliability, and performance.                             \\ \hline
\end{tabular}


\paragraph{Self-consciousness}


\paragraph{Difference between the three selves}
\begin{table}[h!]
    \centering
    \begin{tabular}{|l|l|l|}
        \hline
        \textbf{Concept}
        & \textbf{Definition}
        & \textbf{Role of Self-Identification} \\ \hline
        \textbf{Self-Awareness}
        & The ability to perceive oneself as distinct from the environment.
        & Fundamental: Recognizing one’s own body and actions. \\ \hline
        \textbf{Self-Cognition}
        & The ability to think about oneself, including internal states and missions.
        & Essential: Forming an internal model of one’s capabilities. \\ \hline
        \textbf{Self-Consciousness}
        & The ability to reflect on oneself as an entity with thoughts and emotions.
        & Supporting: Higher-level reasoning about oneself over time. \\ \hline
    \end{tabular}
    \caption{Comparison of Self-Identification in Different Concepts}
\end{table}


\paragraph{Todo}
\todo{In which category self-identification should be put? self-awareness? Self-cognition or self conciouness?}

% Source: from_phd_review/robotics/sense.tex

\subsubsection{Sense}
One of the three main constituents of a robot functionality along with plan and action

\paragraph{Body}

\subparagraph{Actuator}

% Source: from_phd_review/robotics/robot/parts/body/actuator/action.tex

\ProjectDeepHeading{2}{Action}

\subparagraph{Hardware}

% Source: from_phd_review/robotics/robot/parts/body/hardware/pixhawk.tex

something similar to Arduino. Its a microcontroller board and PX4 is installed on it

% Source: from_phd_review/robotics/robot/parts/mind/mind.tex

Maybe befor writing this section, I should first check the basics of \textbf{Theory of mind}

% Source: from_phd_review/robotics/robot/parts/mind/language.tex

\ProjectDeepHeading{1}{Language}

What is alphabet? "Edmund Husserl" says that "Sound is Awareness" Thomas Nagel has explored the concept of aural awareness in his philosophical discussions in his paper "What is it like to be a bat?"

\ProjectDeepHeading{2}{History}

\ProjectDeepHeading{2}{Formation process}
    \ProjectDeepHeading{3}{My thoughts}
        
        \ProjectDeepHeading{5}{Development}
            
            \begin{figure}[H]
                \raggedright
                \begin{tikzpicture}[
                    >=Stealth,
                    every node/.style={
                        draw,
                        rectangle,
                        align=center,
                        inner sep=3pt
                        }]
                    \graph [
                        layered layout,
                        grow=down,
                        level distance=16mm,
                        edges={-Stealth},
                        ]
                        {
                            "Snesory data" -> "Abstraction \\ Clustering" -> "Inner speech"-> "Concensus with others";
                        };
                \end{tikzpicture}
            \end{figure}
            
        \ProjectDeepHeading{5}{Entities}
            
            words and alphabets and sentences and larger parts
            
        
        
    
\ProjectDeepHeading{2}{Structure}
    \ProjectDeepHeading{3}{alphabet}
    \ProjectDeepHeading{3}{words or dictionary}
        First words are formed
        
        the length of the words are related to their usage frequency. The more they are needed the less lengthy they must be
        
        We need a new words for each detected anomally
        
    \ProjectDeepHeading{3}{Composite structures}
    
    
    

\ProjectDeepHeading{2}{language relations}
    Linguistic relations (syntagmatic \& paradigmatic)
    \begin{itemize}
        \item Paradigmatic / Relation de substitution
        \item Syntagmatic / Relation de combinaison
    \end{itemize}

\ProjectDeepHeading{2}{Communication emergence in MRS}

\ProjectDeepHeading{2}{Inner speech}
    \textbf{The importance of inner speech in self-awareness} \todo{Check this paper: \url{https://www.frontiersin.org/journals/robotics-and-ai/articles/10.3389/frobt.2020.00016/full}}

\ProjectDeepHeading{2}{Language grounding with mutual presence and sense the same concept or object at the same time}

\ProjectDeepHeading{2}{Language grounding by comparing pattern in the mind}
%this is just possible for machines in human it must be described

\ProjectDeepHeading{2}{Concept composition}
    \todo{build corner, straight,  build left and associate turn left or straight left }

\ProjectDeepHeading{2}{Formation and refinement}
    \textbf{Formation}

    \textbf{Refinement} Everytime an experience is repeated the language must become more granular
    
\ProjectDeepHeading{2}{My thoughts}
first find the words by clustering. then
\begin{itemize}
    \item noun: is a cluster over the boundries of 0th derivative
        \begin{itemize}
            \item adjective: decription of the noun such as intensity
        \end{itemize}
    \item verb: is relied to actuators and commands. clustring over the first derivative. bverb must move the nouns in front of sensors of the robot
        \begin{itemize}
            \item adverb: decription of the verb such as intensity
        \end{itemize}
\end{itemize}

\ProjectDeepHeading{5}{Language formation phases}
First an internal language is developed then by a signalling system one robot tries to check if another robot has approximately found the same cluster. if it doesnt find then the word dies in the robot.

\ProjectDeepHeading{2}{Anomaly detection}
Whenever an anomally is higher than what it should be probably based on the 3 sigma rule, then its time to create new language compositions. for example a word or a sentence.


\ProjectDeepHeading{2}{Todo}
    \ProjectDeepHeading{5}{Formation idea}
        each robot  forms clusters and relatation but has only one communication channel such as a red light that can create morse code so that another robot in vicinity can see it. this blinking red light can be simulated by a channel in which a rebot emmits signals to simulate
    
    \ProjectDeepHeading{5}{Gradually}
        language and its entitity form gradually, but how could this be implemented or modeled?


\begin{itemize}
    \item what is descriprtion?
    \item I amp a weak language as map is path with between nodes between start and end? How much is it substituable?
    \item Like OOP, in natural language, the distance between concept and concrete is determinable or what ever is said is concerete or maybe concept?
    \item Did first word existed or alphabet? I think word.
    \item "time passage" is not a repeated pattern, then how can a word be corresponded to it?
    \item Should a self-aware robot talk to itself?
    \item Does talking to self needs another agent?
    \item Is abnormality detection over a continuous time interval a word or an alphabet?
    \item Can robot make sense of a word or a sentence without having experience it?
    \item When should a robot consider a set of ordered alphabets as a meaningful word?
    \item When does a robot consider a set of ordered words as a sentence?
    \item Can a robot build meaningful language structures that it knows will never occur in the real world for any agent? words such as beautiful, ugly, god
    \item Maybe emotional words can be built by growth rate or first derivative.
    \item Does the alphabet have anything to do with different levels of derivative as growth rates?
    \item Is one aspect of language to build future probable meaningful structures?
    \item What is the definition of structure?
    \item Does a human, before saying something to another, say it to himself?
    \item Is keeping accidental composition of words or alphabets in mind to continuously assess their meaningfulness an aspect?
    \item what about two opposite words or sentences?
    \item Is there a knowledge in the reality in alphabets we don't have opposite or similarity concepts?
    \item Can neuro-symbolic AI be related to self-awareness? \href{https://en.wikipedia.org/wiki/Neuro-symbolic_AI}{url}
    \item Thinking of robots thinking dependency to the programming language that builds it and machine learning methods it uses.
    \item There must be a relation between an optimal discovered language in a robot and a good programming labguage
    \item Define Optimal for an emergent language in a robot
        \begin{itemize}
            \item Is being optimal equals forming a language that uses its hierarchy and Janeshini and jaygozini methods it detects anomaly better
        \end{itemize}
\end{itemize}

% Source: from_phd_review/robotics/robot/parts/mind/state.tex

\url{https://en.wikipedia.org/wiki/Mental_state}

\subparagraph{Autopilot}

% Source: from_phd_review/robotics/robot/parts/mind/autopilot/px4.tex

PX4 is a peice of software that is used Pixhawk

\subparagraph{Cognition}

\ProjectDeepHeading{1}{Process}

% Source: from_phd_review/robotics/robot/parts/mind/cognition/process/process.tex

\url{https://en.wikipedia.org/wiki/Category:Mental_processes}

\ProjectDeepHeading{2}{Memory}

% Source: from_phd_review/robotics/robot/parts/mind/cognition/process/memory/memory_old.tex

\ProjectDeepHeading{4}{Memory}

\ProjectDeepHeading{5}{Memory architecure}
Regarding any architecture about memory the following information are important:

    \ProjectDeepHeading{6}{Smallest memory units}
        \begin{enumerate}
            \item sensory data
            \item
        \end{enumerate}
    \ProjectDeepHeading{6}{Memorizing}

    \ProjectDeepHeading{6}{Storage}

    \ProjectDeepHeading{6}{Forgetting}
        \ProjectDeepHeading{7}{Applications of Forgetting}
        Forgetting is not merely a failure of memory; it is a functional process essential for efficient cognitive performance \cite{hardt-2013-decay-happens}.
        The human brain receives vast amounts of sensory and conceptual information daily.
        Retaining all of it would overload neural storage and impair the retrieval of relevant memories .
        Forgetting serves several adaptive purposes:

        \begin{itemize}
            \item \textbf{Preventing Cognitive Overload} \\
                By eliminating non-essential information, forgetting reduces interference between similar memories, making recall of important information more accurate and efficient .
            \item \textbf{Facilitating Learning and Adaptation} \\
                Forgetting outdated or conflicting information allows the brain to update its internal models, improving decision-making in changing environments .
            \item \textbf{Enhancing Generalization} \\
                Removing specific, irrelevant details helps the brain abstract general rules from experiences, enabling better transfer of knowledge to new situations .
            \item \textbf{Emotional Regulation} \\
                Forgetting can diminish the emotional intensity of traumatic experiences over time, aiding psychological resilience and reducing chronic stress responses .
        \end{itemize}

        \ProjectDeepHeading{7}{Memory Loss Due to Natural Forgetting (Decay \& Interference)}
        Neural pathways that store a memory or skill, if not used regularly, undergo synaptic weakening.

        Additionally, new information can disrupt previous pathways (interference), making it impossible to retrieve them .

        \ProjectDeepHeading{7}{Targeted Forgetting or Active Suppression (Active Forgetting)}
        Recent research on proteins such as \textit{Rac1} in the hippocampus has shown that the brain has active biological mechanisms for erasing certain memories.

        This process is also important in learning and re-training, as it removes unnecessary pathways and frees up space for new information .



    \ProjectDeepHeading{6}{Review}
        In each review the we expect the following
        \begin{itemize}
            \item Hierrarachy (nested clustes)
        \end{itemize}
        \ProjectDeepHeading{7}{Hierrarachy}


    \ProjectDeepHeading{6}{Remembering}

    \ProjectDeepHeading{6}{Recalling}




    \ProjectDeepHeading{6}{Memory Architectures for Robots}
        \textbf{Episodic Memory} refers to the ability of robots to store specific events, including contextual information like time, place, and sensory details. This type of memory is crucial for missions requiring robots to recall past experiences and apply them to future interactions. For instance, a robot equipped with episodic memory can remember previous locations it visited or interactions with specific humans, which helps improve navigation or social behavior. The CRAM cognitive architecture integrates episodic memory for everyday mission execution in human environments, allowing robots to adapt to new scenarios by recalling past experiences.
    
        \textbf{Semantic Memory} stores general knowledge about the world, such as object properties, relationships, and actions. Unlike episodic memory, it is not tied to specific experiences but rather to facts and concepts. Robots use semantic memory to recognize objects and understand their functions. This type of memory is essential for language grounding, as robots can associate words with objects and actions. The OpenCog framework implements semantic memory by building knowledge graphs that robots can access to reason about the world.
    
        \textbf{Working Memory} functions as a short-term memory system that temporarily holds information for immediate processing. It allows robots to keep track of ongoing missions, manage missions, and adapt to dynamic environments. For example, a robot performing a multi-step mission can use working memory to remember intermediate steps and adjust its behavior based on real-time feedback. The ACT-R cognitive architecture incorporates working memory to enable robots to perform complex missions that require simultaneous handling of multiple pieces of information.
    
        \textbf{Procedural Memory} stores knowledge of how to perform missions and actions, enabling robots to execute learned skills without conscious recall. This memory type is critical for repetitive missions, such as industrial robots performing assembly line missions or humanoid robots learning to walk. Procedural memory is built through reinforcement learning or imitation learning, where the robot refines its motor skills through trial and error. The ROS-TM system, for instance, implements procedural memory to automate robot missions based on previously learned actions.
    
        \textbf{Associative Memory} allows robots to link different pieces of information through learned associations. This type of memory is crucial for pattern recognition, decision-making, and adaptive behavior. Associative memory helps robots recognize objects or situations based on partial information by recalling previous associations between sensory inputs and outcomes. The Neural Turing Machine is a notable implementation of associative memory that enables robots to learn and recall complex sequences and relationships between data points.
    
        \textbf{Long-Term Memory} refers to the permanent storage of knowledge and experiences. Robots use long-term memory to recall information over extended periods, which is essential for personalized interactions and knowledge retention. For example, a customer service robot with long-term memory can remember user preferences and interactions to provide more tailored responses. The Soar cognitive architecture incorporates long-term memory to store symbolic knowledge and use it for problem-solving and reasoning.
    
        \textbf{Hierarchical Temporal Memory (HTM)} is a memory model inspired by the structure of the neocortex. HTM focuses on learning spatial and temporal patterns in sensory data, allowing robots to predict future events based on past experiences. It is particularly useful for missions that require the robot to process sequences, such as speech recognition, predictive maintenance, and motion prediction. The Numenta HTM framework has been applied in various robotic applications to enhance predictive capabilities.



\ProjectDeepHeading{5}{Remembering}
    \ProjectDeepHeading{8}{Remebering}
        Remembering refers to the general process of retrieving stored information, whether it is triggered actively or occurs spontaneously. It encompasses a broad range of memory activities, including implicit and explicit retrieval mechanisms. 
        
    \ProjectDeepHeading{8}{Recalling}
     Recalling highlights the deliberate and mission-oriented act of retrieving particular information from memory. For instance, a robot remembering involves the spontaneous recognition of a previously encountered environment, while recalling might involve actively retrieving a stored navigation strategy to solve a specific mission. These distinctions are essential when designing cognitive architectures or memory systems in robotics to differentiate between general memory access and targeted retrieval operations .

    Remembering can occur through two primary mechanisms: (1) from current sensory observations and (2) through internally triggered processes. Both mechanisms are essential for forming, retrieving, and utilizing memories in both natural and artificial systems.

\ProjectDeepHeading{5}{Models for Remembering}

    \ProjectDeepHeading{6}{Remembering from Current Sensory Observations}
    This category includes mechanisms that rely on sensory inputs to trigger memory recall or to update stored information. These models are fundamental for learning and adapting to new environments.

        \ProjectDeepHeading{7}{Models Requiring Sensory Sequences}
            \textbf{Sequence Learning Models:} Sequence learning encodes temporal patterns in sensory data, enabling the system to predict upcoming events. For instance, in language processing, sequence learning helps anticipate the next word in a sentence, while robots use this model to predict subsequent actions in mission execution.

            \textbf{Hierarchical Temporal Memory (HTM):} HTM models store hierarchical representations of sensory sequences, recognizing increasingly abstract patterns over time. These models are effective for understanding spatial and temporal regularities in dynamic environments, such as a robot navigating a recurring pattern of obstacles. 

        \ProjectDeepHeading{7}{Models Requiring Single Sensory Observations}
            \textbf{Predictive Coding Models:} Predictive coding uses sensory inputs to generate and refine predictions. By minimizing the \textit{prediction error} between expected and observed inputs, the brain adapts its internal model for future use. This process allows the formation of increasingly accurate memory-based predictions, such as anticipating a ball’s motion after observing its trajectory multiple times. 

            \textbf{Bayesian Frameworks:} Bayesian models depend on sensory data as evidence to update prior knowledge into posterior distributions. This allows systems to generalize knowledge and maintain dynamic probabilistic representations of the environment. For example, Bayesian reasoning enables robots to predict object trajectories from noisy sensory observations. 

            \textbf{Hebbian Learning and Associative Memory:} Associative memory is strengthened through simultaneous sensory observations, linking inputs into robust patterns. For instance, learning to associate a dog’s bark with its visual appearance demonstrates how sensory inputs trigger associative recall. 

    \ProjectDeepHeading{6}{Internally Triggered Remembering}
    These models enable memory retrieval and utilization without the need for immediate sensory input. Instead, internal processes such as thoughts, queries, or intrinsic motivations drive the recall of stored information.

        \textbf{Associative Memory Models:} Memories stored as networks of associations can be recalled using internally generated cues. For instance, recalling a mission plan may trigger retrieval of related steps even without external sensory stimuli. This mechanism supports reasoning and problem-solving in both biological and artificial systems.

        \textbf{Cognitive Architectures:} Architectures like ACT-R and Soar retrieve information based on internal missions or queries. These frameworks allow systems to adaptively access stored plans or knowledge based on their current objectives. For example, a robot can recall a stored navigation strategy when its internal system determines a similar context.

        \textbf{Generative Models:} Generative models, such as VAEs and GANs, simulate scenarios by recalling or creating representations from memory. These models internally generate patterns, enabling predictive plan and decision-making. For instance, robots may simulate mission outcomes based on past data to guide actions.

        \textbf{Episodic Memory Models:} Episodic memory allows systems to retrieve specific events or experiences using internal queries. Memories tagged with contextual details, such as temporal or spatial markers, enable targeted recall. For instance, a robot can recall the last occurrence of an obstacle to decide on its current navigation. 

        \textbf{Intrinsic Motivation Systems:} Intrinsic motivations like curiosity or novelty detection provoke memory recall to resolve uncertainties or achieve learning missions. For example, a robot might recall interactions from similar situations to address inconsistencies in its understanding. These systems enable self-driven learning in dynamic environments.

        \textbf{Metacognition and Self-Reflection:} Metacognitive systems retrieve memories based on self-evaluation, enabling the adjustment of strategies or correction of errors. For example, a robot analyzing a failed mission might recall relevant past experiences to improve its approach. This self-reflective capability fosters long-term autonomy and adaptability.


\ProjectDeepHeading{5}{Collective memory}



\ProjectDeepHeading{5}{Questions}
Mathematical models of remembering

remembering by models? Words or just raw sensory data?

Can we call an agent that can just remember or review sensory data self-aware?

What is reviewing?

Can we call a a robot that can only remember/invoke/provoke sensory data from just a given time interval aware or intelligent?

\chapter{Onthology}

% Source: onthology/onthology.tex

\cite{goncalves-etal-2023-ieee-autonomous-robotics-ontology}

% Source: robot/robot.tex

\subsubsection{Autonomous agent}
        An autonomous agent is a system situated within and a part of an environment that senses that environment and acts on it, over time, in pursuit of its own agenda and so as to effect what it senses in the future
        \cite{franklin-1997-is-it-an-agent-or-just-a-program-a-taxonomy-for-autonomous-agents}

% Source: robot/part/mind/mind.tex

\url{https://en.wikipedia.org/wiki/Mind}
    Mind body problem

% Source: robot/telemetry.tex

Telemetry is the in situ collection of measurements or other data at remote points and their automatic transmission to receiving equipment (telecommunication) for monitoring.
    \seehere{https://en.wikipedia.org/wiki/Telemetry}
    \semtag{mavlink}

\section{Anatomy}

% Source: robot/anatomy/anatomy.tex

\begin{forest}
for tree={
    edge={<-},
    align=center
}
[Robot
    [Mind
        [Cognition
            [Object\\level
                [Processes
                    [Memory
                        [Explicit
                            [Long-term
                                [Episodic
                                    [Autobiographical
                                        [Event-specific\\knowledge]
                                        [Lifetime]
                                        [General\\events]
                                    ]
                                    [Context
                                        [Remembering
                                            [Recollection]
                                        ]
                                    ]
                                ]
                                [Semantic]
                                [Item
                                    [knowing
                                        [Familiarity]
                                    ]
                                ]
                            ]
                        ]
                        [Implicit
                            [Skills\\(Procedural)]
                            [Repetition\\priming]
                        ]
                    ]
                    [Attention]
                    [Perception]
                    [Thinking]
                    [Learning]
                    [Language]
                ]
            ]
            [Meta cognition
                [Self-consciousness
                    [Self-awareness
                        [Body\\awareness]
                    ]
                ]
                [Self-monitoring]
                [Self-evaluation]
                [Self-regulation]
            ]
        ]
        [Psyche
            % This is the whole realm of mind. That is concious and unconcious
            [Consciousness]% some cognitive processes are concious some are unconcious
            [Unconsciousness]
        ]
    ]
    [Body
        [Nervous\\system
            [Nerve]
            [Neuron,
                l sep=18mm
                [Sensory, edge label={node[midway, fill=white]{is-a}}]
                [Motor, edge label={node[midway, fill=white]{is-a}}]
                [Interneuron, edge label={node[midway, fill=white]{is-a}}]
                [Projection, edge label={node[midway, fill=white]{is-a}}]
            ]
            [Central
                [Brain\\(CPU/FCU)]
                [Retina\\(GPU)]
                [Spinal cord\\(BUS)]
            ]
            [Peripheral
                [Autonomic]
                [Enteric]
                [Somatic]
            ]
        ]
        [Limb\\(Actuator)]
        [Organ\\(Link)]
    ]
]
\end{forest}

\subsection{Humanoid}

\subsubsection{Booster}

\paragraph{K1}

% Source: robot/kind/humanoid/booster/k1/k1.tex

Booster Robotics K1


    The K1 is also a humanoid research robot, typically used for:
    
    
    \begin{itemize}
    
        \item
        Dynamic walking and balance control
        
        
        
        
        \item
        Whole-body motion planning
        \item
        Real-time control architectures
        \item
        Sensor fusion (IMU, cameras, possibly LiDAR)
        \item
        SDK-level integration for autonomy stacks
    
    
    \end{itemize}

\subsubsection{Utree}

\paragraph{Kind}

% Source: robot/kind/humanoid/utree/kind/g1_h1.tex

Unitree G1 / H1

    These are humanoid robots developed by Unitree Robotics.
    
    Technologies involved:
    
    Biped locomotion control (dynamic balance, whole-body control)
    
    High-DOF joint actuation systems
    
    Real-time motor control and torque control
    
    IMU-based state estimation
    
    Visual and LiDAR perception systems
    
    ROS/SDK integration for autonomy and navigation
    
    Onboard compute (usually x86 or ARM platforms with GPU support)
    
    The H1 is a full-sized humanoid, while G1 is a smaller research-oriented humanoid platform.

\subsection{Uav}

\subsubsection{Kind}

\paragraph{Quadcopter}

\subparagraph{Building toturial}

% Source: robot/kind/uav/kind/quadcopter/building_toturial/building_tutorial.tex

\url{https://www.instructables.com/How-to-Make-a-Drone-With-Pixhawk-4-and-Spektrum/}
    \\
    \url{https://www.youtube.com/watch?v=wsrRNqihjE4&list=PLp8NnRiCCsXs-PI_jt96-bB4bOuiXRLKg&index=1}
    
    % 
    %

\subparagraph{Kind}

\ProjectDeepHeading{1}{Crazyfly}

% Source: robot/kind/uav/kind/quadcopter/kind/crazyfly/crazyfly.tex

\ProjectDeepHeading{4}{Crazyflie 2.1+ Overview}
        
        \textbf{Platform:} Crazyflie 2.1+ is a palm-sized open-source micro quadcopter designed primarily for research and development rather than consumer applications.
        
        \ProjectDeepHeading{5}{Physical Characteristics}
            
            \begin{itemize}
                \item Weight: approximately 29 grams
                \item Motor-to-motor distance: approximately 92 mm
                \item Fully palm-sized form factor
            \end{itemize}
        
        \ProjectDeepHeading{5}{Onboard Sensors (IMU)}
            
            The main board includes:
            
            \begin{itemize}
                \item 3-axis accelerometer and 3-axis gyroscope (BMI088)
                \item 3-axis magnetometer
                \item Barometer (BMP388)
            \end{itemize}
            
            These sensors enable state estimation, attitude control, and onboard filtering (e.g., EKF).
        
        \ProjectDeepHeading{5}{Processing Units}
            
            \begin{itemize}
                \item STM32F405 (Cortex-M4) for flight control
                \item nRF51822 for radio communication
            \end{itemize}
        
        \ProjectDeepHeading{5}{Communication Interfaces}
            
            \begin{itemize}
                \item 2.4 GHz radio communication (Crazyradio)
                \item USB interface
                \item Optional WiFi via AI Deck
            \end{itemize}
        
        \ProjectDeepHeading{5}{Expansion Capability (Deck System)}
            
            Crazyflie uses a modular expansion deck system, allowing additional hardware to be mounted directly on top of the platform.
            
            \ProjectDeepHeading{6}{Multi-Ranger Deck}
            
            \begin{itemize}
                \item Five Time-of-Flight (ToF) laser sensors
                \item Measures distance in front, back, left, right, and upward directions
                \item Range up to approximately 4 meters
                \item Functions as a lightweight multi-directional ranging system
                \end{itemize}
            
            \ProjectDeepHeading{6}{Flow Deck v2}
            
            \begin{itemize}
                \item Optical flow sensor
                \item Downward-facing ToF sensor for altitude measurement
                \item Enables stable indoor flight without GPS
                \end{itemize}
            
            \ProjectDeepHeading{6}{AI Deck 1.1}
            
            \begin{itemize}
                \item GAP8 RISC-V processor
                \item 320×320 monochrome camera (Himax HM01B0)
                \item ESP32 for WiFi connectivity
                \item Supports onboard AI workloads and custom vision processing
                \end{itemize}
        
        \ProjectDeepHeading{5}{Software Characteristics}
            
            \begin{itemize}
                \item Fully open-source firmware
                \item Python API available
                \item Firmware-level modification possible
                \item Compatible with ROS
                \item Simulation support (e.g., Gazebo)
            \end{itemize}
        
        \ProjectDeepHeading{5}{Advantages}
            
            \begin{itemize}
                \item Extremely lightweight and compact
                \item Fully programmable at firmware level
                \item Modular sensor expansion
                \item Suitable for swarm research
                \item Suitable for indoor autonomy experiments
            \end{itemize}
        
        \ProjectDeepHeading{5}{Limitations}
            
            \begin{itemize}
                \item Very limited payload capacity
                \item Short flight time (approximately 7–10 minutes)
                \item Poor wind resistance
                \item Ranging sensors are ToF-based rather than scanning LiDAR
            \end{itemize}
        
        \ProjectDeepHeading{5}{Conclusion}
            
            For a minimal fully-programmable quadcopter platform with onboard IMU, expandable ranging sensors, and optional camera-based AI processing, Crazyflie 2.1+ with Multi-Ranger and AI Deck represents one of the smallest practical research platforms currently available.

\section{Life}

% Source: robot/life/life.tex

\cite{friston-2013-life-as-we-know-it} suggests that life—or biological self-organization—is an inevitable and emergent property of any (ergodic) random dynamical system that possesses a Markov blanket.

\subsection{Allostasis}

% Source: robot/life/allostasis/allostasis.tex

Allostasis is a physiological mechanism of regulation in which an organism anticipates and adjusts its energy use according to environmental demands. \url{https://en.wikipedia.org/wiki/Allostasis}
    %
    %

\subsection{Autopoiesis}

% Source: robot/life/autopoiesis/autopoiesis.tex

Autopoiesis one of several current theories of life, refers to a system capable of producing and maintaining itself by creating its own parts
    \seeherefooter{https://en.wikipedia.org/wiki/Autopoiesis}
    
    \cite{maturana-varela-1980-autopoiesis-and-cognition-the-realization-of-the-living}
    defines for the first time a living being as a system that continuously produces and reproduces itself. Its parts generate the very network that, in turn, produces those same parts.
    So, being alive means having a self-producing and self-maintaining organization, not merely possessing a certain material or behavior.

\subsection{Homeostasis}

% Source: robot/life/homeostasis/homeostasis.tex

Homeostasis is the state of steady internal physical and chemical conditions maintained by living organisms
    \seeherefooter{https://en.wikipedia.org/wiki/Homeostasis}

\subsection{Self organisiation}

% Source: robot/life/self_organisiation/language_related.tex

\cite{witzany-2014-biological-self-organization} argues that “biological self-organization” should not be explained merely through mechanical cause-and-effect models. The author argues that in living beings, organization does not occur only through the physical exchange of matter, but rather depends on communicative processes; that is, cells, tissues, organs, and organisms interact with one another in a “meaningful” way through signals, and without these interactions biological coordination is not possible.
    
    The core of the paper is that biological communication does not consist only of a sender and a receiver, but, like natural language, has three layers: syntax, meaning, and context-dependent use. In other words, a biological signal does not produce only one fixed and mechanical reaction; the same signal can take on different meanings depending on the context, time, place, and condition of the living being. According to the author, it is precisely this dependence on context that separates life from non-living processes.
    
    Then the paper moves into genetics and says that DNA should not be regarded merely as a passive storehouse of information. The author emphasizes the role of non-coding RNAs, mobile genetic elements, and viruses, and says that these have active roles in the regulation, alteration, and rearrangement of the genome. In this view, biological novelty is not only the result of “random mutation,” but is to some extent the result of the action of agents that can manipulate and re-regulate the genetic text.
    
    Another important part of the paper is its critique of mathematical and formal views of language. The author says that models which see language or the genetic code only as a formal syntax or a computable structure cannot explain the main feature of natural language: namely, that meaning is formed within social context and real use. He extends this point to biology and says that biological codes, too, are not understood only through formal rules, but depend on living agents and the practical context in which signs are used.
    
    The overall conclusion of the paper is that biological self-organization is fundamentally a “communicative-semiotic” process, not merely a physical-mechanical one. That is, life should be understood as a network of living agents that coordinate themselves through signs, rules, and context-dependent interpretation. For this reason, the author defends a non-reductionist and non-mechanistic description of life.
    
    If I want to put it very simply, the paper’s point is this:
    A living being does not merely “react”; it “interprets signs,” and this interpretation is the basis of the organization of life.

% Source: robot/life/self_organisiation/self_organisation.tex

\cite{ashby-1947-principles-of-the-self-organizing-dynamic-system} presents one of the classic formulations of self-organization and argues that some dynamical systems can move toward order and stability without an external designer.
    His focus is mainly cybernetic and systemic: how a system, through its own internal dynamics, can settle into more organized states.
    This work lays a theoretical foundation for self-organization, rather than explaining life directly in the language of modern biology.
    
    Biological organization is the organization of complex biological structures and systems that define life using a reductionistic approach \seehere{https://en.wikipedia.org/wiki/Biological_organisation}.
    
    
    
    \cite{kauffman-1993-the-origins-of-order-self-organization-and-selection-in-evolution} argues that biological order does not arise from natural selection alone; self-organization also plays a fundamental role.
    He shows that complex networks can spontaneously generate patterns, structures, and stable constraints, and that these provide the basis on which evolution operates.
    In this view, life emerges from the interplay of selection and self-organization, not from only one of them.

\section{Mind body transition}

% Source: robot/mind_body_transition/mind_body_problem.tex

Cehvck for dualism
    
    The mind–body problem is a philosophical problem concerning the relationship between thought and consciousness in the human mind and body. It addresses the nature of consciousness, mental states, and their relation to the physical brain and nervous system.
    \url{https://en.wikipedia.org/wiki/Mind%E2%80%93body_problem}
    
    \cite{georgiev-2020-quantum-information-theoretic-approach-to-the-mind-brain-problem}  argues that consciousness can be understood as unobservable quantum information integrated in quantum brain states. It distinguishes the private, unobservable mind from the observable brain by using the quantum distinction between states and observables. It also uses Holevo’s theorem to explain why conscious experience is only partially communicable through classical information.

% Source: robot/mind_body_transition/property_dualism.tex

Property dualism describes a category of positions in the philosophy of mind which hold that, although the world is composed of just one kind of substance—the physical kind—there exist two distinct kinds of properties: physical properties and mental properties.
    \url{https://en.wikipedia.org/wiki/Property_dualism}
    % 
    %

\subsection{Body}

\subsubsection{Action}

% Source: robot/part/body/action/action.tex

In philosophy, an action is something an agent does. Actions contrast with events which merely happen to someone and are typically performed for a purpose and guided by an intention. \seehere{https://en.wikipedia.org/wiki/Action_(philosophy)}
    %
    %

\subsubsection{Part}

% Source: robot/part/body/part/part.tex

\begin{figure}[H]
        \centering
        \ProjectSafeIncludeGraphics[width=0.9\textwidth]{robot/part/body/body_parts_hierarchy.png}
    \end{figure}
    %
    %

\paragraph{Limb}

% Source: robot/part/body/part/limb/limb.tex

A limb is a body appendage used mainly for movement (arm or leg)
    
    A limb contains multiple organs and tissues, but it is classified as an anatomical region rather than a single organ.
    
    A limb is a body appendage used mainly for movement (arm or leg).
    A limb contains multiple organs and tissues, but it is classified as an anatomical region rather than a single organ.

\subparagraph{Actuator}

% Source: robot/part/body/part/limb/actuator/actuator.tex

Actuator is whatever that converts a control signal to physical force

\ProjectDeepHeading{1}{Rotor}

\ProjectDeepHeading{2}{Brusless}

% Source: robot/part/body/part/limb/actuator/rotor/brusless/brushless.tex

Uses PWM to increase and decrease volatage frequency.

\paragraph{Nervous system}

\subparagraph{Neurorobotics}

% Source: robot/part/body/part/nervous_system/neurorobotics/neurorobotics.tex

Neurorobotics is the combined study of neuroscience, robotics, and artificial intelligence. It is the science and technology of embodied autonomous neural systems \seeherefooter{https://en.wikipedia.org/wiki/Neurorobotics}.

\subparagraph{Neuroscience}

% Source: robot/part/body/part/nervous_system/neuroscience/neuroscience.tex

\ProjectDeepHeading{2}{Neuroscience}
    
\ProjectSafeIncludePdf[pages=-]{../assets/review/biology/neuroscience.pdf}

\subparagraph{Part}

\ProjectDeepHeading{1}{Central}

\ProjectDeepHeading{2}{Brain}

\ProjectDeepHeading{3}{Part}

\ProjectDeepHeading{4}{Cerebellum}

% Source: robot/part/body/part/nervous_system/part/central/brain/part/cerebellum/cerebellum.tex

The cerebellum, a prominent brain structure at the back of the brain below
    the neocortex, is involved in fine-tuning motor sequences and motor learning.
    Within the cerebellum are repeated neural circuits forming modules that deal
    with motor planning, motor execution, and balance.

\ProjectDeepHeading{3}{Region}

% Source: robot/part/body/part/nervous_system/part/central/brain/region/region.tex

For a list of brain regions \seeherefooter{https://en.wikipedia.org/wiki/List_of_regions_in_the_human_brain}

\ProjectDeepHeading{4}{Network}

\ProjectDeepHeading{5}{Large scale}

% Source: robot/part/body/part/nervous_system/part/central/brain/region/network/large_scale/large_scale.tex

Large-scale brain networks (also known as intrinsic brain networks) are collections of widespread brain regions showing functional connectivity by statistical analysis of the fMRI BOLD signal or other recording methods such as EEG PET and MEG \seeherefooter{https://en.wikipedia.org/wiki/Large-scale_brain_network}
    
    \cite{bressler-2010-large-scale-brain-networks-in-cognition-emerging-methods-and-principles} argue that cognition cannot be explained well by isolated brain modules, but by dynamic interactions among distributed large-scale brain networks. They distinguish structural networks, based on anatomical connections, from functional networks, based on coordinated activity among brain regions. The paper reviews methods such as diffusion imaging, fMRI, EEG/MEG, Granger causality, dynamic causal modeling, and graph theory for studying these networks. It highlights major networks such as the default-mode network, central-executive network, and salience network, and explains their roles in self-related cognition, control, attention, and switching. The main conclusion is that large-scale brain network analysis offers a systematic framework for understanding cognition, learning, and cognitive disorders.

\ProjectDeepHeading{6}{Mode}

% Source: robot/part/body/part/nervous_system/part/central/brain/region/network/large_scale/mode/mode.tex

Large-scale functional brain networks are distributed sets of brain regions that show coordinated activity across tasks or during rest. In the resting-state fMRI literature, these networks are also called intrinsic connectivity networks. A widely used parcellation is the seven-network organization proposed by \citet{yeo-2011-the-organization-of-the-human-cerebral-cortex}, which includes the visual network, somatomotor network, dorsal attention network, ventral attention network, limbic network, frontoparietal control network, and default mode network.
    
    
    \ProjectDeepHeading{9}{Visual Network}
        
        The visual network is mainly associated with the processing of visual sensory information, such as shape, color, motion, and spatial structure.
    
    
    \ProjectDeepHeading{9}{Somatomotor Network}
        
        The somatomotor network is associated with bodily sensation and motor control. It includes cortical regions involved in movement execution and somatosensory processing.
    
    
    \ProjectDeepHeading{9}{Dorsal Attention Network}
        
        The dorsal attention network is associated with goal-directed external attention. It is active when attention is deliberately oriented toward external objects, spatial locations, or task-relevant sensory information.
    
    
    \ProjectDeepHeading{9}{Ventral Attention Network}
        
        The ventral attention network is associated with stimulus-driven attention. It is involved when unexpected or behaviorally relevant stimuli interrupt the current focus of attention.
    
    
    \ProjectDeepHeading{9}{Limbic Network}
        
        The limbic network is associated with emotion, valuation, motivation, and affective aspects of cognition.
    
    
    \ProjectDeepHeading{9}{Frontoparietal Control Network}
        
        The frontoparietal control network is associated with cognitive control, task demands, working memory, decision-making, and flexible regulation of behavior. In some literature, it overlaps conceptually with the central executive network.
    
    
    \ProjectDeepHeading{9}{Default Mode Network}
        
        The default mode network is associated with internally oriented cognition, including autobiographical memory, future thinking, self-referential processing, social cognition, and mind wandering. It was originally described as a baseline mode of brain function that is reduced during many externally demanding tasks \citep{raichle-2001-a-default-mode-of-brain-function}.
    
    
    \ProjectDeepHeading{9}{Salience Network}
        
        The salience network is not always separated in the seven-network Yeo parcellation, where it is often close to the ventral attention network. However, in the triple-network model, it is treated as a central network for detecting behaviorally relevant events and switching between the default mode network and the central executive network \citep{menon-2011-large-scale-brain-networks-and-psychopathology,menon-2010-saliency-switching-attention-and-control}.
    
    
    \ProjectDeepHeading{9}{Triple-Network Model}
        
        The triple-network model focuses on the interaction between the default mode network, the salience network, and the central executive network. In this view, the salience network helps regulate switching between internally oriented cognition and externally oriented cognitive control \citep{menon-2011-large-scale-brain-networks-and-psychopathology}.

\ProjectDeepHeading{7}{Kind}

% Source: robot/part/body/part/nervous_system/part/central/brain/region/network/large_scale/mode/kind/default_mode_network.tex

Being active when a person is not focused on the outside world and the brain is at wakeful rest, such as during daydreaming and mind-wandering \seeherefooter{https://en.wikipedia.org/wiki/Default_mode_network}
    
    \cite{raichle-2015-the-brains-default-mode-network} reviews the discovery and development of the concept of the brain’s default mode network. The paper explains that the DMN was first noticed through task-induced decreases in activity during attention-demanding tasks.
    It describes the DMN as a set of bilateral cortical regions involving medial prefrontal, posterior cingulate/precuneus, parietal, and temporal areas. The review connects the DMN to resting-state functional connectivity, intrinsic brain activity, memory, self-related processing, and emotional regulation. Raichle argues that the DMN is not merely a mind-wandering network, but a central part of the brain’s intrinsic organization in health and disease.
    
    
    \ProjectDeepHeading{10}{Relation with reviewing memories}
        \cite{andrews-hanna-2010-evidence-for-the-default-networks-role-in-spontaneous-cognition}
        examine whether the default mode network supports spontaneous cognition. They show that activity in default-network regions increases when people report internally directed thoughts. These thoughts often include autobiographical memories, future simulations, and self-related mental content. The paper argues that the default network is not simply “resting,” but actively supports mental processes detached from the external task. Its main contribution is linking DMN activity to spontaneous thought, especially personal past and future thinking.
    
    
    \ProjectDeepHeading{10}{Relation with internal narrative or internal language}
        \cite{menon-2023-20-years-of-the-default-mode-network-a-review-and-synthesis} reviews two decades of research on the Default Mode Network and explains how its role has expanded from a “resting-state” network to a central system for self-reference, social cognition, episodic memory, semantic memory, language, and mind wandering. The paper argues that the Default Mode Network is not a single-purpose system, but a set of interacting nodes and subnetworks whose functions depend on dynamic interactions with salience and frontoparietal control networks. Its main synthesis is that the Default Mode Network integrates memory, language, and semantic representations to construct a coherent internal narrative, which supports the sense of self and subjective continuity.

% Source: robot/part/body/part/nervous_system/part/central/brain/region/network/large_scale/mode/kind/dorsal_attention_network.tex

\seeherefooter{}
    \cite{corbetta-2002-control-of-goal-directed-and-stimulus-driven-attention-in-the-brain} distinguishes two interacting attention systems in the human brain: a dorsal frontoparietal system for goal-directed/top-down attention and a ventral system for stimulus-driven reorienting. The dorsal system, including frontal eye fields and intraparietal/superior parietal regions, is active when attention is voluntarily directed according to current goals. The paper is useful because it gives the classical framework for separating intentional attention from attention captured by unexpected or salient stimuli.
    
    \cite{vossel-2014-dorsal-and-ventral-attention-systems} updates the dorsal/ventral attention distinction and emphasizes that the two systems are anatomically distinct but functionally cooperative. The dorsal attention system supports controlled spatial attention and task-based selection, while the ventral system helps detect behaviorally relevant events and redirect attention. The paper is useful because it presents a more modern view: attention is not simply split into two isolated networks, but depends on interaction between them.
    
    \cite{corbetta-2008-the-reorienting-system-of-the-human-brain-from-environment-to-theory-of-mind} explains the human brain’s reorienting system, especially the right-lateralized ventral frontoparietal network, which helps redirect attention toward behaviorally relevant unexpected stimuli. It distinguishes this ventral network from the dorsal attention network: the dorsal network maintains goal-directed attention, while the ventral network interrupts or resets ongoing activity when important new events appear. The paper also extends the role of the ventral attention network beyond simple environmental attention, linking it to network switching, internally directed thought, and theory-of-mind cognition.

% Source: robot/part/body/part/nervous_system/part/central/brain/region/network/large_scale/mode/kind/mind_wandering.tex

\seeherefooter{https://en.wikipedia.org/wiki/Default_mode_network}

    \cite{mason-2007-wandering-minds-the-default-network-and-stimulus-independent-thought} examine the relation between the default network and stimulus-independent thought, especially mind wandering. The paper argues that when people are not fully occupied by external tasks, internally generated thoughts become more likely. Using fMRI evidence, the authors link activity in the default network to these wandering, self-generated mental states. The paper is important because it connects the default mode network not only to rest, but also to spontaneous cognition and task-unrelated thought.

% Source: robot/part/body/part/nervous_system/part/central/brain/region/network/large_scale/mode/kind/ventral_attention_system.tex

\cite{scalf-2014-trial-history-effects-in-the-ventral-attentional-network}  investigates whether the ventral attentional network preserves information from previous trials and is affected when previously ignored stimuli become task-relevant. Using fMRI and the distractor preview effect, the authors found stronger activation in VAN-related regions, especially the right middle frontal gyrus and right supramarginal gyrus, when current targets had been rejected in the previous trial. The study argues that VAN is not only a stimulus-driven reorienting system, but also carries trial-history information that shapes current attentional selection.

\ProjectDeepHeading{3}{Theory}

\ProjectDeepHeading{4}{Epistemic foraging}

% Source: robot/part/body/part/nervous_system/part/central/brain/theory/epistemic_foraging/epistemic_foraging.tex

\url{https://en.wikipedia.org/wiki/Foraging}
    
    Optimal foraging theory (OFT) is a behavioral ecology model that helps predict how an animal behaves when searching for food.
    \url{https://en.wikipedia.org/wiki/Optimal_foraging_theory}
    
    \url{https://www.envisioning.com/vocab/epistemic-foraging}

\ProjectDeepHeading{4}{Kind}

% Source: robot/part/body/part/nervous_system/part/central/brain/theory/kind/kinds.tex

%
    %

\ProjectDeepHeading{5}{Bayesian}

% Source: robot/part/body/part/nervous_system/part/central/brain/theory/kind/bayesian/bayesian.tex

Bayesian approaches to brain function investigate the capacity of the nervous system to operate in situations of uncertainty in a fashion that is close to the optimal prescribed by Bayesian statistics. \url{https://en.wikipedia.org/wiki/Bayesian_approaches_to_brain_function}
    %
    %

\ProjectDeepHeading{6}{Active inference}

% Source: robot/part/body/part/nervous_system/part/central/brain/theory/kind/bayesian/active_inference/active_inference.tex

\ProjectDeepHeading{12}{Active inference}
    \url{https://en.wikipedia.org/wiki/Free_energy_principle\#Active_inference}
    \url{https://en.wikipedia.org/wiki/Free_energy_principle}
    \ProjectSafeIncludePdf[pages=-, fitpaper=true]{../assets/proposed/robot/structure/mind/active_inference.pdf}
    Active inference can explain cognitive proccesses such as perception, action selection and

% Source: robot/part/body/part/nervous_system/part/central/brain/theory/kind/bayesian/active_inference/free_energy.tex

% 
    %

\ProjectDeepHeading{7}{Suprise}

% Source: robot/part/body/part/nervous_system/part/central/brain/theory/kind/bayesian/active_inference/suprise/suprise.tex

\cite{weissglass-2019-greatest-surprise-reduction-semantics-an-information-theoretic-solution-to-misrepresentation-and-disjunction} This paper proposes a theory of signal content according to which the content of a signal is the event whose surprisal is reduced the most by that signal. The author uses this idea to address two major problems in psychosemantics: **misrepresentation** and the **disjunction problem.

% Source: robot/part/body/part/nervous_system/part/central/brain/theory/kind/bayesian/active_inference/suprise/suprise_reduction_semantics.tex

\

\ProjectDeepHeading{8}{Ai}

% Source: robot/part/body/part/nervous_system/part/central/brain/theory/kind/bayesian/active_inference/suprise/ai/ai.tex

\ProjectDeepHeading{11}{Executive Summary}
        In active inference, surprise is not a casual psychological notion; it is the negative log probability of an outcome under the agent's generative model. An observation becomes surprising when the model assigns it low probability. Because exact surprise is usually intractable, agents minimize variational free energy, which serves as a computable upper bound on surprise.
        
        Surprise can increase when sensory input deviates from top-down predictions, when priors are too rigid or poorly calibrated, when precision weighting amplifies prediction errors, or when the model is misspecified relative to the environment. Friston's framework explains surprise reduction mainly through two coupled routes: changing beliefs and predictions through perception or learning, and changing sensory input through action.
    
    
    \ProjectDeepHeading{11}{Background: Surprise and Free Energy in Active Inference}
        In Friston's formulation, surprise is defined as
        \begin{equation}
            I(o) = -\ln p(o),
        \end{equation}
        where $o$ is an observation and $p(o)$ is the probability assigned to it by the generative model. If the model expects an outcome, surprise is low; if the outcome is improbable, surprise is high.
        
        Because the exact marginal probability $p(o)$ is often difficult to compute, active inference works with variational free energy,
        \begin{equation}
            F = D_{KL}\bigl(q(s) \parallel p(s \mid o)\bigr) - \ln p(o),
        \end{equation}
        which can also be written in accuracy-complexity form. By minimizing $F$, the agent indirectly minimizes surprise. This is the basis for perception, action, and learning in the free-energy principle.
    
    
    \ProjectDeepHeading{11}{Factors Increasing Surprise}
        Several factors increase surprise under active inference:
        \begin{itemize}[leftmargin=1.8em]
	\item \textbf{Prediction mismatch:} sensory input differs from the predicted input generated by higher-level beliefs.
    \item \textbf{Rigid or poor priors:} the model strongly expects one class of states and cannot accommodate alternatives.
    \item \textbf{Incorrect precision weighting:} prediction errors are overweighted or underweighted.
    \item \textbf{Model misspecification:} the generative model lacks the latent causes needed to explain the observation.
    \item \textbf{Novel structure in the environment:} the world has changed in a way that learned regularities no longer capture.
        \end{itemize}
        
        \begin{table}[H]
            \centering
            \small
            \begin{tabularx}{\textwidth}{>{\raggedright\arraybackslash}p{3.2cm} >{\raggedright\arraybackslash}X >{\raggedright\arraybackslash}X}
		\toprule
        \rowcolor{gray!15}
        \textbf{Factor} & \textbf{Mathematical role} & \textbf{Effect on surprise} \\
        \midrule
        Prediction mismatch & Large sensory prediction error under the likelihood model & Raises free energy and signals poor explanatory fit \\
        Rigid prior beliefs & Strong prior concentration on a narrow state space & Makes deviations appear highly improbable \\
        Misweighted precision & Changes the gain on prediction errors & Can exaggerate or suppress the impact of mismatch \\
        Model misspecification & Missing hidden causes or wrong dependencies & Prevents accurate inference even with updating \\
        Environmental novelty & True change in external structure or contingencies & Produces sustained surprise until learning occurs \\
        \bottomrule
            \end{tabularx}
            \caption{Main factors that increase surprise in active inference.}
        \end{table}
    
    
    \ProjectDeepHeading{11}{Likelihood, Prior, Posterior, and Precision}
        Surprise depends on the relation between four core elements.
        
        \ProjectDeepHeading{12}{Likelihood}
            The likelihood $p(o \mid s)$ specifies how hidden states $s$ generate observations $o$. If an observation is not well explained by likely states, surprise rises.
        
        \ProjectDeepHeading{12}{Prior}
            The prior $p(s)$ expresses what the agent expects before seeing the current data. Strong priors can stabilize inference, but if they are inaccurate they can also magnify surprise.
        
        \ProjectDeepHeading{12}{Posterior}
            The posterior $p(s \mid o)$ updates beliefs after data arrive. In variational inference, the agent uses an approximate posterior $q(s)$ and minimizes divergence between $q(s)$ and the true posterior.
        
        \ProjectDeepHeading{12}{Precision}
            Precision is inverse variance. In predictive coding accounts of active inference, precision weights prediction errors and controls their effective influence on belief updating. Precision therefore modulates not the raw mismatch itself, but the importance assigned to that mismatch.
    
    
    \ProjectDeepHeading{11}{Hierarchical Models and Precision-Weighted Prediction Errors}
        Friston's account is hierarchical. Higher levels encode more abstract or slower-changing causes, while lower levels encode faster or more immediate sensory causes. Higher levels send predictions downward; lower levels return prediction errors upward.
        
        A simplified relation is
        \begin{equation}
            \ensuremath{\varepsilon} = o - \hat{o},
        \end{equation}
        where $o$ is observed input and $\hat{o}$ is predicted input. Precision-weighted prediction error is then
        \begin{equation}
            \ensuremath{\xi} = \ensuremath{\Pi} \, \ensuremath{\varepsilon},
        \end{equation}
        with $\ensuremath{\Pi}$ denoting precision. Large $\ensuremath{\xi}$ means the mismatch is both substantial and treated as reliable.
        
        \begin{center}
            \fbox{%
                \parbox{0.9\textwidth}{%
                    \textbf{Hierarchical message flow}\newline
                    Higher-level beliefs generate top-down predictions of lower-level states or sensory causes.\newline
                    Lower levels compare those predictions with actual input and compute prediction errors.\newline
                    Prediction errors are weighted by precision and sent upward to update beliefs.}}
        \end{center}
    
    
    \ProjectDeepHeading{11}{Example Calculations of Surprise}
        If an outcome has probability $p(o)=0.5$, then
        \begin{equation}
            I(o) = -\ln(0.5) \approx 0.693.
        \end{equation}
        If instead $p(o)=0.1$, then
        \begin{equation}
            I(o) = -\ln(0.1) \approx 2.303.
        \end{equation}
        Thus, less probable observations are more surprising.
        
        \begin{table}[H]
            \centering
            \small
            \begin{tabular}{>{\centering\arraybackslash}p{3.0cm} >{\centering\arraybackslash}p{3.0cm} >{\centering\arraybackslash}p{3.0cm}}
                \toprule
                \rowcolor{gray!15}
                \textbf{Observation probability $p(o)$} & \textbf{Surprise $-\ln p(o)$} & \textbf{Interpretation} \\
                \midrule
                0.50                                    & 0.693                         & Mildly surprising       \\
                0.10                                    & 2.303                         & Clearly surprising      \\
                0.01                                    & 4.605                         & Strongly surprising     \\
                \bottomrule
            \end{tabular}
            \caption{Simple numerical examples of surprise.}
        \end{table}
    
    
    \ProjectDeepHeading{11}{Empirical and Simulation Evidence}
        Predictive coding and active inference models explain mismatch responses, repetition suppression, perceptual updating, and several forms of action selection. When the incoming signal violates a learned regularity, prediction error rises and drives belief revision or action.
        
        Simulation work in Friston's papers also shows that agents can reduce expected surprise either by updating hidden-state beliefs or by acting so that sampled inputs better fit prior expectations. In later formulations, expected free energy extends this logic to future-oriented choice under uncertainty.
    
    
    \ProjectDeepHeading{11}{Implications for Perception, Learning, and Psychopathology}
        \textbf{Perception} reduces surprise by updating current beliefs about hidden causes.\newline
        \textbf{Learning} reduces future surprise by changing model parameters and priors.\newline
        \textbf{Action} reduces surprise by changing what the agent samples from the world.\newline
        \textbf{Psychopathology} can be interpreted as maladaptive inference, for example when priors are too strong, precision is misallocated, or the generative model cannot flexibly explain input.
    
    
    \ProjectDeepHeading{11}{Comparison of Factors Affecting Surprise}
        \begin{longtable}{>{\raggedright\arraybackslash}p{3.0cm} >{\raggedright\arraybackslash}p{5.0cm} >{\raggedright\arraybackslash}p{5.2cm}}
            \toprule
            \rowcolor{gray!15}
            \textbf{Component} & \textbf{What changes} & \textbf{Result for surprise} \\
            \midrule
            \endfirsthead
            \toprule
            \rowcolor{gray!15}
            \textbf{Component} & \textbf{What changes} & \textbf{Result for surprise} \\
            \midrule
            \endhead
            Perception & Current beliefs about hidden causes & Lowers surprise by improving posterior fit to present input \\
            Learning & Parameters, structure, or priors of the model & Lowers future surprise by improving the model itself \\
            Action & Sampled sensory states or environmental relation & Lowers surprise by making input conform better to predictions \\
            Attention / precision control & Gain on prediction error signals & Changes which mismatches dominate updating \\
            Novelty & Environmental contingencies or hidden causes & Raises surprise until accommodated by inference or learning \\
            \bottomrule
        \end{longtable}
    
    
    \ProjectDeepHeading{11}{How Surprise Arises and Is Resolved}
        Surprise arises when the agent encounters observations that its current model treats as improbable. This usually appears locally as prediction error and globally as elevated free energy.
        
        In Friston's framework, surprise is mainly resolved in two coupled ways. First, the agent can change its beliefs, predictions, and model parameters through inference and learning. Second, it can change the sensory data it samples by acting on the world. Precision control, attention, and hierarchical message passing are not separate fundamental routes; they are mechanisms that shape how the two main routes operate.
    
    
    \ProjectDeepHeading{11}{References}
        \begin{itemize}[leftmargin=1.8em]
	\item Friston, K. J. (2010). The free-energy principle: a unified brain theory? \textit{Nature Reviews Neuroscience}, 11(2), 127--138.
    \item Friston, K., Kilner, J., and Harrison, L. (2006). A free energy principle for the brain. \textit{Journal of Physiology - Paris}, 100(1--3), 70--87.
    \item Feldman, H., and Friston, K. (2010). Attention, uncertainty, and free-energy. \textit{Frontiers in Human Neuroscience}, 4, 215.
    \item Friston, K., Daunizeau, J., and Kiebel, S. (2009). Reinforcement learning or active inference? \textit{PLoS ONE}, 4(7), e6421.
    \item Parr, T., and Friston, K. J. (2017). Working memory, attention, and salience in active inference. \textit{Scientific Reports}, 7, 14678.
        \end{itemize}

\ProjectDeepHeading{6}{Predictive coding}

% Source: robot/part/body/part/nervous_system/part/central/brain/theory/kind/bayesian/predictive_coding/predictive_coding.tex

\ProjectDeepHeading{9}{Other definitions}
        \cite{spratling-2017-a-review-of-predictive-coding-algorithms} reviews several different algorithms that have all been labeled predictive coding. It argues that predictive coding is better understood as a family of related methods, not a single algorithm. Across these methods, the shared idea is that a system uses a generative model to explain sensory input by reducing prediction error. The paper compares the algorithms in terms of their generative assumptions, optimization rules, and biological plausibility.
        Overall, it presents predictive coding mainly as a framework for perceptual inference in cortical processing.
        
        \clickurlfooter{https://en.wikipedia.org/wiki/Predictive_coding} In neuroscience, predictive coding (also known as predictive processing) is a theory of brain function which postulates that the brain is constantly generating and updating a ``mental model'' of the environment. According to the theory, such a mental model is used to predict input signals from the senses that are then compared with the actual input signals from those senses. Predictive coding is member of a wider set of theories that follow the Bayesian brain hypothesis.
        \begin{figure}[H]
            \centering
            \ProjectSafeIncludeGraphics[width=0.9\textwidth]{robot/part/body/part/nervous_system/part/central/brain/theory/kind/bayesian/predictive_coding/predictive_coding.png}
        \end{figure}
    
    
    \ProjectDeepHeading{9}{Parts}
        
        \ProjectDeepHeading{10}{Local model}
        
        \ProjectDeepHeading{10}{Prediction model}
        
        \ProjectDeepHeading{10}{Prediction}
        
        \ProjectDeepHeading{10}{Prediction error}
    
    
    %
    %

\ProjectDeepHeading{5}{Dynamical systems theories}

% Source: robot/part/body/part/nervous_system/part/central/brain/theory/kind/dynamical_systems_theories/dynamical_systems_theories.tex

\cite{kitzbichler-2009-broadband-criticality-of-human-brain-network-synchronization} argues that large-scale human brain activity shows signatures of critical dynamics across multiple frequency bands. Using human brain network synchronization data and computational modeling, it suggests that the brain operates near a critical point rather than in a purely ordered or disordered regime. This near-critical state helps explain scale-free temporal structure, metastability, and flexible coordination across distributed brain systems.

\ProjectDeepHeading{5}{Efficient coding}

% Source: robot/part/body/part/nervous_system/part/central/brain/theory/kind/efficient_coding/efficient_coding.tex

Efficient coding\cite{barlow-1961-possible-principles-underlying-the-transformation-of-sensory-messages} is a theoretical model of sensory neuroscience in the brain. Within the brain, neurons communicate with one another by sending electrical impulses referred to as action potentials or spikes. Barlow hypothesized that the spikes in the sensory system formed a neural code for efficiently representing sensory information. By efficient it is understood that the code minimized the number of spikes needed to transmit a given signal.

\ProjectDeepHeading{5}{Hebbian}

% Source: robot/part/body/part/nervous_system/part/central/brain/theory/kind/hebbian/hebbian.tex

\clickurlfooter{https://en.wikipedia.org/wiki/Hebbian_theory} The theory is often summarized as``Neurons that fire together, wire together``
    %
    %

\ProjectDeepHeading{5}{Neural darwinism}

% Source: robot/part/body/part/nervous_system/part/central/brain/theory/kind/neural_darwinism/neural_darwinism.tex

\clickurlfooter{https://en.wikipedia.org/wiki/Neural_Darwinism} suggests the brain develops many neuronal groups and circuits, and experience strengthens some of them while others weaken.
    So, brain function emerges through a kind of selection process, where useful neural patterns are kept because they help perception, action, and adaptation. It is called “Darwinism” because, like natural selection in biology, variation exists first and selection acts on that variation.
    %
    %

\ProjectDeepHeading{5}{Optimal control theory}

% Source: robot/part/body/part/nervous_system/part/central/brain/theory/kind/optimal_control_theory/optimal_control_theory.tex

\clickurlfooter{https://en.wikipedia.org/wiki/Optimal_control} Optimal control theory is a branch of control theory that deals with finding a control for a dynamical system over a period of time such that an objective function is optimized
    
    \cite{diedrichsen-2009-the-coordination-of-movement-optimal-feedback-control-and-beyond} explains how optimal control theory, especially optimal feedback control, helps account for flexible and task-dependent movement coordination. It argues that movement is organized by distributing control across multiple effectors, such as muscles, joints, and limbs, in a way that fits task demands.
    The authors show that this framework can explain patterns of feedback correction, movement variability, and coordination across effectors. They also point out that hierarchical control and movement initiation still need further development for a fuller theory of voluntary motor control.

\ProjectDeepHeading{5}{Reinforcement learning}

% Source: robot/part/body/part/nervous_system/part/central/brain/theory/kind/reinforcement_learning/reinforcement_learning.tex

\clickurlfooter{https://en.wikipedia.org/wiki/Reinforcement} defines reinforcement refers to consequences that increase the likelihood of an organism's future behavior, typically in the presence of a particular antecedent stimulus.
    %
    %

\ProjectDeepHeading{1}{Peripheral}

% Source: robot/part/body/part/nervous_system/part/peripheral/peripheral.tex

The PNS consists of nerves and ganglia, which lie outside the brain and the spinal cord.[1] The main function of the PNS is to connect the CNS to the limbs and organs, essentially serving as a relay between the brain and spinal cord and the rest of the body.\url{https://en.wikipedia.org/wiki/Peripheral_nervous_system}
    \begin{enumerate}
        \item Somatic sensory-motor nerves that connect the central nervous system to the limbs and organs
        \item Autonomic nervous system (which regulates body processes such as digestion and heart rate)
        \item Enteric nervous system (which controls the gastrointestinal system)
    \end{enumerate}

\ProjectDeepHeading{1}{Sesnsory}

% Source: robot/part/body/part/nervous_system/part/sesnsory/sensory.tex

Sensory neurons, also known as afferent neurons, are in the nervous system which convert a specific type of stimulus, via their receptors, into action potentials or graded receptor potentials.
    \url{https://en.wikipedia.org/wiki/Sensory_neuroscience}

\ProjectDeepHeading{2}{Modality}

\ProjectDeepHeading{3}{Stimulus}

% Source: robot/part/body/part/nervous_system/part/sesnsory/modality/stimulus/stimulus.tex

In physiology, a stimulus[1] is a change in a living thing's internal or external environment.
    \url{https://en.wikipedia.org/wiki/Stimulus_(physiology)}

\ProjectDeepHeading{4}{Modality}

% Source: robot/part/body/part/nervous_system/part/sesnsory/modality/stimulus/modality/modality.tex

Stimulus modality, also called sensory modality, is one aspect of a stimulus or what is perceived after a stimulus.
    \url{https://en.wikipedia.org/wiki/Stimulus_modality}
    % 
    %

\ProjectDeepHeading{5}{Kind}

% Source: robot/part/body/part/nervous_system/part/sesnsory/modality/stimulus/modality/kind/antecedent.tex

\clickurlfooter{https://en.wikipedia.org/wiki/Antecedent_(behavioral_psychology)} defines an antecedent is a stimulus that cues an organism to perform a learned behavior. When an organism perceives an antecedent stimulus, it behaves in a way that maximizes reinforcing consequences and minimizes punishing consequences
    %
    %

\ProjectDeepHeading{2}{Multi}

\ProjectDeepHeading{3}{Binding problem}

% Source: robot/part/body/part/nervous_system/part/sesnsory/multi/binding_problem/binding_problem.tex

The binding problem refers to the overall encoding of our brain circuits for the combination of decisions, actions, and perception. It is considered a "problem" because no complete model exists \seeherefooter{https://en.wikipedia.org/wiki/Binding_problem}.

\ProjectDeepHeading{3}{Integration}

\ProjectDeepHeading{4}{Cue}

% Source: robot/part/body/part/nervous_system/part/sesnsory/multi/integration/cue/cue_integration.tex

\begin{itemize}
        \item Statistically Optimal Multisensory Cue Integration: A Practical Tutorial \seeherefooter{https://www.researchgate.net/publication/282575320_Statistically_Optimal_Multisensory_Cue_Integration_A_Practical_Tutorial}
        \item An Analysis of a Ring Attractor Model for Cue Integration \seeherefooter{https://www.researchgate.net/publication/326227083_An_Analysis_of_a_Ring_Attractor_Model_for_Cue_Integration}
        \item Cue integration as a common mechanism for action and outcome bindings \seeherefooter{https://www.sciencedirect.com/science/article/abs/pii/S0010027720302420}
    \end{itemize}
    \seeherefooter{https://pmc.ncbi.nlm.nih.gov/articles/PMC3130032/}

\ProjectDeepHeading{4}{Kind}

\ProjectDeepHeading{5}{Bayesian}

% Source: robot/part/body/part/nervous_system/part/sesnsory/multi/integration/kind/bayesian/bayesian.tex

This study investigates how visual optic-flow and vestibular signals are combined when humans estimate their direction of self-motion. Eight participants completed heading-discrimination tasks under visual-only, vestibular-only, and combined conditions containing spatial conflicts of $\pm 6^{\circ}$ or $\pm 10^{\circ}$.
    Combining the two sensory cues reduced perceptual variance approximately as predicted by maximum-likelihood estimation, even when the cues indicated conflicting heading directions. However, the observed cue weights were not determined solely by the reliabilities measured in the unimodal conditions, because participants consistently gave greater weight to vestibular information. The results suggest that visual--vestibular integration is statistically near-optimal in terms of precision, but that its systematic vestibular bias requires a Bayesian prior or a more comprehensive model to explain \cite{butler-2010-bayesian-integration-of-visual-and-vestibular-signals-for-heading}.

\ProjectDeepHeading{6}{Kalman}

% Source: robot/part/body/part/nervous_system/part/sesnsory/multi/integration/kind/bayesian/kalman/kalman.tex

\cite{kalman-1960-a-new-approach-to-linear-filtering-and-prediction-problems} reformulates linear filtering and prediction using state-space models and conditional estimation for the first time which was later coned as Kalman filter.
    It derives recursive equations for the optimal estimate and for the covariance of the estimation error.
    The method applies to stationary and nonstationary processes, as well as finite-memory and growing-memory filters.
    The optimal filter coefficients are obtained directly from the recursively updated error covariance matrix.
    The paper also establishes a duality between optimal estimation and the noise-free optimal regulator problem.

\ProjectDeepHeading{7}{Biology}

% Source: robot/part/body/part/nervous_system/part/sesnsory/multi/integration/kind/bayesian/kalman/biology/biology.tex

\cite{deneve-2007-optimal-sensorimotor-integration-in-recurrent-cortical-networks-a-neural-implementation-of-kalman-filters} proposes recurrent basis-function networks as a biologically plausible implementation of Kalman filtering in cortical circuits.
    The network combines noisy sensory feedback with predictions generated by an internal model of sensorimotor dynamics.
    Simulations of visual-object tracking and arm-state estimation show that the network approaches the accuracy of an optimal Kalman filter.
    Its recurrent connections predict future sensorimotor states, while reliability-dependent sensory gains determine how strongly feedback influences the estimate.
    The model also predicts context-dependent changes in neuronal tuning curves, including variations caused by arm position, movement velocity, and sensory reliability.

\ProjectDeepHeading{4}{Part}

\ProjectDeepHeading{5}{Visual vestibular}

% Source: robot/part/body/part/nervous_system/part/sesnsory/multi/integration/part/visual_vestibular/visual_vestibular.tex

This study investigates how visual optic-flow and vestibular signals are combined when humans estimate their direction of self-motion.
    Eight participants completed heading-discrimination tasks under visual-only, vestibular-only, and combined conditions containing spatial conflicts of $\pm 6^{\circ}$ or $\pm 10^{\circ}$.
    Combining the two sensory cues reduced perceptual variance approximately as predicted by maximum-likelihood estimation, even when the cues indicated conflicting heading directions.
    However, the observed cue weights were not determined solely by the reliabilities measured in the unimodal conditions, because participants consistently gave greater weight to vestibular information.
    The results suggest that visual--vestibular integration is statistically near-optimal in terms of precision, but that its systematic vestibular bias requires a Bayesian prior or a more comprehensive model to explain \cite{butler-2010-bayesian-integration-of-visual-and-vestibular-signals-for-heading}.

\ProjectDeepHeading{3}{Processing}

% Source: robot/part/body/part/nervous_system/part/sesnsory/multi/processing/processing.tex

\begin{itemize}
        \item Structural Basis of Multisensory Processing \seeherefooter{https://www.ncbi.nlm.nih.gov/books/NBK92880}
    \end{itemize}

\ProjectDeepHeading{2}{Part}

\ProjectDeepHeading{3}{Neuron}

% Source: robot/part/body/part/nervous_system/part/sesnsory/part/neuron/neuron.tex

\url{https://en.wikipedia.org/wiki/Sensory_nervous_system}
    
    
    % 
    % 



    
    
    
    Sensory neurons, also known as afferent neurons, are in the nervous system which convert a specific type of stimulus, via their receptors, into action potentials or graded receptor potentials .

% Source: robot/part/body/part/nervous_system/part/sesnsory/part/neuron/from_action_potential_to_perception.tex

\ProjectDeepHeading{6}{from touch to perception}

    \cite{romo-2020-turning-touch-into-perception} claims the flow of data from stimuli sensing until making perception is as follows
    
    Touch stimulus → ventral posterior lateral nucleus of the somatosensory thalamus → primary somatosensory cortex → secondary somatosensory cortex → premotor and frontal cortical areas → perceptual report → perception

% Source: robot/part/body/part/nervous_system/part/sesnsory/part/neuron/processing.tex

\url{https://en.wikipedia.org/wiki/Stimulus_modality}
    % 
    %

% Source: robot/part/body/part/nervous_system/part/sesnsory/part/neuron/sense.tex

\url{https://en.wikipedia.org/wiki/Sense}
    % 
    %

\ProjectDeepHeading{3}{Receptor}

% Source: robot/part/body/part/nervous_system/part/sesnsory/part/receptor/receptor.tex

\textbf{This entry is irrevalent to sensory receptor}  In biochemistry and pharmacology, receptors are proteins that receive and transduce signals that may be integrated into biological systems \seeherefooter{https://en.wikipedia.org/wiki/Receptor_(biochemistry)}.
    
    A sensory receptor converts the energy in a stimulus into an electrical signal \seeherefooter{https://en.wikipedia.org/wiki/Transduction_(physiology)}
    
    Receptors are broadly split into two main categories: exteroceptors, which receive external sensory stimuli, and interoceptors, which receive internal sensory stimuli.
    %
    %

\ProjectDeepHeading{4}{Kind}

% Source: robot/part/body/part/nervous_system/part/sesnsory/part/receptor/kind/imu.tex

IMU
    
    \ProjectDeepHeading{8}{Giroscope}
        aa
    
    \ProjectDeepHeading{8}{Accelometer}
        aa
    
    \ProjectDeepHeading{8}{Magnometer}
        aa
    %
    %

% Source: robot/part/body/part/nervous_system/part/sesnsory/part/receptor/kind/kinds.tex

\ProjectDeepHeading{8}{Motion sensors}
    
    \ProjectDeepHeading{8}{GNSS}
        aa
    
    \ProjectDeepHeading{8}{IMU}
        IMU
        
        \ProjectDeepHeading{9}{Giroscope}
            aa
        
        \ProjectDeepHeading{9}{Accelometer}
            aa
        
        \ProjectDeepHeading{9}{Magnometer}
            aa
    
    \ProjectDeepHeading{8}{Infrared light (passive and active sensors)}
    
    \ProjectDeepHeading{8}{Visible light (video and camera systems)}
    
    \ProjectDeepHeading{8}{Radio frequency energy (radar, microwave and tomographic motion detection)}
    
    \ProjectDeepHeading{8}{Sound (microphones, other acoustic sensors)}
    
    \ProjectDeepHeading{8}{Kinetic energy (triboelectric, seismic, and inertia-switch sensors)}
    
    \ProjectDeepHeading{8}{Magnetism (magnetic sensors, magnetometers)}
    
    \ProjectDeepHeading{8}{Wi-Fi Signals (WiFi Sensing)}
    
    
    \ProjectDeepHeading{7}{Barometer}
    something
    
    \ProjectDeepHeading{7}{Camera}
    %
    %

\ProjectDeepHeading{5}{Proprioceptive}

% Source: robot/part/body/part/nervous_system/part/sesnsory/part/receptor/kind/proprioceptive/proprioceptive.tex

\ProjectDeepHeading{8}{Proprioception}
        
        \clickurlfooter{https://en.wikipedia.org/wiki/Proprioception} defines is the sense of self-movement, force, and body position.
        
        \cite{tuthill-2018-proprioception} defines proprioception as the sense that lets the body monitor position, movement, and force, even when we are not consciously aware of it.
        The article explains that this sense is essential for coordinated action, and shows that losing proprioception can severely disrupt even basic movement.
        It also reviews the sensory receptors and neural pathways that allow the nervous system to build an internal representation of the body.

\ProjectDeepHeading{4}{Potential}

% Source: robot/part/body/part/nervous_system/part/sesnsory/part/receptor/potential/potential.tex

A receptor potential is a type of graded potential, is the transmembrane potential difference produced by activation of a sensory receptor \seeherefooter{https://en.wikipedia.org/wiki/Receptor_potential}.

\ProjectDeepHeading{4}{Temporal alignment}

% Source: robot/part/body/part/nervous_system/part/sesnsory/part/receptor/temporal_alignment/temporal_alignment.tex

Aligning multi-rate sensors

\ProjectDeepHeading{5}{Kind}

% Source: robot/part/body/part/nervous_system/part/sesnsory/part/receptor/temporal_alignment/kind/gaussian_process_regression.tex

\cite{barfoot-2014-batch-continuous-time-trajectory-estimation-as-exactly-sparse-gaussian-process-regression} propose a continuous-time approach for batch trajectory estimation using Gaussian Process regression. The trajectory is treated as a one-dimensional Gaussian Process where time is the independent variable. This is useful for asynchronous or multi-rate robotic sensors, such as scanning laser rangefinders, rolling-shutter cameras, and other sensors that produce observations at different timestamps. Instead of using direct linear or spline interpolation between discrete poses, the method represents the whole trajectory as a continuous-time probabilistic function. The key contribution is that, for Gaussian Process priors generated by linear time-varying stochastic differential equations, the inverse kernel matrix becomes exactly sparse and block-tridiagonal. This allows efficient GP regression and interpolation, making it possible to query the robot state at any time of interest with low computational cost.

% Source: robot/part/body/part/nervous_system/part/sesnsory/part/receptor/temporal_alignment/kind/interpolation.tex

\cite{babu-2005-comparative-study-of-interpolation-techniques-for-ultra-tight-integration-of-gps-ins-pl-sensors} study interpolation techniques for ultra-tight GPS/INS/Pseudolite integration. The problem is that the navigation filter usually runs at a low rate, typically between 1 and 100 Hz, while the GPS carrier tracking loop requires updates at about 1000 Hz. Running the Kalman filter at 1000 Hz would be computationally expensive, so the authors adopt an interpolation-based solution. Their method increases the sampling rate of low-rate INS-derived Doppler measurements using multirate signal processing. They compare Polyphase and CIC interpolators and report that CIC is computationally simpler because it avoids multipliers and gives better preliminary performance in their experiment.

% Source: robot/part/body/part/nervous_system/part/sesnsory/part/receptor/temporal_alignment/kind/spline_interpolation.tex

\seeherefooter{https://en.wikipedia.org/wiki/Spline_interpolation}
    
    \cite{lovegrove-2013-spline-fusion-a-continuous-time-representation-for-visual-inertial-fusion-with-application-to-rolling-shutter-cameras} propose a continuous-time spline-based framework for visual-inertial SLAM and calibration. Instead of representing the sensor trajectory only at discrete timestamps, they model it as a continuous curve using cubic B-splines. This allows measurements from unsynchronized and high-rate sensors, such as cameras and IMUs, to be fused in a common temporal framework. The method is also useful for rolling-shutter cameras, because each image scanline can be associated with its own measurement time.

\ProjectDeepHeading{2}{Transduction}

% Source: robot/part/body/part/nervous_system/part/sesnsory/transduction/transduction.tex

In physiology, transduction is the translation of arriving stimulus into an action potential by a sensory receptor \seeherefooter{https://en.wikipedia.org/wiki/Transduction_(physiology)}.
    
    In this video the lecturer explains the simplest possible course of events signal in the form of a molcule evolves to a electric signal in the form of a receptor potential that may causes a spiking in the connected sensory neuron.
    \seeherefooter{https://www.youtube.com/watch?v=d7kitk-knjA}
    This following Figure is derived from the aformentioned video that shows the stages of sensory transduction
    \begin{figure}[H]
        \centering
        \ProjectSafeIncludeGraphics[width=0.9\textwidth]{robot/part/body/part/nervous_system/part/sesnsory/transduction/transduction_stages.png}
    \end{figure}
    
%    
%

\ProjectDeepHeading{1}{Shared}

\ProjectDeepHeading{2}{Motor system}

% Source: robot/part/body/part/nervous_system/part/shared/motor_system/motor_system.tex

\url{https://en.wikipedia.org/wiki/Motor_system}
    % 
    %

% Source: robot/part/body/part/nervous_system/part/shared/motor_system/motor_skill.tex

\url{https://en.wikipedia.org/wiki/Motor_skill}
    % 
    %

\ProjectDeepHeading{2}{Nerve}

% Source: robot/part/body/part/nervous_system/part/shared/nerve/nerve.tex

A bundle of axons (nerve fibers) wrapped together by connective tissue that carries electrical signals between the central nervous system and the body. \url{https://en.wikipedia.org/wiki/Nerve}
    \ProjectSafeIncludeGraphics[scale=0.2]{robot/part/body/part/nervous_system/part/shared/nerve/nerve.png}

\ProjectDeepHeading{2}{Neuron}

% Source: robot/part/body/part/nervous_system/part/shared/neuron/neuron.tex

A neuron is an excitable system \seeherefooter{https://mcb.berkeley.edu/courses/mcb137/exercises/Lesson\%206\%20E&O\%20Systems.pdf} and \seeherefooter{https://en.wikipedia.org/wiki/Oscillation} and \seeherefooter{https://en.wikipedia.org/wiki/Neural_oscillation} that lives in a exciatble environment \seeherefooter{https://en.wikipedia.org/wiki/Excitable_medium}.

\ProjectDeepHeading{3}{Kind}

% Source: robot/part/body/part/nervous_system/part/shared/neuron/kind/kinds.tex

\ProjectDeepHeading{6}{Sensory Neurons}
        These neurons tell the rest of the brain about the external and internal environment.
        
    \ProjectDeepHeading{6}{Motor (and Other Output) Neurons}
        Motor neurons contract muscles and mediate behavior, and other output neurons stimulate glands and organs.
        
    \ProjectDeepHeading{6}{Projection Neurons}
        Projection (communication) neurons transmit signals from one brain area to another.
        
    \ProjectDeepHeading{6}{Interneurons}
        The vast majority of neurons in vertebrates are interneurons involved in local computations. Computational interneurons extract and process information coming in from the senses, compare that information to what is in memory, and use the information to plan and execute behavior. Each of the several hundred distinguishable brain regions contains several dozen distinct types or classes of computational interneurons that mediate the function of that brain area.

\ProjectDeepHeading{4}{Interneuron}

% Source: robot/part/body/part/nervous_system/part/shared/neuron/kind/interneuron/interneuron.tex

Interneurons (also called internuncial neurons, association neurons, connector neurons, or intermediate neurons) are neurons that are not specifically motor neurons or sensory neurons. \url{https://en.wikipedia.org/wiki/Interneuron}
    % 
    %

\ProjectDeepHeading{4}{Motor}

% Source: robot/part/body/part/nervous_system/part/shared/neuron/kind/motor/motor.tex

A motor neuron (or motoneuron), also known as efferent neuron is a neuron that allows for both voluntary and involuntary movements of the body through muscles and glands \seeherefooter{https://en.wikipedia.org/wiki/Motor_neuron}.

% Source: robot/part/body/part/nervous_system/part/shared/neuron/kind/motor/goal.tex

A motor goal is a neurally planned motor outcome that is used to organize motor control.
    Motor goals are experimentally shown to exist since planned movements can when disrupted adjust to achieve their planned outcome. If, for example, a person makes a movement of their hand to touch or grasp something and unexpected their arm is pushed – their brain automatically reorganizes the movement so it so achieves its intended aim. This also occurs if an arm is perturbed which results in an automatic correction that enables it to fulfill its planned spatial-temporal target \seeherefooter{https://en.wikipedia.org/wiki/Motor_goal}

\ProjectDeepHeading{3}{Model}

% Source: robot/part/body/part/nervous_system/part/shared/neuron/model/model.tex

Biological neuron models, also known as spiking neuron models,[1] are mathematical descriptions of the conduction of electrical signals in neurons \seeherefooter{https://en.wikipedia.org/wiki/Biological_neuron_model}

\ProjectDeepHeading{3}{Part}

\ProjectDeepHeading{4}{Action potential}

% Source: robot/part/body/part/nervous_system/part/shared/neuron/part/action_potential/action_potential.tex

\seeherefooter{https://en.wikipedia.org/wiki/Action_potential}
    Also called as spike
    
    
    
    
    
    \ProjectSafeIncludeGraphics[width=\paperwidth]{robot/part/body/part/nervous_system/part/shared/neuron/part/action_potential/action_potential.png}
    %
    %

% Source: robot/part/body/part/nervous_system/part/shared/neuron/part/action_potential/sorting.tex

Spike sorting is to find out from which neuron a spike is emitted.
    A review on spike sorting/classification is \cite{lewicki-1998-a-review-of-methods-for-spike-sorting-the-detection-and-classification-of-neural-action-potentials} has reviewed different spike sorting methods
    
    \clickurlfooter{https://pub.ista.ac.at/~gtkacik/PLOSOne_SpikeSorting.pdf}

% Source: robot/part/body/part/nervous_system/part/shared/neuron/part/action_potential/train.tex

The temporal sequence of action potentials generated by a neuron is called its ``spike train``
    \cite{vandrongelen-2007-spike-train-analysis}

\ProjectDeepHeading{5}{Cable theory}

% Source: robot/part/body/part/nervous_system/part/shared/neuron/part/action_potential/cable_theory/cable_theory.tex

\seeherefooter{https://en.wikipedia.org/wiki/Cable_theory}

\ProjectDeepHeading{5}{Encoding and decoding}

\ProjectDeepHeading{6}{Kind}

% Source: robot/part/body/part/nervous_system/part/shared/neuron/part/action_potential/encoding_and_decoding/kind/neural_coding.tex

Neural coding (or neural representation) refers to the relationship between a stimulus and its respective neuronal responses, and the signaling relationships among networks of neurons in an ensemble. \seeherefooter{https://en.wikipedia.org/wiki/Neural\_coding}
    
    \cite{shimazaki-2025-neural-coding-foundational-concepts-statistical-formulations-and-recent-advances} explains neural coding as the process by which stimuli are transformed into neural activity that supports perception and behavior, combining classical concepts with statistical encoding and decoding frameworks.

% Source: robot/part/body/part/nervous_system/part/shared/neuron/part/action_potential/encoding_and_decoding/kind/representation_with_autoencoder.tex

\cite{granley-2022-hybrid-neural-autoencoders} propose a topographic variational autoencoder to model proprioceptive neural coding in somatosensory cortex.
    The model is trained on natural human arm-movement kinematics, not directly on neural recordings, but reproduces several coding properties observed in monkey cortical neurons.
    The paper supports the idea that autoencoders can be used to model sensory representations and neural coding when the latent space is interpreted as a biologically constrained neural population.
    \begin{figure}[H]
        \centering
        \ProjectSafeIncludeGraphics[width=0.9\textwidth]{robot/part/body/part/nervous_system/part/shared/neuron/part/action_potential/encoding_and_decoding/kind/granley-2022-hybrid-neural-autoencoders_figure_5.png}
    \end{figure}

\ProjectDeepHeading{7}{Decoding}

% Source: robot/part/body/part/nervous_system/part/shared/neuron/part/action_potential/encoding_and_decoding/kind/decoding/decoding.tex

\url{https://en.wikipedia.org/wiki/Neural_decoding}
    Neural decoding is a neuroscience field concerned with the hypothetical reconstruction of sensory and other stimuli from information that has already been encoded and represented in the brain by networks of neurons.
    
    in

\ProjectDeepHeading{7}{Encoding}

% Source: robot/part/body/part/nervous_system/part/shared/neuron/part/action_potential/encoding_and_decoding/kind/encoding/efficient.tex

The efficient coding hypothesis was proposed by Horace Barlow in 1961 as a theoretical model of sensory neuroscience in the brain. Within the brain, neurons communicate with one another by sending electrical impulses referred to as action potentials or spikes.
    
    Barlow hypothesized that the spikes in the sensory system formed a neural code for efficiently representing sensory information. By efficient it is understood that the code minimized the number of spikes needed to transmit a given signal. This is somewhat analogous to transmitting information across the internet, where different file formats can be used to transmit a given image. Different file formats require different numbers of bits for representing the same image at a given distortion level, and some are better suited for representing certain classes of images than others. According to this model, the brain is thought to use a code which is suited for representing visual and audio information which is representative of an organism's natural environment .
    \seeherefooter{https://en.wikipedia.org/wiki/Efficient_coding_hypothesis}

\ProjectDeepHeading{8}{Code}

\ProjectDeepHeading{9}{Kind}

\ProjectDeepHeading{10}{Population}

% Source: robot/part/body/part/nervous_system/part/shared/neuron/part/action_potential/encoding_and_decoding/kind/encoding/code/kind/population/population.tex

Population coding is a method to represent stimuli by using the joint activities of a number of neurons. In population coding, each neuron has a distribution of responses over some set of inputs, and the responses of many neurons may be combined to determine some value about the inputs \seeherefooter{https://en.wikipedia.org/wiki/Neural_coding\#Population_coding}.
    
    Population coding is a method to represent stimuli by using the joint activities of a number of neurons. In population coding, each neuron has a distribution of responses over some set of inputs, and the responses of many neurons may be combined to determine some value about the inputs.
    \seeherefooter{https://en.wikipedia.org/wiki/Neural_coding\#Population_coding}
    
    \cite{shimazaki-2025-neural-coding-foundational-concepts-statistical-formulations-and-recent-advances} emphasizes that neural information is often represented at population level wehere
    
    \begin{itemize}
        \item firing rate
        \item temporal structure
        \item correlations between neurons
    \end{itemize}
    jointly affect how accurately stimuli can be encoded and decoded. The paper also reviews recent large-scale recording studies, showing that understanding information limits in neural populations now depends on analyzing high-dimensional covariance structure rather than only small-scale pairwise correlations.
    \cite{shimazaki-2025-neural-coding-foundational-concepts-statistical-formulations-and-recent-advances}

\ProjectDeepHeading{10}{Sparse}

% Source: robot/part/body/part/nervous_system/part/shared/neuron/part/action_potential/encoding_and_decoding/kind/encoding/code/kind/sparse/sparse.tex

The sparse code is when each item is encoded by the strong activation of a relatively small set of neurons. For each item to be encoded, this is a different subset of all available neurons \seeherefooter{https://en.wikipedia.org/wiki/Neural_coding\#Sparse_coding}.

\ProjectDeepHeading{10}{Temporal}

% Source: robot/part/body/part/nervous_system/part/shared/neuron/part/action_potential/encoding_and_decoding/kind/encoding/code/kind/temporal/temporal.tex

When precise spike timing or high-frequency firing-rate fluctuations are found to carry information, the neural code is often identified as a temporal code  \seeherefooter{https://en.wikipedia.org/wiki/Neural_coding\#Temporal_coding}.

\ProjectDeepHeading{5}{Kind}

\ProjectDeepHeading{6}{Graded or binary}

\ProjectDeepHeading{7}{Kind}

% Source: robot/part/body/part/nervous_system/part/shared/neuron/part/action_potential/kind/graded_or_binary/kind/all_or_none.tex

\url{https://en.wikipedia.org/wiki/All-or-none_law}
    In physiology, the all-or-none law (sometimes the all-or-none principle or all-or-nothing law) is the principle that if a single nerve fibre is stimulated, it will always give a maximal response and produce an electrical impulse of a single amplitude.

\ProjectDeepHeading{8}{Graded}

% Source: robot/part/body/part/nervous_system/part/shared/neuron/part/action_potential/kind/graded_or_binary/kind/graded/graded.tex

Graded potentials are changes in membrane potential that vary according to the size of the stimulus, as opposed to being all-or-none.
    \url{https://en.wikipedia.org/wiki/Graded_potential}
    \begin{figure}[H]
        \centering
        \ProjectMissingAsset{/home/donkarlo/Dropbox/repo/nd\_robotic\_ai\_learning\_project/src/robot/part/body/part/nervous\_system/action\_potential/graded\_or\_binary/kinds/graded\_potential/graded\_potential.png}
    \end{figure}
    
    \cite{juusola-2007-coding-with-spike-shapes-and-graded-potentials-in-cortical-networks} is mainly about action-potential waveform coding. Graded potentials are important because they modulate spike shapes and synaptic transmission. So the paper is about the interaction between digital spikes and analog graded membrane potentials.
    
    
    \ProjectDeepHeading{11}{Waveform Information in Graded Potentials and Graded Spikes}
        
        The safer term in the literature is usually graded potential or graded synaptic transmission, not graded action potential. A classical action potential is often treated as an all-or-none event, but several studies show that spike waveform can still vary and modulate synaptic transmission. In particular, spike amplitude, spike width, and spike shape can affect calcium entry and neurotransmitter release \citep{zbili-2019-past-and-future-of-analog-digital-modulation-of-synaptic-transmission}. A more direct case of amplitude-dependent graded spiking is reported in cerebrospinal-fluid-contacting neurons, where different spike amplitudes can signal different neurotransmitter inputs \citep{johnson-2023-graded-spikes-differentially-signal-neurotransmitter-input-in-cerebrospinal-fluid-contacting-neurons-of-the-mouse-spinal-cord}.
        
        \ProjectDeepHeading{12}{Amplitude or Peak Size}
            
            The amplitude of a graded voltage signal can encode stimulus strength, contrast, release strength, or, in specific cells, the type of input. At a visual synapse, contrast can be transmitted not only by the frequency of release events but also by their amplitude \citep{james-2019-an-amplitude-code-transmits-information-at-a-visual-synapse}. In mammalian cone photoreceptors, membrane voltage continuously represents illumination and controls graded vesicle fusion at ribbon synapses \citep{grabner-2023-mechanisms-of-simultaneous-linear-and-nonlinear-computations-at-the-mammalian-cone-photoreceptor-synapse}. In cerebrospinal-fluid-contacting neurons, smaller and larger graded spikes can distinguish purinergic and cholinergic inputs \citep{johnson-2023-graded-spikes-differentially-signal-neurotransmitter-input-in-cerebrospinal-fluid-contacting-neurons-of-the-mouse-spinal-cord}.
            
            \begin{equation}
                A
                =
                \max_t |V(t)-V_{\mathrm{rest}}|
            \end{equation}
            
            Therefore, the amplitude can be interpreted as
            
            \begin{equation}
                A
                \rightarrow
                \text{stimulus intensity, contrast, release strength, or input identity}.
            \end{equation}
        
        \ProjectDeepHeading{12}{Polarity and Sign}
            
            The polarity of a graded potential is related to excitation and inhibition. Excitatory postsynaptic potentials increase the probability of postsynaptic firing, whereas inhibitory postsynaptic potentials decrease it. However, the biological effect depends on the reversal potential and its relation to the threshold, not only on the visual direction of voltage change \citep{purves-2001-excitatory-and-inhibitory-postsynaptic-potentials,purves-2001-summation-of-synaptic-potentials}.
            
            \begin{equation}
                \ensuremath{\Delta} V(t)
                =
                V(t)-V_{\mathrm{rest}}
            \end{equation}
            
            A safe interpretation is
            
            \begin{equation}
                \ensuremath{\Delta} V(t)
                \rightarrow
                \text{excitatory or inhibitory tendency, depending on reversal potential and threshold}.
            \end{equation}
        
        \ProjectDeepHeading{12}{Rise Time, Slope, and Temporal Precision}
            
            The slope and rise time of a graded membrane potential can encode temporal properties of the stimulus. In graded-potential neurons of the blowfly visual system, active membrane properties affect how precisely stimulus velocity is encoded in membrane potential \citep{haag-1998-active-membrane-properties-and-signal-encoding-in-graded-potential-neurons}.
            
            \begin{equation}
                S_{\uparrow}
                =
                \max_t \frac{dV(t)}{dt}
            \end{equation}
            
            Thus, slope and rise time can be interpreted as
            
            \begin{equation}
                S_{\uparrow},\ t_{\mathrm{rise}}
                \rightarrow
                \text{stimulus dynamics, velocity, and temporal precision}.
            \end{equation}
        
        \ProjectDeepHeading{12}{Duration, Decay, and Adaptation}
            
            The duration and decay of a graded potential can carry information about stimulus duration, adaptation, and temporal integration. In zebrafish lateral-line hair cells, different cells show different sensitivities, operating ranges, and adaptive properties, allowing weak stimuli to be encoded while preserving dynamic range for stronger deflections \citep{pichler-2019-the-transfer-characteristics-of-hair-cells-encoding-mechanical-stimuli-in-the-lateral-line-of-zebrafish}.
            
            \begin{equation}
                W_{\ensuremath{\theta}}
                =
                t_2-t_1
            \end{equation}
            
            \begin{equation}
                V(t_1)
                =
                V(t_2)
                =
                V_{\mathrm{rest}}+\ensuremath{\theta} A
            \end{equation}
            
            A safe interpretation is
            
            \begin{equation}
                W_{\ensuremath{\theta}},\ \ensuremath{\tau}_{\mathrm{decay}}
                \rightarrow
                \text{stimulus duration, adaptation, temporal integration, or synaptic release dynamics}.
            \end{equation}
        
        \ProjectDeepHeading{12}{Area Under the Waveform}
            
            The area under the voltage deviation summarizes amplitude and duration together. This should be treated as a mathematical descriptor of the waveform, not as a claim that the neuron explicitly computes this scalar as a separate biological code.
            
            \begin{equation}
                Q_V
                =
                \int_{t_0}^{t_1}
                \left(V(t)-V_{\mathrm{rest}}\right)
                dt
            \end{equation}
            
            A safe interpretation is
            
            \begin{equation}
                Q_V
                \rightarrow
                \text{total voltage deviation over time}.
            \end{equation}
        
        \ProjectDeepHeading{12}{Timing, Latency, and Phase}
            
            The timing of the graded response can encode when the stimulus occurred and how quickly the sensory cell or synapse responded. In zebrafish lateral-line hair cells, sustained outputs, transient outputs, and transient return-to-rest signals show that timing and response time course are part of mechanical stimulus encoding \citep{pichler-2019-the-transfer-characteristics-of-hair-cells-encoding-mechanical-stimuli-in-the-lateral-line-of-zebrafish}.
            
            \begin{equation}
                L
                =
                t_{\mathrm{onset}}-t_{\mathrm{stimulus}}
            \end{equation}
            
            A safe interpretation is
            
            \begin{equation}
                L,\ t_{\mathrm{peak}},\ \ensuremath{\phi}
                \rightarrow
                \text{stimulus timing, latency, and phase-like timing relationships}.
            \end{equation}
        
        \ProjectDeepHeading{12}{Waveform Variability as Cellular-State Information}
            
            Waveform variability can also reflect internal cellular state. Spike waveform depends on voltage-gated channel density, channel inactivation, neuromodulation, recent firing history, and axonal or synaptic state. These waveform variations can change calcium entry and synaptic release, which means that spike waveform should not always be treated as biologically irrelevant \citep{zbili-2019-past-and-future-of-analog-digital-modulation-of-synaptic-transmission}.
            
            \begin{equation}
                \text{waveform shape}
                \rightarrow
                \text{channel state, membrane state, neuromodulation, firing history, and synaptic strength}.
            \end{equation}
        
        \ProjectDeepHeading{12}{Compact Conclusion}
            
            For graded potentials, the waveform can carry information through amplitude, polarity, slope, width, decay, timing, and adaptation. For classical action potentials, the main code is usually spike timing and spike rate, but waveform variation can still modulate synaptic transmission. The strongest directly relevant evidence for a true amplitude-based graded spike code is the study of cerebrospinal-fluid-contacting neurons, where different spike amplitudes signal different neurotransmitter inputs \citep{johnson-2023-graded-spikes-differentially-signal-neurotransmitter-input-in-cerebrospinal-fluid-contacting-neurons-of-the-mouse-spinal-cord}.

\ProjectDeepHeading{9}{Kind}

\ProjectDeepHeading{10}{Excitatory postsynaptic potential}

% Source: robot/part/body/part/nervous_system/part/shared/neuron/part/action_potential/kind/graded_or_binary/kind/graded/kind/excitatory_postsynaptic_potential/excitatory_postsynaptic_potential.tex

In neuroscience, an excitatory postsynaptic potential (EPSP) is a postsynaptic potential that makes the postsynaptic neuron more likely to fire an action potential \seeherefooter{https://en.wikipedia.org/wiki/Excitatory_postsynaptic_potential}.

\ProjectDeepHeading{10}{Inhibitory postsynaptic potential}

% Source: robot/part/body/part/nervous_system/part/shared/neuron/part/action_potential/kind/graded_or_binary/kind/graded/kind/inhibitory_postsynaptic_potential/inhibitory_postsynaptic_potential.tex

An inhibitory postsynaptic potential (IPSP) is a kind of synaptic potential that makes a postsynaptic neuron less likely to generate an action potential \seeherefooter{https://en.wikipedia.org/wiki/Inhibitory_postsynaptic_potential}.

\ProjectDeepHeading{5}{Neural computation}

% Source: robot/part/body/part/nervous_system/part/shared/neuron/part/action_potential/neural_computation/neural_computation.tex

Neural computation is the information processing performed by networks of neurons. Neural computation is affiliated with the philosophical tradition known as Computational theory of mind, also referred to as computationalism, which advances the thesis that neural computation explains cognition.
    \url{https://en.wikipedia.org/wiki/Neural_computation}

\ProjectDeepHeading{6}{Neuoral ensemble}

% Source: robot/part/body/part/nervous_system/part/shared/neuron/part/action_potential/neural_computation/neuoral_ensemble/neuoral_ensemble.tex

\url{https://en.wikipedia.org/wiki/Neuronal_ensemble}
    A neuronal ensemble is a population of nervous system cells (neurons) (or cultured neurons) involved in a particular neural computation.
    
    an alternative is \clickurlfooter{https://en.wikipedia.org/wiki/Grandmother_cell}  is a hypothetical neuron that represents a complex but specific concept or object.

% Source: robot/part/body/part/nervous_system/part/shared/neuron/part/action_potential/neural_computation/neuoral_ensemble/grandmother_cell.tex

is a hypothetical neuron that represents a complex but specific concept or object \seeherefooter{https://en.wikipedia.org/wiki/Grandmother_cell}

\ProjectDeepHeading{5}{Population}

% Source: robot/part/body/part/nervous_system/part/shared/neuron/part/action_potential/population/spike.tex

\clickurlfooter{https://en.wikipedia.org/wiki/Population_spike}
    
    \cite{mcnaughton-1978-synaptic-enhancement-in-fascia-dentata-cooperativity-among-coactive-afferents} studies synaptic enhancement in the fascia dentata of the hippocampus and emphasizes cooperativity among simultaneously active afferents.
    It is relevant to population spikes because that phenomenon is commonly observed in hippocampal field recordings where many nearby neurons fire in a coordinated way.
    Conceptually, the paper helps connect collective extracellular signals to underlying synaptic input and coordinated neuronal discharge.
    So, for your purpose, it is a useful early reference if you want a source tied to hippocampal population-level spiking rather than to a single neuron’s action potential

\ProjectDeepHeading{5}{Stages}

\ProjectDeepHeading{6}{Part}

% Source: robot/part/body/part/nervous_system/part/shared/neuron/part/action_potential/stages/part/depolarzation.tex

In biology, depolarization or hypopolarization is a change within a cell, during which the cell undergoes a shift in electric charge distribution, resulting in less negative charge inside the cell compared to the outside \seeherefooter{https://en.wikipedia.org/wiki/Depolarization}.
    
    An action potential is a rapid, all-or-none change in membrane voltage across an excitable cell membrane. Depolarization is mainly caused by sodium influx, while repolarization is mainly caused by potassium efflux. Ion gradients and the Na/K-ATPase pump help restore and maintain the resting membrane potential after excitation \cite{grider-2023-physiology-action-potential}.

\ProjectDeepHeading{5}{Summetion}

% Source: robot/part/body/part/nervous_system/part/shared/neuron/part/action_potential/summetion/summation.tex

\seeherefooter{https://en.wikipedia.org/wiki/Summation_(neurophysiology)}
    \seeherefooter{https://www.lecerveau.ca/flash/capsules/articles_pdf/temporal_summation.pdf} defines neural summation as Summation, which includes both spatial summation and temporal summation, is the process that determines whether or not an action potential will be generated by the combined effects of excitatory and inhibitory signals, both from multiple simultaneous inputs (spatial summation), and from repeated inputs (temporal summation). Depending on the sum total of many individual inputs, summation may or may not reach the threshold voltage to trigger an action potential.

\ProjectDeepHeading{6}{Kind}

% Source: robot/part/body/part/nervous_system/part/shared/neuron/part/action_potential/summetion/kind/spatial.tex

Spatial summation: several different inputs are summed at the same time.

\ProjectDeepHeading{5}{Wave form}

% Source: robot/part/body/part/nervous_system/part/shared/neuron/part/action_potential/wave_form/wave_form.tex

\seeherefooter{https://en.wikipedia.org/wiki/Spike_response_model}

\ProjectDeepHeading{6}{Kind}

\ProjectDeepHeading{7}{Fitzhugh nagumo}

% Source: robot/part/body/part/nervous_system/part/shared/neuron/part/action_potential/wave_form/kind/fitzhugh_nagumo/fitzhugh_nagumo.tex

The FitzHugh--Nagumo model (FHN)\cite{cebrian-lacasa-2025-six-decades-of-the-fitzhugh-nagumo-model-a-guide-through-its-spatio-temporal-dynamics-and-influence-across-disciplines} describes a prototype of an excitable system, such as a neuron.
    It is a simplified dynamical model of neuronal excitability and action-potential-like behavior.
    
    \begin{equation}
        \dot{v}
        =
        c
        \Bigl(
        v
        -
        \frac{v^{3}}{3}
        +
        w
        +
        I
        \Bigr)
    \end{equation}
    
    \begin{equation}
        \dot{w}
        =
        -
        \frac{1}{c}
        \Bigl(
        v
        -
        a
        +
        bw
        \Bigr)
    \end{equation}
    
    \cite{cebrian-lacasa-2025-six-decades-of-the-fitzhugh-nagumo-model-a-guide-through-its-spatio-temporal-dynamics-and-influence-across-disciplines}
    reviews six decades of research on the FitzHugh--Nagumo model as a simplified but powerful description of excitable dynamics.
    It shows that the model captures not only neuronal excitability and oscillations, but also bistability, traveling waves, synchronization, and spatial pattern formation.
    The paper emphasizes that the model has become useful beyond neuroscience, including cardiac dynamics, biology, electronics, optics, mathematics, and physics.
    Its main contribution is to organize the known dynamical regimes and bifurcations of the model, especially in spatially extended and coupled systems.

\ProjectDeepHeading{7}{Frequency modulated mobius}

% Source: robot/part/body/part/nervous_system/part/shared/neuron/part/action_potential/wave_form/kind/frequency_modulated_mobius/frequency_modulated_mobius.tex

\seeherefooter{https://journal.r-project.org/articles/RJ-2022-015/}
    
    \cite{rueda-2021-a-novel-wave-decomposition-for-oscillatory-signals}
    and
    \cite{rodriguez-collado-2021-identification-of-action-potential-waveforms}
    propose a wave decomposition method for oscillatory signals based on Frequency Modulated M\"obius (FMM) waves.
    The method represents complex oscillatory patterns as a sum of interpretable wave components with parameters controlling amplitude, phase, location, and shape.
    For action-potential-like curves, its main value is that it provides a citable mathematical basis for approximating signal waveforms by combining several FMM components.
    
    A single FMM wave can be written as follows:
    
    \begin{equation}
        X(t)
        =
        M
        +
        A
        \cos
        \Biggl[
            \ensuremath{\beta}
            +
            2
            \arctan
            \Biggl(
            \ensuremath{\omega}
            \tan
            \Bigl(
            \frac{t-\ensuremath{\alpha}}{2}
            \Bigr)
            \Biggr)
            \Biggr]
    \end{equation}
    
    \begin{itemize}
        \item $X(t)$ is the value of the FMM wave at time $t$.
        
        \item $M$ is the baseline parameter. Bound: $M \in \mathbb{R}$.
        It controls the vertical shift of the FMM signal without changing
        the waveform shape, phase, or temporal deformation.
        
        \item $t$ is time. Bound: $t \in [0, 2\ensuremath{\pi})$.
        
        \item $A$ is the amplitude parameter. Bound: $A \geq 0$.
        It controls the vertical scale of the waveform.
        
        \item $\ensuremath{\alpha}$ is the location parameter. Bound:
        $\ensuremath{\alpha} \in [0, 2\ensuremath{\pi})$. It controls the horizontal position of
        the waveform on the time axis.
        
        \item $\ensuremath{\beta}$ is the phase parameter. Bound:
        $\ensuremath{\beta} \in [0, 2\ensuremath{\pi})$. It controls the internal phase of the wave.
        Together with $\ensuremath{\alpha}$ and $\ensuremath{\omega}$, it determines the temporal
        position of the maximum and minimum points of the waveform.
        
        \item $\ensuremath{\omega}$ is the shape parameter. Bound:
        $\ensuremath{\omega} \in (0, 1]$. It controls the sharpness, skewness, and
        temporal deformation of the waveform. When $\ensuremath{\omega}$ is close to
        zero, the waveform becomes more spike-like; when $\ensuremath{\omega}$ is close
        to one, the waveform becomes more similar to a sinusoidal wave.
        
        \item $\cos(\cdot)$ is the cosine function used to generate the
        oscillatory waveform.
        
        \item $\arctan(\cdot)$ is the inverse tangent function used inside
        the phase transformation.
        
        \item $\tan(\cdot)$ is the tangent function used to transform the
        time variable before applying the inverse tangent.
        
        \item $\ensuremath{\beta} + 2 \arctan\bigl(\ensuremath{\omega}
        \tan\bigl(\frac{t-\ensuremath{\alpha}}{2}\bigr)\bigr)$ is the phase function of
        the FMM wave.
    \end{itemize}
    
    
    \ProjectDeepHeading{10}{Properties}
        \begin{itemize}
            \item The period of FMM wave is $$2\ensuremath{\pi}$$
        \end{itemize}
    
    
    \ProjectDeepHeading{10}{Implementations}
        For python implementation \seeherefooter{https://github.com/FMMGroupVa/fmmpy} and for R package \seeherefooter{https://cran.r-project.org/web/packages/FMM/vignettes/FMMVignette.html}

% Source: robot/part/body/part/nervous_system/part/shared/neuron/part/action_potential/wave_form/kind/frequency_modulated_mobius/multi.tex

\ProjectDeepHeading{11}{Multi-component FMM waveform}
    
    A complex waveform can be modeled as the sum of several FMM waves.
    The general $m$-component model is:
    
    \begin{equation}
        V(t)
        =
        M
        +
        \sum_{j=1}^{m}
        W_j(t)
    \end{equation}
    
    Equivalently:
    
    \begin{equation}
        V(t)
        =
        M
        +
        \sum_{j=1}^{m}
        A_j
        \cos
        \Biggl[
            \ensuremath{\beta}_j
            +
            2
            \arctan
            \Biggl(
            \ensuremath{\omega}_j
            \tan
            \Bigl(
            \frac{t-\ensuremath{\alpha}_j}{2}
            \Bigr)
            \Biggr)
            \Biggr]
    \end{equation}
    
    Here, $M$ is the baseline or vertical offset of the full signal.
    In the case of an action potential, $V(t)$ represents the membrane voltage at time $t$.
    \cite{rueda-2021-a-novel-wave-decomposition-for-oscillatory-signals}

% Source: robot/part/body/part/nervous_system/part/shared/neuron/part/action_potential/wave_form/kind/frequency_modulated_mobius/superposition.tex

\ProjectDeepHeading{10}{A sample equation we are trying to solve}
        \[
            \begin{aligned}
                &M
                +
                A\cos\left(
                         \ensuremath{\beta}
                         +
                         2\arctan\left(
                                     \ensuremath{\omega}\tan\left(\frac{t-\ensuremath{\alpha}}{2}\right)
                    \right)
                \right)
                \\
                &=
                \\
                &M_1
                +
                A_1\cos\left(
                           \ensuremath{\beta}_1
                    +
                    2\arctan\left(
                                \ensuremath{\omega}_1\tan\left(\frac{t-\ensuremath{\alpha}_1}{2}\right)
                    \right)
                \right)
                \\
                &+
                M_2
                +
                A_2\cos\left(
                           \ensuremath{\beta}_2
                    +
                    2\arctan\left(
                                \ensuremath{\omega}_2\tan\left(\frac{t-\ensuremath{\alpha}_2}{2}\right)
                    \right)
                \right)
            \end{aligned}
        \]
    
    
    \ProjectDeepHeading{10}{Single FMM Approximation of a Two-Component FMM Signal}
        
        The goal is to approximate a signal generated by the sum of two FMM components using only one FMM signal. The original signal is assumed to have the following form:
        
        \[
            y(t_i)
            =
            M_0
            +
            f(t_i; A_1,\ensuremath{\alpha}_1,\ensuremath{\beta}_1,\ensuremath{\omega}_1)
            +
            f(t_i; A_2,\ensuremath{\alpha}_2,\ensuremath{\beta}_2,\ensuremath{\omega}_2)
        \]
        
        where each FMM component is defined as
        
        \[
            f(t; A,\ensuremath{\alpha},\ensuremath{\beta},\ensuremath{\omega})
            =
            A
            \cos
            \left(
                \ensuremath{\beta}
                +
                2
                \arctan
                \left(
                    \ensuremath{\omega}
                    \tan
                    \left(
                        \frac{t-\ensuremath{\alpha}}{2}
                    \right)
                \right)
            \right)
        \]
        
        The approximation replaces the two-component signal with a single five-parameter FMM signal:
        
        \[
            \hat{y}(t_i; \ensuremath{\theta})
            =
            M
            +
            f(t_i; A,\ensuremath{\alpha},\ensuremath{\beta},\ensuremath{\omega})
        \]
        
        where the unknown parameter vector is
        
        \[
            \ensuremath{\theta}
            =
            \left[
                M,
                A,
                \ensuremath{\alpha},
                \ensuremath{\beta},
                \ensuremath{\omega}
            \right]^{T}
        \]
        
        The fitting problem is formulated as a nonlinear least-squares problem. For each sampled time point \(t_i\), the residual is the difference between the original two-component signal and the fitted single-component signal:
        
        \[
            r_i(\ensuremath{\theta})
            =
            y(t_i)
            -
            \hat{y}(t_i; \ensuremath{\theta})
        \]
        
        The objective function is the sum of squared residuals:
        
        \[
            J(\ensuremath{\theta})
            =
            \sum_{i=1}^{N}
            r_i(\ensuremath{\theta})^2
            =
            \sum_{i=1}^{N}
            \left(
                y(t_i)
                -
                \hat{y}(t_i; \ensuremath{\theta})
            \right)^2
        \]
        
        The optimal parameter vector is obtained by minimizing this objective function:
        
        \[
            \ensuremath{\theta}^{*}
            =
            \arg\min_{\ensuremath{\theta}}
            J(\ensuremath{\theta})
        \]
    
    
    \ProjectDeepHeading{10}{Pseudo-Code}
        
        \begin{algorithm}[htbp]
		\caption{Single FMM Approximation of a Two-Component FMM Signal}
        \begin{algorithmic}[1]
            \Require Time points \(T = \{t_1,t_2,\dots,t_N\}\)
            \Require Target signal values \(Y = \{y(t_1),y(t_2),\dots,y(t_N)\}\)
            \Require Parameter bounds for \(M\), \(A\), \(\ensuremath{\alpha}\), \(\ensuremath{\beta}\), and \(\ensuremath{\omega}\)
            \Ensure Best single FMM parameter vector \(\ensuremath{\theta}^{*}\)
            
            \State Define the single FMM model:
            \[
                \hat{y}(t_i; \ensuremath{\theta})
                =
                M
                +
                A
                \cos
                \left(
                    \ensuremath{\beta}
                    +
                    2
                    \arctan
                    \left(
                        \ensuremath{\omega}
                        \tan
                        \left(
                            \frac{t_i-\ensuremath{\alpha}}{2}
                        \right)
                    \right)
                \right)
            \]
            
            \State Define the residual vector:
            \[
                R(\ensuremath{\theta})
                =
                \left[
                    y(t_1)-\hat{y}(t_1;\ensuremath{\theta}),
                y(t_2)-\hat{y}(t_2;\ensuremath{\theta}),
                    \dots,
                y(t_N)-\hat{y}(t_N;\ensuremath{\theta})
                \right]^{T}
            \]
            
            \State Define the scalar error function:
            \[
                J(\ensuremath{\theta})
                =
                R(\ensuremath{\theta})^{T}R(\ensuremath{\theta})
            \]
            
            \State Use a global optimization method to find an initial parameter vector:
            \[
                \ensuremath{\theta}_{\mathrm{global}}
                =
                \arg\min_{\ensuremath{\theta}}
                J(\ensuremath{\theta})
            \]
            
            \State Use bounded nonlinear least-squares optimization starting from \(\ensuremath{\theta}_{\mathrm{global}}\):
            \[
                \ensuremath{\theta}^{*}
                =
                \arg\min_{\ensuremath{\theta}}
                \|R(\ensuremath{\theta})\|_2^2
            \]
            
            \State Compute the fitted signal:
            \[
                \hat{Y}
                =
                \left[
                    \hat{y}(t_1;\ensuremath{\theta}^{*}),
                    \hat{y}(t_2;\ensuremath{\theta}^{*}),
                    \dots,
                    \hat{y}(t_N;\ensuremath{\theta}^{*})
                \right]
            \]
            
            \State \Return \(\ensuremath{\theta}^{*}\), \(\hat{Y}\), and \(J(\ensuremath{\theta}^{*})\)
        \end{algorithmic}
        \end{algorithm}
    
    
    \ProjectDeepHeading{10}{Explanation of the Approximation}
        
        The original signal contains two FMM components. In general, the sum of two FMM waves is not exactly equal to a single FMM wave. Therefore, the method does not find an exact analytical replacement. Instead, it finds the best numerical approximation under a least-squares criterion.
        
        The single FMM model has only five parameters: baseline, amplitude, alpha, beta, and omega. These parameters are adjusted so that the generated single FMM curve becomes as close as possible to the target signal over all sampled time points.
        
        The approximation works by comparing the target value and the fitted value at every time point. The difference between them is called the residual. If the fitted curve is close to the target curve, the residuals are small. If the fitted curve is far from the target curve, the residuals are large.
        
        The optimization method minimizes the sum of squared residuals. Squaring the residuals has two effects. First, positive and negative errors do not cancel each other. Second, larger errors receive a stronger penalty than smaller errors.
        
        The fitting procedure has two stages. The first stage uses global optimization to search broadly through the parameter space and find a reasonable starting point. This reduces the chance of starting from a poor local solution. The second stage uses nonlinear least-squares optimization to refine the parameters locally and reduce the residuals more precisely.
        
        The final result is a single FMM wave whose shape, phase, amplitude, and baseline are chosen so that it approximates the two-component signal with minimum squared error over the sampled interval.


%        
%

\ProjectDeepHeading{7}{Hodgkin}

% Source: robot/part/body/part/nervous_system/part/shared/neuron/part/action_potential/wave_form/kind/hodgkin/hodgkin.tex

\cite{hodgkin-1952-a-quantitative-description-of-membrane-current-and-its-application-to-conduction-and-excitation-in-nerve} present a quantitative model of nerve membrane current based on sodium, potassium, leakage current, and membrane capacitance.
    The paper formulates voltage-dependent conductance equations that explain action potential generation, propagation, excitation, and refractory behavior.
    Its main contribution is the mathematical foundation of the Hodgkin-Huxley model, connecting ionic membrane dynamics to nerve conduction.

\ProjectDeepHeading{7}{Morris lecar model}

% Source: robot/part/body/part/nervous_system/part/shared/neuron/part/action_potential/wave_form/kind/morris_lecar_model/morris_lecar_model.tex

The Morris–Lecar model is a biological neuron model developed by professors Catherine Morris and Harold Lecar to reproduce the variety of oscillatory behavior in relation to Ca++ and K+ conductance in the muscle fiber of the giant barnacle \seeherefooter{https://en.wikipedia.org/wiki/Morris\%E2\%80\%93Lecar_model}

\ProjectDeepHeading{4}{Axon}

% Source: robot/part/body/part/nervous_system/part/shared/neuron/part/axon/axon.tex

axons are the outputs

\ProjectDeepHeading{4}{Dendrite}

% Source: robot/part/body/part/nervous_system/part/shared/neuron/part/dendrite/dendrite.tex

The input of the neuron

\ProjectDeepHeading{4}{Noise}

% Source: robot/part/body/part/nervous_system/part/shared/neuron/part/noise/noise.tex

\clickurlfooter{https://en.wikipedia.org/wiki/Neuronal_noise} defines neunoral the random intrinsic electrical fluctuations within neuronal networks. These fluctuations are not associated with encoding a response to internal or external stimuli and can be from one to two orders of magnitude

\ProjectDeepHeading{4}{Soma}

% Source: robot/part/body/part/nervous_system/part/shared/neuron/part/soma/soma.tex

Soma or cell body, is the bulbous, non-process portion of a neuron or glial cell that contains the cell nucleus \seeherefooter{https://en.wikipedia.org/wiki/Soma_(biology)}.

\ProjectDeepHeading{3}{Pathway}

% Source: robot/part/body/part/nervous_system/part/shared/neuron/pathway/pathway.tex

In neuroanatomy, a neural pathway is the connection formed by axons that project from neurons to make synapses onto neurons in another location, to enable neurotransmission (the sending of a signal from one region of the nervous system to another)   \url{https://en.wikipedia.org/wiki/Neural_pathway}

\ProjectDeepHeading{3}{Population}

\ProjectDeepHeading{4}{Kind}

\ProjectDeepHeading{5}{Network}

% Source: robot/part/body/part/nervous_system/part/shared/neuron/population/kind/network/network.tex

A neural network, also called a neuronal network, is an interconnected population of neurons (typically containing multiple neural circuits) \url{https://en.wikipedia.org/wiki/Neural_network_(biology)}
    % 
    %

\ProjectDeepHeading{6}{Shared}

\ProjectDeepHeading{7}{Circuit}

% Source: robot/part/body/part/nervous_system/part/shared/neuron/population/kind/network/shared/circuit/circuit.tex

A neural circuit is a population of neurons interconnected by synapses to carry out a specific function when activated. Multiple neural circuits interconnect with one another to form large scale brain networks. \seeherefooter{https://en.wikipedia.org/wiki/Neural_circuit}.
    %
    %

\ProjectDeepHeading{8}{Kind}

% Source: robot/part/body/part/nervous_system/part/shared/neuron/population/kind/network/shared/circuit/kind/converging.tex

\begin{itemize}
        \item Usually happens in from prepheral neurons toward the brain
    \end{itemize}
    
    \ProjectDeepHeading{11}{Important brains}
    \begin{itemize}
        \item On the neuronal basis for multisensory convergence: a brief overview \seeherefooter{https://www.sciencedirect.com/science/article/abs/pii/S0926641002000599}
    \end{itemize}

\ProjectDeepHeading{5}{Pool}

% Source: robot/part/body/part/nervous_system/part/shared/neuron/population/kind/pool/pool.tex

A neural pool is a group of functionally related neurons that work together to process or represent specific information. It is different from neural circuit in the sense that a circuit discusses only a set of interconnected neurons while a pool is focused on a set of neurons with a shared or common functionality
    
    
    \ProjectDeepHeading{8}{Some sources}
        For Transmission and Processing of Signals in Neuronal Pools see
        \seeherefooter{https://www.brainkart.com/article/Transmission-and-Processing-of-Signals-in-Neuronal-Pools_19636/}
        A video on neuronal pools and neural processing
        \seeherefooter{https://www.youtube.com/watch?v=QJ8AW5pi2T4}

\ProjectDeepHeading{4}{Model}

\ProjectDeepHeading{5}{Kind}

% Source: robot/part/body/part/nervous_system/part/shared/neuron/population/model/kind/efficient_coding_hypothesis.tex

The efficient coding hypothesis was proposed by Horace Barlow in 1961 as a theoretical model of sensory neuroscience in the brain.[1] Within the brain, neurons communicate with one another by sending electrical impulses referred to as action potentials or spikes.
    
    Barlow hypothesized that the spikes in the sensory system formed a neural code for efficiently representing sensory information. By efficient it is understood that the code minimized the number of spikes needed to transmit a given signal. This is somewhat analogous to transmitting information across the internet, where different file formats can be used to transmit a given image. Different file formats require different numbers of bits for representing the same image at a given distortion level, and some are better suited for representing certain classes of images than others. According to this model, the brain is thought to use a code which is suited for representing visual and audio information which is representative of an organism's natural environment \seeherefooter{https://en.wikipedia.org/wiki/Efficient_coding_hypothesis}.
%    
%

\ProjectDeepHeading{3}{Saltatory}

% Source: robot/part/body/part/nervous_system/part/shared/neuron/saltatory/saltatory.tex

In neuroscience, saltatory conduction is the propagation of action potentials along myelinated axons from one node of Ranvier to the next, increasing the conduction velocity of action potentials \seeherefooter{https://en.wikipedia.org/wiki/Saltatory_conduction}.

\ProjectDeepHeading{3}{Transmission}

% Source: robot/part/body/part/nervous_system/part/shared/neuron/transmission/transmission.tex

Neurotransmission (Latin: transmissio "passage, crossing" from transmittere "send, let through") is the process by which signaling molecules called neurotransmitters are released by the axon terminal of a neuron (the presynaptic neuron), and bind to and react with the receptors on the dendrites of another neuron (the postsynaptic neuron) a short distance away \seeherefooter{https://en.wikipedia.org/wiki/Neurotransmission}.

\ProjectDeepHeading{2}{Neurotransmission}

% Source: robot/part/body/part/nervous_system/part/shared/neurotransmission/neurotransmission.tex

Neurotransmission (Latin: transmissio "passage, crossing" from transmittere "send, let through") is the process by which signaling molecules called neurotransmitters are released by the axon terminal of a neuron (the presynaptic neuron), and bind to and react with the receptors on the dendrites of another neuron (the postsynaptic neuron) a short distance away \seeherefooter{https://en.wikipedia.org/wiki/Neurotransmission}.

\ProjectDeepHeading{3}{Neurotransmitter}

% Source: robot/part/body/part/nervous_system/part/shared/neurotransmission/neurotransmitter/neurotransmitter.tex

A neurotransmitter is a signaling molecule secreted by a neuron to affect another cell across a synapse \seeherefooter{https://en.wikipedia.org/wiki/Neurotransmitter}.

\ProjectDeepHeading{2}{Synapse}

% Source: robot/part/body/part/nervous_system/part/shared/synapse/synapse.tex

a synapse is a structure that allows a neuron (or nerve cell) to pass an electrical or chemical signal to another neuron or a target effector cell \url{https://en.wikipedia.org/wiki/Synapse}.
    % 
    %

\ProjectDeepHeading{3}{Part}

\ProjectDeepHeading{4}{Chemical}

% Source: robot/part/body/part/nervous_system/part/shared/synapse/part/chemical/chemical.tex

Chemical synapses are biological junctions through which neurons' signals can be sent to each other and to non-neuronal cells such as those in muscles or glands \seeherefooter{https://en.wikipedia.org/wiki/Chemical_synapse}.

\ProjectDeepHeading{4}{Potential}

% Source: robot/part/body/part/nervous_system/part/shared/synapse/part/potential/potential.tex

Synaptic potential refers to the potential difference across the postsynaptic membrane that results from the action of neurotransmitters at a neuronal synapse \seeherefooter{https://en.wikipedia.org/wiki/Synaptic_potential}.

\ProjectDeepHeading{5}{Postsynaptic potential}

% Source: robot/part/body/part/nervous_system/part/shared/synapse/part/potential/postsynaptic_potential/postsynaptic_potential.tex

Postsynaptic potentials are changes in the membrane potential of the postsynaptic terminal of a chemical synapse.
    \seeherefooter{https://en.wikipedia.org/wiki/Postsynaptic_potential}

\ProjectDeepHeading{6}{Kind}

% Source: robot/part/body/part/nervous_system/part/shared/synapse/part/potential/postsynaptic_potential/kind/kind.tex

An excitatory postsynaptic potential (EPSP) is a postsynaptic potential that makes the postsynaptic neuron more likely to fire an action potential \footnote{}

\ProjectDeepHeading{4}{Transmission}

% Source: robot/part/body/part/nervous_system/part/shared/synapse/part/transmission/transmission.tex

\ProjectDeepHeading{7}{Resources}
        \seeherefooter{https://www.youtube.com/watch?v=fhiJo1QYGtA}

\subparagraph{Platform}

\ProjectDeepHeading{1}{Flight controller}

% Source: robot/part/body/part/nervous_system/platform/flight_controller/flight_controller.tex

SpeedyBee

\paragraph{Organ}

% Source: robot/part/body/part/organ/organ.tex

An organ is a distinct part of a living organism that performs a specific function or a set of related functions.


    It is made of different types of tissues that work together in a coordinated way.
    
    A structurally identifiable body part composed of multiple tissues that cooperate to carry out a particular physiological function.
    
    Heart → pumps blood

    Lungs → exchange oxygen and carbon dioxide
    
    Liver → processes nutrients and detoxifies chemicals
    
    Kidneys → filter blood and produce urine
    
    Brain → processes and integrates information
    
    Organ = group of different tissues working together

\subparagraph{Tissue}

% Source: robot/part/body/part/organ/tissue/tissue.tex

Tissue = group of similar cells with a common function

\paragraph{Shared}

\subparagraph{Cell}

\ProjectDeepHeading{1}{Part}

\ProjectDeepHeading{2}{Membrane}

% Source: robot/part/body/part/shared/cell/part/membrane/membrane.tex

The cell membrane separates and protects the interior of a cell from the outside environment
    \seeherefooter{https://en.wikipedia.org/wiki/Cell_membrane}

% Source: robot/part/body/part/shared/cell/part/membrane/potential.tex

Membrane potential (also transmembrane potential or membrane voltage) is the difference in electric potential between the interior and the exterior of a biological cell \seeherefooter{https://en.wikipedia.org/wiki/Membrane_potential}
    
    \begin{itemize}
        \item usually the within the cell the voltage is -70 and the voltage should reach to 050 so that the neuron fires
    \end{itemize}

\ProjectDeepHeading{2}{Signaling}

% Source: robot/part/body/part/shared/cell/part/signaling/signaling.tex

Cell signaling is the biological process by which a cell interacts with itself, with other cells, and with the environment. Cell signaling is a fundamental property of all forms of life. . Typically, the signaling process involves three components: the first messenger (the ligand), the receptor, and the signal itself \seeherefooter{https://en.wikipedia.org/wiki/Cell_signaling}.
    
    \begin{itemize}
        \item Neuron signaling at the end of the axon is an example
        \item Insoline release is another example of cell signalling
    \end{itemize}
    
    \textbf{Other resources}
    \begin{itemize}
        \item Cell to Cell Communication || Types of signaling \seeherefooter{https://www.youtube.com/watch?v=i3bY-JCYs4A}
        \item intracell communication \seeherefooter{https://www.youtube.com/watch?v=aIZQ3ker0KE}
    \end{itemize}

% Source: robot/part/body/part/shared/cell/part/signaling/transduction.tex

Signal transduction is the process by which a chemical or physical signal is transmitted through a cell as a series of molecular events \seeherefooter{https://en.wikipedia.org/wiki/Signal_transduction}

\subparagraph{Genetic}

\ProjectDeepHeading{1}{Dna}

% Source: robot/part/body/part/shared/genetic/dna/dna.tex

DNA is the molucule that carries the genes

\subsubsection{Process}

% Source: robot/part/body/process/process.tex

A biological process is a processes necessary for an organism to live, and therefore shapes its capacity to interact with its environment. Biological processes are made of many chemical reactions or other events that are involved in the persistence and transformation of life forms
    \seeherefooter{https://en.wikipedia.org/wiki/Biological_process}
    %
    %

% Source: robot/part/mind/mental_image.tex

a mental image is an experience that, on most occasions, significantly resembles the experience of "perceiving" some object, event, or scene but occurs when the relevant object, event, or scene is not actually present to the senses. \seehere{https://en.wikipedia.org/wiki/Mental_image}
    % 
    %

% Source: robot/part/mind/mental_representation.tex

A mental representation (or cognitive representation), in philosophy of mind, cognitive psychology, neuroscience, and cognitive science, is a hypothetical internal cognitive symbol that represents external reality or its abstractions \seehere{https://en.wikipedia.org/wiki/Mental_representation}
    % 
    %

\subsubsection{Affordance}

% Source: robot/part/mind/affordance/affordance.tex

\url{https://en.wikipedia.org/wiki/Affordance\#Affordance_in_robotics}
    
    A review
    \url{https://www.tandfonline.com/doi/full/10.1080/01691864.2017.1394912}
    % 
    %

% Source: robot/part/mind/cognition/cognition.tex

\cite{revlin-2012-cognition-theory-and-practice} defines cognition as mental processes that deal with knowledge, involving the acquisition, transformation, storage, retrieval, and use of information.
    %
    %

% Source: robot/part/mind/cognition/computational.tex

\url{https://en.wikipedia.org/wiki/Computational_cognition}
    
    %
    %

% Source: robot/part/mind/cognition/neuroscience.tex

\url{https://en.wikipedia.org/wiki/Cognitive_neuroscience}
    
    %
    %

\paragraph{Architecture}

% Source: robot/part/mind/cognition/architecture/architecture.tex

A cognitive architecture is both a theory about the structure of the human mind and a computational instantiation of such a theory used in the fields of artificial intelligence (AI) and computational cognitive science. \url{https://en.wikipedia.org/wiki/Cognitive_architecture}

\subparagraph{Kinds}

\ProjectDeepHeading{1}{Act r}

% Source: robot/part/mind/cognition/architecture/kinds/act_r/act_r.tex

"Adaptive Control of Thought—Rational") is a cognitive architecture mainly developed by John Robert Anderson and Christian Lebiere at Carnegie Mellon University. Like any cognitive architecture, ACT-R aims to define the basic and irreducible cognitive and perceptual operations that enable the human mind. In theory, each task that humans can perform should consist of a series of these discrete operations. \seehere{https://en.wikipedia.org/wiki/ACT-R}
% 
%

\ProjectDeepHeading{1}{Dual}

% Source: robot/part/mind/cognition/architecture/kinds/dual/dual.tex

DUAL is a general cognitive architecture integrating the connectionist and symbolic approaches at the micro level. \seehere{https://en.wikipedia.org/wiki/DUAL_(cognitive_architecture)}
    % 
    %

\ProjectDeepHeading{1}{Soar}

% Source: robot/part/mind/cognition/architecture/kinds/soar/soar.tex

Soar[1] is a cognitive architecture,[2] originally created by John Laird, Allen Newell, and Paul Rosenbloom at Carnegie Mellon University.
    
    The goal of the Soar project is to develop the fixed computational building blocks necessary for general intelligent agents – agents that can perform a wide range of tasks and encode, use, and learn all types of knowledge to realize the full range of cognitive capabilities found in humans, such as decision making, problem solving, planning, and natural-language understanding. \url{}
    % 
    %

\paragraph{Bias}

% Source: robot/part/mind/cognition/bias/bias.tex

\clickurlfooter{https://en.wikipedia.org/wiki/Cognitive_bias} defines cognitive bias as a systematic pattern of deviation from norm or rationality in judgment.
    
    \begin{figure}[H]
        \centering
        \ProjectSafeIncludeGraphics[width=0.9\textwidth]{robot/part/mind/cognition/bias/cognitive_bias_codex.png}
    \end{figure}
    %
    %

\subparagraph{Kinds}

% Source: robot/part/mind/cognition/bias/kinds/salience.tex

\clickurlfooter{https://en.wikipedia.org/wiki/Salience_(neuroscience)} defines slaience as the property by which something stands out. Salient events are an attentional mechanism by which organisms learn and survive; those organisms can focus their limited perceptual and cognitive resources on the pertinent (that is, salient) subset of the sensory data available to them.
    %
    %

\paragraph{Map}

% Source: robot/part/mind/cognition/map/map.tex

A cognitive map is a type of mental representation used by an individual to order their personal store of information about their everyday or metaphorical spatial environment, and the relationship of its component parts. \url{https://en.wikipedia.org/wiki/Cognitive_map}
    % 
    %

\paragraph{Module}

% Source: robot/part/mind/cognition/module/module.tex

A cognitive module in cognitive psychology is a specialized tool or sub-unit that can be used by other parts to resolve cognitive tasks. \seehere{https://en.wikipedia.org/wiki/Cognitive_module}
    % 
    %

\subparagraph{Meta cognition}

% Source: robot/part/mind/cognition/part/meta_cognition/meta_cognition.tex

\cite{metcalfe-1994-metacognition-knowing-about-knowing} defines metacognition is an awareness of one's thought processes and an understanding of the patterns behind them.
    \url{https://en.wikipedia.org/wiki/Metacognition}
    
    Also this book for a neuro science perspective
    \url{https://link.springer.com/book/10.1007/978-3-642-45190-4}
    \\
    \url{https://www.amazon.de/-/en/Cognitive-Neuroscience-Metacognition-Stephen-Fleming/dp/3662511711}
    
    \cite{fleur-2021-metacognition-ideas-and-insights-from-neuro-and-educational-sciences}
    This review compares how educational sciences and cognitive neuroscience study metacognition. It distinguishes metacognitive knowledge, meaning awareness and monitoring of one’s cognition, from metacognitive control, meaning regulation, planning, and adaptation of behaviour. The paper argues that future research should connect lab-based neural measures with educational protocols, especially for training offline metacognition such as reflection and learning strategies.
    
    \ProjectDeepHeading{6}{Image}
        \ProjectSafeIncludeGraphics[scale=1]{robot/part/mind/cognition/part/meta_cognition/meta.png}
    
    
    \ProjectDeepHeading{3}{Metacognitive Knowledge}
    \begin{itemize}
        \item \textbf{Definition:} Knowledge that a system has about its own cognitive processes, such as what it knows, what it does not know, and how it learns.
        \item \textbf{Typical components:}
        \begin{itemize}
            \item Knowledge of knowledge states
            \item Knowledge of strategies
            \item Knowledge of task difficulty
        \end{itemize}
    \end{itemize}
    
    
    \ProjectDeepHeading{3}{Metacognitive Monitoring}
    \begin{itemize}
        \item \textbf{Definition:} The process by which a system observes and evaluates its ongoing cognitive activities.
        \item \textbf{Typical functions:}
        \begin{itemize}
            \item Error detection
            \item Confidence estimation
            \item Uncertainty estimation
        \end{itemize}
    \end{itemize}
    
    
    \ProjectDeepHeading{3}{Metacognitive Control}
    \begin{itemize}
        \item \textbf{Definition:} The ability of a system to regulate or adjust its cognitive processes based on monitoring results.
        \item \textbf{Typical actions:}
        \begin{itemize}
            \item Changing strategies
            \item Allocating attention
            \item Requesting more information
        \end{itemize}
    \end{itemize}
    
    
    \ProjectDeepHeading{3}{Self-Awareness}
    \begin{itemize}
        \item \textbf{Definition:} The ability of a system to represent and recognize its own internal states and identity.
        \item \textbf{Typical aspects:}
        \begin{itemize}
            \item Self-model
            \item Body awareness
            \item Internal state awareness
        \end{itemize}
    \end{itemize}
    
    
    \ProjectDeepHeading{3}{Meta-Learning}
    \begin{itemize}
        \item \textbf{Definition:} The ability of a system to improve its own learning process through experience.
        \item \textbf{Typical aspects:}
        \begin{itemize}
            \item Learning to learn
            \item Adaptation of learning strategies
            \item Adjustment of learning parameters
        \end{itemize}
    \end{itemize}
    
    
    \ProjectDeepHeading{3}{Introspection}
    \begin{itemize}
        \item \textbf{Definition:} The capacity of a system to internally examine its own reasoning or cognitive states.
        \item \textbf{Typical aspects:}
        \begin{itemize}
            \item Examination of reasoning steps
            \item Evaluation of internal beliefs
            \item Reflection on cognitive states
        \end{itemize}
    \end{itemize}
    
    
    \ProjectDeepHeading{3}{Uncertainty Awareness}
    \begin{itemize}
        \item \textbf{Definition:} The ability of a system to estimate the reliability or uncertainty of its own predictions or decisions.
        \item \textbf{Typical aspects:}
        \begin{itemize}
            \item Confidence estimation
            \item Uncertainty estimation
            \item Reliability assessment of predictions
        \end{itemize}
    \end{itemize}
    
    \ProjectDeepHeading{6}{Examples}
        I know that I don't know

\ProjectDeepHeading{1}{Compuational}

% Source: robot/part/mind/cognition/part/meta_cognition/compuational/compuational.tex

\url{https://arxiv.org/pdf/2201.12885}
    % 
    %

\ProjectDeepHeading{1}{Parts}

\ProjectDeepHeading{2}{Self conciouness}

% Source: robot/part/mind/cognition/part/meta_cognition/parts/self_conciouness/self_conciouness.tex

\url{https://en.wikipedia.org/wiki/Self-consciousness}
    % 
    % 



    
    \cite{campbell-1995-the-body-image-and-self-consciousness} argues that body image contributes to self-consciousness by allowing a person to understand themselves as a concrete bodily object rather than merely as a subject of psychological states. He distinguishes body schema from self-consciousness: an animal or child may use bodily representations for posture and action without yet possessing full self-consciousness. The chapter connects bodily self-representation to causal understanding, arguing that grasping one’s own bodily and psychological states as internally connected is central to the conception of oneself.

\ProjectDeepHeading{3}{Self awareness}

% Source: robot/part/mind/cognition/part/meta_cognition/parts/self_conciouness/self_awareness/self_awareness.tex

\url{https://selfawarepatterns.com/2020/01/25/recurrent-processing-theory-and-the-function-of-consciousness/}
    
    see active inference on markov blanket that helps with drawing a boundry between self and environment
    
    In the philosophy of self, self-awareness is the awareness and reflection of one's own personality or individuality, including traits, feelings, and behaviors. \url{https://en.wikipedia.org/wiki/Self-awareness}
    
    \ProjectDeepHeading{6}{From the neuroscience point of view}
    Modern neuroscience treats self-awareness not as the product of a single ``center`` but as the emergent behavior of interacting brain systems. Functional MRI, lesion, and connectivity studies implicate a distributed network—particularly the medial prefrontal cortex, anterior cingulate cortex, posterior cingulate cortex, and temporoparietal junction—in processes of self-reflection and self-modelling. These regions overlap substantially with the brain's default mode network and engage in metacognitive monitoring—feedback loops in which predictions about internal and external states are continuously compared with incoming sensory and emotional input. Such neural machinery underlies the experience of being aware that one is aware.\cite{fleming-frith-2014-metacognition}
    
    \ProjectDeepHeading{6}{Corrolation predictive brain on sensory data}
    Experimental evidence suggests that self-awareness depends on the brain's capacity for metacognition—the monitoring and evaluation of its own processes. Such ``monitoring systems`` continually compare predicted sensory and emotional states with actual input, generating the experience of being aware of awareness itself. This feedback architecture allows the brain to notice discrepancies between expectation and perception, forming the foundation of conscious self-reflection. This recursive feedback process gives rise to the sensation of being aware of awareness, sometimes described as a ``mirror of mirrors`` within consciousness.\cite{garcia-2013-stuck-in-the-moment}


    
    It is the ability of the agent of itself and its distinction from the world
    
    compare with self conciouness
    \begin{minted}{python}
        class Awareness:
            def __init__(self):
                pass
            
            def discovered_actuators(self) -> list:
                pass
            
            def discovered_sensors(self) -> list:
                pass
            
            def what_am_i_doing_now(self) -> str:
                pass
    \end{minted}

% Source: robot/part/mind/cognition/part/meta_cognition/parts/self_conciouness/self_awareness/concept.tex

\url{https://en.wikipedia.org/wiki/Self-concept}
    % 
    %

% Source: robot/part/mind/cognition/part/meta_cognition/parts/self_conciouness/self_awareness/gpt_selfawarness_kinds_and_levels.tex

\ProjectDeepHeading{6}{Self-Awareness}
        
        Self-awareness generally refers to a system's capacity to represent, monitor, and reason about its own internal states, processes, and actions. In cognitive science and artificial intelligence, it is typically considered a higher-order cognitive capability emerging from several underlying components and organized into different forms and levels.
    
    
    \ProjectDeepHeading{6}{Components of Self-Awareness}
        
        Self-awareness is composed of several interacting cognitive mechanisms:
        
        \begin{enumerate}[label=\arabic*.]
\item \textbf{Self-representation}

The ability of an agent to maintain an internal model of itself. This includes representations of its body, capabilities, goals, and internal states.

\item \textbf{Self-monitoring}

Continuous observation of internal processes such as perceptions, actions, decisions, and internal states.

\item \textbf{Self-evaluation}

The ability to assess one's own performance, actions, or internal states against goals, expectations, or norms.

\item \textbf{Self-regulation}

The ability to modify behavior or internal processes based on self-evaluation.

\item \textbf{Meta-representation}

Representation of one's own cognitive processes (thinking about thinking).

\item \textbf{Temporal self-continuity}

Awareness that the self persists across time (past, present, and anticipated future states).
        \end{enumerate}
        
        These mechanisms together allow an agent to treat itself as an object of cognition.
    
    
    \ProjectDeepHeading{6}{Kinds of Self-Awareness}
        
        Different types of self-awareness have been identified in psychology, philosophy, and cognitive science.
        
        \ProjectDeepHeading{7}{Bodily Self-Awareness}
            
            Awareness of one's own physical body and its position in space.
            
            Examples:
            \begin{itemize}
                \item proprioception
                \item body ownership
                \item sense of agency over movements
            \end{itemize}
        
        \ProjectDeepHeading{7}{Perceptual Self-Awareness}
            
            Awareness that one's perceptions belong to oneself.
            
            Example:
            
            ``I am the one seeing this object.''
        
        \ProjectDeepHeading{7}{Reflective Self-Awareness}
            
            The ability to reflect on one's own mental states.
            
            Examples:
            
            \begin{itemize}
                \item awareness of beliefs
                \item awareness of emotions
                \item awareness of intentions
            \end{itemize}
        
        \ProjectDeepHeading{7}{Narrative Self-Awareness}
            
            Awareness of oneself as a temporally extended entity with a personal history.
            
            This includes autobiographical memory and personal identity.
        
        \ProjectDeepHeading{7}{Social Self-Awareness}
            
            Awareness of how one appears to others and how one is evaluated socially.
            
            Examples:
            
            \begin{itemize}
                \item embarrassment
                \item reputation management
                \item theory of mind about how others perceive the self
            \end{itemize}
    
    
    \ProjectDeepHeading{6}{Levels of Self-Awareness}
        
        Self-awareness can also be organized into hierarchical levels of cognitive complexity.
        
        \ProjectDeepHeading{7}{Level 0: No Self-Model}
            
            Systems react to stimuli but have no representation of themselves.
            
            Examples:
            
            \begin{itemize}
                \item simple reflex systems
                \item basic reactive robots
            \end{itemize}
        
        \ProjectDeepHeading{7}{Level 1: Bodily Awareness}
            
            The agent represents its body and physical state.
            
            Examples:
            
            \begin{itemize}
                \item proprioception
                \item spatial self-localization
            \end{itemize}
        
        \ProjectDeepHeading{7}{Level 2: Agent Awareness}
            
            The system represents itself as the cause of actions.
            
            Examples:
            
            \begin{itemize}
                \item sense of agency
                \item action ownership
            \end{itemize}
        
        \ProjectDeepHeading{7}{Level 3: Self-Monitoring}
            
            The agent observes its own internal processes.
            
            Examples:
            
            \begin{itemize}
                \item monitoring internal states
                \item detecting errors in one's own actions
            \end{itemize}
        
        \ProjectDeepHeading{7}{Level 4: Meta-Cognitive Awareness}
            
            The agent can reason about its own cognition.
            
            Examples:
            
            \begin{itemize}
                \item evaluating confidence in decisions
                \item reasoning about knowledge or ignorance
            \end{itemize}
        
        \ProjectDeepHeading{7}{Level 5: Narrative or Reflective Self}
            
            The agent constructs a temporally extended model of itself.
            
            Examples:
            
            \begin{itemize}
                \item autobiographical memory
                \item long-term self-identity
            \end{itemize}
    
    
    \ProjectDeepHeading{6}{Relationship Between the Concepts}
        
        The hierarchy can be summarized as:
        
        \[
            \text{Body Awareness}
            \rightarrow
            \text{Agency Awareness}
            \rightarrow
            \text{Self-Monitoring}
            \rightarrow
            \text{Meta-Cognition}
            \rightarrow
            \text{Reflective/Narrative Self}
        \]
        
        Self-awareness emerges when a system forms internal models of itself and uses these models to monitor, evaluate, and regulate its behavior.

% Source: robot/part/mind/cognition/part/meta_cognition/parts/self_conciouness/self_awareness/knowledge.tex

\url{https://en.wikipedia.org/wiki/Self-knowledge_(psychology)}
    % 
    %

% Source: robot/part/mind/cognition/part/meta_cognition/parts/self_conciouness/self_awareness/neural.tex

\begin{verbatim}
class Neural:
    """
    \url{https://en.wikipedia.org/wiki/Neural_basis_of_self}
    """
\end{verbatim}

% Source: robot/part/mind/cognition/part/meta_cognition/parts/self_conciouness/self_awareness/self.tex

\ProjectDeepHeading{6}{Neural basis of self}
		\seehere{https://en.wikipedia.org/wiki/Neural_basis_of_self}
	
	\ProjectDeepHeading{6}{From autobiographical memory}
		In \cite{brewer-1986-what-is-autobiographical-memory} page 26, the parts of self is mentioned and is summarized as
		The working self, often referred to as just the 'self', is a set of active personal goals or self-images organized into goal hierarchies. These personal goals and self-images work together to modify cognition and the resulting behavior so an individual can operate effectively in the world.
		\seehere{https://en.wikipedia.org/wiki/Autobiographical_memory?utm_source=chatgpt.com\#Working_self}
	
	\cite{brewer-1986-what-is-autobiographical-memory} includes the self in the following hierarchy:
	
	\begin{figure}[H]
		\centering
		\begin{forest}
			for tree={
				grow=south,
				align=center
			}
			[Individual
				[Self
					[Ego\\(experience stream)]
					[Self-schema\\(knowledge about self)]
					[Autobiographical memory\\(episodes + facts)]
				]
			]
		\end{forest}
	\end{figure}
	
	\begin{itemize}
		\item \textbf{Autobiographical memory}
		\begin{itemize}
			\item Defined as memory for information related to the self.
		\end{itemize}
		
		\item \textbf{Structure of the self}
		\begin{itemize}
			\item The self consists of:
			\begin{itemize}
				\item an experiencing ego,
				\item a self-schema,
				\item and associated personal memories and autobiographical facts.
			\end{itemize}
		\end{itemize}
		
		\item \textbf{Ego}
		\begin{itemize}
			\item The conscious experiencing entity.
			\item Represents moment-to-moment subjective experience.
			\item Personal memories correspond to the ego’s experiences over time.
		\end{itemize}
		
		\item \textbf{Self-schema}
		\begin{itemize}
			\item A cognitive structure containing general knowledge about the self.
			\item Organized, relatively stable, and partly unconscious.
			\item Integrates private and publicly observable information about the self.
		\end{itemize}
		
		\item \textbf{Self}
		\begin{itemize}
			\item A complex structure including ego, self-schema, and self-related long-term memories.
		\end{itemize}
		
		\item \textbf{Individual}
		\begin{itemize}
			\item A broader entity including the self, body, depersonalized knowledge, and various cognitive and motor skills.
		\end{itemize}
	\end{itemize}

\ProjectDeepHeading{4}{Development stages}

% Source: robot/part/mind/cognition/part/meta_cognition/parts/self_conciouness/self_awareness/development_stages/level.tex

\url{https://en.wikipedia.org/wiki/Self-awareness\#Human_development}
    \ProjectDeepHeading{7}{Levels}
        \begin{itemize}
            \item Level 0—\href{https://en.wikipedia.org/wiki/Confusion}{Confusion}: The person is unaware of any mirror reflection or the mirroring itself; they perceive a mirror image as an extension of their environment.
        \end{itemize}
        \begin{itemize}
            \item Level 1—Differentiation: The individual realizes the mirror is able to reflect things. They see that what is in the mirror is of a different nature from what is surrounding them. At this level they can differentiate between their own movement in the mirror and the movement of the surrounding environment.
        \end{itemize}
        \begin{itemize}
            \item Level 2—Situation: The individual can link the movements on the mirror to what is perceived within their own body.
        \end{itemize}
        \begin{itemize}
            \item Level 3—\href{https://en.wikipedia.org/wiki/Identification_(psychology)}{Identification}: An individual can now see that what's in the mirror is not another person but actually them.
        \end{itemize}
        \begin{itemize}
            \item Level 4—Permanence: The individual is able to identify the self in previous pictures looking different or younger. A "permanent self" is now experienced.
        \end{itemize}
        \begin{itemize}
            \item Level 5—\href{https://en.wikipedia.org/wiki/Self-consciousness}{Self-consciousness} or "meta" self-awareness: At this level, not only is the self seen from a first person view but it is realized that it is also seen from a third person's view. A person who develops self consciousness begins to understand they can be in the mind of others: for instance, how they are seen from a public standpoint.
        \end{itemize}

    % 
    %

\ProjectDeepHeading{4}{Kind}

\ProjectDeepHeading{5}{Reflective}

% Source: robot/part/mind/cognition/part/meta_cognition/parts/self_conciouness/self_awareness/kind/reflective/reflective.tex

% 
    %

\ProjectDeepHeading{6}{Kind}

% Source: robot/part/mind/cognition/part/meta_cognition/parts/self_conciouness/self_awareness/kind/reflective/kind/existance.tex

% 
    %

% Source: robot/part/mind/cognition/part/meta_cognition/parts/self_conciouness/self_awareness/kind/reflective/kind/feeling.tex

% 
    %

% Source: robot/part/mind/cognition/part/meta_cognition/parts/self_conciouness/self_awareness/kind/reflective/kind/perception.tex

% 
    %

% Source: robot/part/mind/cognition/part/meta_cognition/parts/self_conciouness/self_awareness/kind/reflective/kind/thought.tex

% 
    %

\ProjectDeepHeading{5}{Social}

% Source: robot/part/mind/cognition/part/meta_cognition/parts/self_conciouness/self_awareness/kind/social/social.tex

% 
    %

\ProjectDeepHeading{6}{Kind}

% Source: robot/part/mind/cognition/part/meta_cognition/parts/self_conciouness/self_awareness/kind/social/kind/self_evaluation.tex

% 
    %

% Source: robot/part/mind/cognition/part/meta_cognition/parts/self_conciouness/self_awareness/kind/social/kind/self_monitoring.tex

% 
    %

% Source: robot/part/mind/cognition/part/meta_cognition/parts/self_conciouness/self_awareness/kind/social/kind/self_regulation.tex

% 
    %

\ProjectDeepHeading{4}{Novelty detection}

% Source: robot/part/mind/cognition/part/meta_cognition/parts/self_conciouness/self_awareness/novelty_detection/relationship_between_self_awareness_and_novelty_detection.tex

% 
    % 


`

% Source: robot/part/mind/cognition/part/object_level/object_level.tex

The object level is the lower cognitive level at which ordinary cognitive operations are performed, such as recognizing and discriminating objects, making decisions, encoding semantic information, and representing space. In the metacognitive model, information from the object level is monitored by the meta-level, while the meta-level can regulate or control object-level processes from above \cite{fleur-2021-metacognition-ideas-and-insights-from-neuro-and-educational-scie}.

\ProjectDeepHeading{1}{Architecture}

% Source: robot/part/mind/cognition/part/object_level/architecture/architecture.tex

\url{https://en.wikipedia.org/wiki/Cognitive_architecture}

\ProjectDeepHeading{1}{Behavior}

% Source: robot/part/mind/cognition/part/object_level/behavior/behavior.tex

\seeherefooter{https://en.wikipedia.org/wiki/Behavior}
    %
    %

% Source: robot/part/mind/cognition/part/object_level/process/process.tex

\ProjectDeepHeading{4}{Automatic vs controlled processes}
    pass

\ProjectDeepHeading{2}{Models}

% Source: robot/part/mind/cognition/part/object_level/process/models/models.tex

A cognitive model is a representation of one or more cognitive processes in humans or other animals for the purposes of comprehension and prediction.
    \seehere{https://en.wikipedia.org/wiki/Cognitive_model}
    % 
    %

\ProjectDeepHeading{3}{Executive functions}

% Source: robot/part/mind/cognition/part/object_level/process/parts/executive_functions/executive_functions.tex

executive functions (collectively referred to as executive function and cognitive control) are a set of cognitive processes that support goal-directed behavior, by regulating thoughts and actions through cognitive control, selecting and successfully monitoring actions that facilitate the attainment of chosen objectives.
    \seehere{https://en.wikipedia.org/wiki/Executive_functions}
    %
    %

\ProjectDeepHeading{4}{Parts}

% Source: robot/part/mind/cognition/part/object_level/process/parts/executive_functions/parts/attentional_control.tex

\seehere{https://en.wikipedia.org/wiki/Attentional_control}
    %
    %

% Source: robot/part/mind/cognition/part/object_level/process/parts/executive_functions/parts/cognitive_felexibility.tex

Cognitive flexibility[note 1] is an intrinsic property of a cognitive system often associated with the mental ability to adjust its activity and content, switch between different task rules and corresponding behavioral responses, maintain multiple concepts simultaneously and shift internal attention between them. \seehere{https://en.wikipedia.org/wiki/Cognitive_flexibility}
    %
    %

% Source: robot/part/mind/cognition/part/object_level/process/parts/executive_functions/parts/cognitive_inhibition.tex

Cognitive inhibition refers to the mind's ability to tune out stimuli that are irrelevant to the task/process at hand or to the mind's current state. Additionally, it can be done either in whole or in part, intentionally or otherwise. \seehere{https://en.wikipedia.org/wiki/Cognitive_inhibition}
    %
    %

% Source: robot/part/mind/cognition/part/object_level/process/parts/executive_functions/parts/inhibitory_control.tex

Inhibitory control is a cognitive process – and, more specifically, an executive function – that permits an individual to inhibit their impulses and natural, habitual, or dominant behavioral responses to stimuli (a.k.a. prepotent responses) in order to select a more appropriate behavior that is consistent with completing their goals. \seehere{https://en.wikipedia.org/wiki/Inhibitory_control}
    %
    %

% Source: robot/part/mind/cognition/part/object_level/process/parts/executive_functions/parts/working_memory.tex

\clickurlfooter{https://en.wikipedia.org/wiki/Working_memory} defines working memory as a cognitive system with a limited capacity that can hold information temporarily
    and mentions the following functions for working memory
    \begin{itemize}
        \item Temporary storage of information
        \item Manipulation of stored information
        \item Support for reasoning
        \item Guidance of decision-making and behavior
        \item Direction of attention toward relevant information
        \item Suppression of irrelevant information and inappropriate actions
        \item Coordination among multiple processes when more than one task is performed simultaneously
    \end{itemize}
    
    
    \ProjectDeepHeading{7}{Theories}
        \cite{baddeley-1974-working-memory} argue that short-term retention is not a single passive store but part of a broader working memory system. They propose that this system temporarily holds and processes information needed for reasoning, comprehension, and learning. Their account distinguishes a control component from temporary storage mechanisms, instead of treating short-term memory as one unitary box. The chapter uses experimental evidence from several tasks to show that immediate memory limits do not fully explain complex cognitive performance. Overall, the paper lays the foundation for viewing working memory as an active, limited-capacity system that supports ongoing cognition.
        the functions in the previous section can be reformulated by \cite{baddeley-1974-working-memory} memory as follows:
        \begin{itemize}
            \item \textbf{Central executive:} controls attention, suppresses irrelevant information, and coordinates concurrent processes and tasks
            \item \textbf{Phonological loop:} temporarily stores auditory and verbal information and refreshes it through rehearsal
            \item \textbf{Visuospatial sketchpad:} temporarily stores and manipulates visual and spatial information, such as mental imagery and mental maps
            \item \textbf{Episodic buffer:} integrates phonological, visual, and spatial information and links working memory with long-term memory
        \end{itemize}
        
        
        see \cite{cowan-2017-the-many-faces-of-working-memory}
        
        \cite{eriksson-2015-neurocognitive-architecture-of-working-memory} says that gola and task set (that we call actions are stored in working memory), see see Fig.~1 of \cite{eriksson-2015-neurocognitive-architecture-of-working-memory}
    
    
    \ProjectDeepHeading{7}{Planning and working memory}
        
        \cite{desposito-2015-the-cognitive-neuroscience-of-working-memory} argues that working memory is the system that keeps task-relevant information active so behavior can be guided when the needed information is not currently in the environment. It explicitly links working memory to plans, goals, and intentions: multiple plans or parts of a plan can be held at once, and interrupted plans remain active so they can later be resumed. Its main view is that the prefrontal cortex does not simply store sensory content; rather, it maintains abstract goals, rules, and task structure and sends top-down control signals that bias action selection. So, in this paper, cognition is goal-directed control: working memory holds and prioritizes goal-relevant representations, while prefrontal systems use them to organize and guide actions.
        
        \cite{das-2015-three-faces-of-cognitive-processes-theory-assessment-and-intervention}  presents the PASS model, which treats cognition as a set of distinct processes: Planning, Attention, Simultaneous processing, and Successive processing. It describes planning as a broad metacognitive and executive function involving decision making, judgment, and organizing behavior toward a goal. The author argues that planning is broader than working memory and explicitly says that limited-capacity working memory cannot contain future planning or a person’s major life goals.
        So, in this view, working memory supports some cognitive operations, but planning is a higher-level process that goes beyond temporary storage and updating.
    
    
    \ProjectDeepHeading{7}{Working memory and awarness}
        \cite{gutzwiller-2013-the-role-of-working-memory-in-levels-of-situation-awareness} examines how working memory relates to different levels of situation awareness in a dynamic decision-making task. Using three working-memory tasks, the authors built a factor-based working-memory measure and tested it in a firefighting simulation. They found no significant relationship between working memory and Level 1 situation awareness, which mainly involved detecting and recalling environmental elements. However, working memory was significantly related to Level 3 situation awareness, especially for prediction-oriented and implicit awareness measures. The study concludes that working memory contributes more to higher-level predictive awareness than to basic perceptual awareness.
    
    
    
    \ProjectDeepHeading{7}{Working memory and motor control}
        \cite{marvel-2019-how-the-motor-system-integrates-with-working-memory}
        
        \cite{mcdougle-2026-motor-working-memory}
    
    \ProjectDeepHeading{7}{Working memory and attention}
    
        \cite{oberauer-2019-working-memory-and-attention-a-conceptual-analysis-and-review}
    
    
    \ProjectDeepHeading{7}{Working memory and perception}
        \cite{wilkes-2005-working-memory-and-perception} treats perception as a functional module in a robotic cognitive architecture, not as a specifically localized brain region.
        This perceptual module identifies objects and events in the robot’s environment and supplies information to working memory, short-term memory, and spatial reasoning.
        Its processes can also be guided by working memory, short-term memory, and long-term memory, making perception part of a bidirectional system.
        In contrast, the paper links working memory more explicitly to neuroscience, especially computational models of prefrontal cortex circuits.
        So, in this article, perception is described mainly as an interacting cognitive module rather than as a named part of the brain or mind.

\ProjectDeepHeading{3}{Language}

% Source: robot/part/mind/cognition/part/object_level/process/parts/language/language.tex

All branches in wikipedia
    \url{https://en.wikipedia.org/wiki/Portal:Linguistics}
    
    
    \ProjectDeepHeading{6}{Definition from Saussure prespective}
        Language is a symbolic system that, in its most basic form, consists of sequences of replaceable components used to convey different meanings \cite{saussure-1916-course-in-general-linguistics}
        
        
    \ProjectDeepHeading{6}{Definition}
        Language is a system of conventional symbols (spoken, written, or signed) that enables members of a community to communicate ideas, experiences, and intentions.
        Encyclopedia Britannica
        
        \ProjectMissingAsset{/figspath language\_definition.png}
        
        \ProjectDeepHeading{7}{Phonetics}
        
        \ProjectDeepHeading{7}{Phonology}
        
        \ProjectDeepHeading{7}{Morphology}
        
        \ProjectDeepHeading{7}{Syntax}
        
        \ProjectDeepHeading{7}{Semantics}
        
        \ProjectDeepHeading{7}{Pragmatics}
    
    
    \ProjectDeepHeading{6}{Applied Linguistics}
    
    
    \ProjectDeepHeading{6}{Interdiciplinary}
        
        \ProjectDeepHeading{7}{Psycholinguistics}
        
        \ProjectDeepHeading{7}{Sociolinguistics}
        
        \ProjectDeepHeading{7}{Historical Linguistics}
        
        \ProjectDeepHeading{7}{Anthropological Linguistics}
        
        \ProjectDeepHeading{7}{Computational Linguistics}
        
        \ProjectDeepHeading{7}{Cognitive Linguistics}
        
        \ProjectDeepHeading{7}{Corpus Linguistics}
        
        \ProjectDeepHeading{7}{Biolinguistics}
            How does language ability roots in body and genetics
            Noam Chomsky ``Language Faculty''
        
        \ProjectDeepHeading{7}{Neurolinguistics}
    
    
    \ProjectDeepHeading{6}{Cognitive Neuroscience of Language}
    
    
    \ProjectDeepHeading{6}{Language Acquisition}
        \url{https://en.wikipedia.org/wiki/Language_acquisition}

% Source: robot/part/mind/cognition/part/object_level/process/parts/language/biological_prespective.tex

\ProjectDeepHeading{6}{Biological perspective}
Refs
\url{https://www.amazon.de/-/en/dp/0199284776/?coliid=I2ZFZ3O9PPZ6RM&colid=662UPZCH1AIN&psc=1&ref_=list_c_wl_lv_ov_lig_dp_it}
    
\url{https://www.amazon.de/-/en/dp/B08B1HS3XZ/?coliid=I2QQPO9D6WJT6H&colid=662UPZCH1AIN&psc=0&ref_=list_c_wl_lv_ov_lig_dp_it}

% Source: robot/part/mind/cognition/part/object_level/process/parts/language/computational_linguistic.tex

\ProjectDeepHeading{6}{Computational linguistic}

% Source: robot/part/mind/cognition/part/object_level/process/parts/language/grounding.tex

\ProjectDeepHeading{6}{Grounding}

% Source: robot/part/mind/cognition/part/object_level/process/parts/language/internal_language.tex

\ProjectDeepHeading{6}{Relation with self-awareness}
        \cite{fernyhough-2023-inner-speech-as-language-process-and-cognitive-tool} presents inner speech as an internal language process and cognitive tool involved in planning, attention, memory, emotion regulation, self-reflection, and abstract thinking. The authors argue that inner speech is not a uniform phenomenon, but appears in different forms such as dialogic, condensed, evaluative, motivational, and other-voice inner speech. A central point is that inner speech can help humans manipulate internal representations and improve cognitive performance, especially in demanding tasks. In the section on the self, the paper explains that inner speech can support self-observation, introspection, metacognition, and the formation of a coherent self-concept. However, the authors also note that inner speech may not fully explain self-awareness or self-regulation, because some evidence suggests that self-awareness itself can be a stronger explanatory factor. The review finally connects inner speech to embodied cognition, neuroscience, artificial intelligence, robotics, and brain–computer interfaces, while emphasizing the need to map different forms of inner speech to different cognitive functions.
    
    
    \ProjectDeepHeading{6}{Neuronal roots of inner speach}
        \cite{alderson-day-2016-the-brains-conversation-with-itself-neural-substrates-of-dialogic-inner-speech} uses fMRI compares dialogic inner speech, meaning inner speech shaped like a conversation, with monologic inner speech, meaning single-speaker internal speech. Dialogic inner speech produced stronger activation in a wider bilateral network, including superior temporal gyri, precuneus, posterior cingulate cortex, and inferior/medial frontal regions. The study suggests that internal dialogue is not merely silent self-talk, but also recruits social-cognitive and Theory-of-Mind-related processes, especially in the right posterior superior temporal gyrus.
    
    
    \ProjectDeepHeading{7}{Relation with self processing}
        \cite{molnar-szakacs-2013-self-processing-and-the-default-mode-network-interactions-with-the-mirror-neuron-system} argues that self-processing cannot be reduced to a single brain region or a homogeneous Default Mode Network. It distinguishes physical self-processing, linked more strongly to the Mirror Neuron System, from psychological self-referential processing, linked more strongly to Default Mode Network nodes such as the medial prefrontal cortex and the posterior cingulate cortex. The authors propose that the self emerges through dynamic interaction between embodied simulation and mentalizing systems.
    
    
    \ProjectDeepHeading{6}{Relation with Default Mode Network}
        \cite{menon-2023-20-years-of-the-default-mode-network-a-review-and-synthesis} reviews two decades of research on the Default Mode Network and explains how its role has expanded from a “resting-state” network to a central system for self-reference, social cognition, episodic memory, semantic memory, language, and mind wandering. The paper argues that the Default Mode Network is not a single-purpose system, but a set of interacting nodes and subnetworks whose functions depend on dynamic interactions with salience and frontoparietal control networks. Its main synthesis is that the Default Mode Network integrates memory, language, and semantic representations to construct a coherent internal narrative, which supports the sense of self and subjective continuity.

% Source: robot/part/mind/cognition/part/object_level/process/parts/language/neuro_linguistic.tex

\ProjectDeepHeading{6}{Neuro linguistc}

% Source: robot/part/mind/cognition/part/object_level/process/parts/language/origins.tex

\ProjectDeepHeading{6}{Origins}
    \begin{figure}
        \centering
        \ProjectMissingAsset{/figspathP(origins.pdf)}
        \caption{}
        \label{fig:}
    \end{figure}

\ProjectDeepHeading{4}{Communication}

\ProjectDeepHeading{5}{Model}

% Source: robot/part/mind/cognition/part/object_level/process/parts/language/communication/model/model.tex

\seeherefooter{https://en.wikipedia.org/wiki/Models_of_communication}

\ProjectDeepHeading{6}{Kind}

\ProjectDeepHeading{7}{Barnlund}

% Source: robot/part/mind/cognition/part/object_level/process/parts/language/communication/model/kind/barnlund/communication.tex

\ProjectDeepHeading{10}{Bragg theory}
        \cite{bragg-2021-transactional-communication-model-quick-look} presents the transactional communication model as an iterative process in which actors are simultaneously senders and receivers who encode, transmit, decode, and respond to messages. It emphasizes that effective communication depends not only on message transmission, but also on whether the intended audience receives and interprets the message as intended. The report highlights how schemas, worldview, cognitive bias, physiological/psychological noise, semantic noise, environmental noise, and other actors can distort encoding and decoding.
        
        \seeherefooter{https://en.wikipedia.org/wiki/Barnlund}
        
        \ProjectDeepHeading{11}{Intrapersonal}
            \seeherefooter{https://en.wikipedia.org/wiki/Barnlund\%27s_model_of_communication\#Intrapersonal_model}
            \begin{figure}[H]
                \centering
                \ProjectSafeIncludeGraphics[width=0.9\textwidth]{robot/part/mind/cognition/part/object_level/process/parts/language/communication/model/kind/barnlund/barnlund_intera_personal.png}
            \end{figure}
        
        \ProjectDeepHeading{11}{Interpersonal}
            \seeherefooter{https://en.wikipedia.org/wiki/Barnlund\%27s_model_of_communication\#Interpersonal_model}
            \begin{figure}[H]
                \centering
                \ProjectSafeIncludeGraphics[width=0.9\textwidth]{robot/part/mind/cognition/part/object_level/process/parts/language/communication/model/kind/barnlund/barnlund_interpersonal.png}
            \end{figure}

\ProjectDeepHeading{4}{Grammar}

% Source: robot/part/mind/cognition/part/object_level/process/parts/language/grammar/grammar.tex

FUG  = Functional Unification Grammar
    
    LFG  = Lexical Functional Grammar
    
    HPSG = Head-driven Phrase Structure Grammar
    
    PATR-II

\ProjectDeepHeading{4}{Grounding}

% Source: robot/part/mind/cognition/part/object_level/process/parts/language/grounding/grounding.tex

\cite{cohen-2024-a-survey-of-robotic-language-grounding-tradeoffs-between-symbols-and-embeddings} reviews robotic language grounding as a tradeoff between symbolic/formal representations and neural embedding-based approaches. Formal methods offer interpretability, lower data needs, and stronger safety guarantees, while end-to-end embedding methods offer flexibility but require more data and computation. The paper argues that future robotic language grounding should combine both: structured symbols for safety and reasoning, plus learned multimodal representations for generalization.

\ProjectDeepHeading{4}{Model}

\ProjectDeepHeading{5}{Kind}

\ProjectDeepHeading{6}{Usage based}

% Source: robot/part/mind/cognition/part/object_level/process/parts/language/model/kind/usage_based/used_based.tex

Usage-based models of language are linguistic approaches within a broader functional or cognitive framework that emerged in the late 1980s. These models assume a profound relationship between linguistic structure and usage \seeherefooter{https://en.wikipedia.org/wiki/Usage-based_models_of_language}

\ProjectDeepHeading{4}{Theory}

% Source: robot/part/mind/cognition/part/object_level/process/parts/language/theory/theory.tex

Theoretical linguistics, or general linguistics, is the branch of linguistics which inquires into the nature of language itself and seeks to answer fundamental questions as to what language is; how it works; how universal grammar (UG) as a domain-specific mental organ operates, if it exists at all; what its unique properties are; how language relates to other cognitive processes, etc. Theoretical linguists are most concerned with constructing models of linguistic knowledge and ultimately developing a linguistic theory.\seeherefooter{https://en.wikipedia.org/wiki/Theoretical_linguistics}

\ProjectDeepHeading{5}{Cognitive}

% Source: robot/part/mind/cognition/part/object_level/process/parts/language/theory/cognitive/cognitive.tex

Guides the assumption that linguistic patterns are patterns of conceptualization. Thus, cognitive linguists consider that the study of language provides insight into other human cognitive functions and vice-versa \seeherefooter{https://en.wikipedia.org/wiki/Cognitive_linguistics}

\ProjectDeepHeading{5}{Dependency grammar}

% Source: robot/part/mind/cognition/part/object_level/process/parts/language/theory/dependency_grammar/dependency_grammar.tex

Dependency grammar (DG) is a class of modern grammatical theories that are all based on the dependency relation (as opposed to the constituency relation of phrase structure) and that can be traced back primarily to the work of Lucien Tesnière. Dependency is the notion that linguistic units, e.g. words, are connected to each other by directed links. The (finite) verb is taken to be the structural center of clause structure. All other syntactic units (words) are either directly or indirectly connected to the verb in terms of the directed links, which are called dependencies \seeherefooter{https://en.wikipedia.org/wiki/Dependency_grammar}.

\ProjectDeepHeading{5}{Distributionalism}

% Source: robot/part/mind/cognition/part/object_level/process/parts/language/theory/distributionalism/distributionalism.tex

Distributionalism is a general theory of language and a discovery procedure for establishing elements and structures of language based on observed usage. The purpose of distributionalism was to provide a scientific basis for syntax as independent of meaning \seeherefooter{https://en.wikipedia.org/wiki/Distributionalism}.

\ProjectDeepHeading{5}{Functional}

% Source: robot/part/mind/cognition/part/object_level/process/parts/language/theory/functional/functional.tex

Functional linguistics is an approach to the study of language characterized by taking systematically into account the speaker's and the hearer's side, and the communicative needs of the speaker and of the given language community. Linguistic functionalism spawned in the 1920s to 1930s from Ferdinand de Saussure's systematic structuralist approach to language (1916) \seeherefooter{https://en.wikipedia.org/wiki/Functional_linguistics}.
    
    \cite{kay-1984-functional-unification-grammar} introduces Functional Unification Grammar (FUG) as a single monotonic formalism for representing linguistic descriptions in machine translation. Linguistic objects are represented as functional descriptions composed of attribute--value features, and compatible descriptions are combined through unification. Because FUG treats analysis, transfer, and synthesis within the same reversible framework, a translation from language A to language B should also support the reverse relation from B to A. A
    
    \cite{kay-1984-functional-unification-grammar} defines ``Feature structure`` as a functional description of a linguistic object. It is built from features, where each feature consists of an attribute and an associated value. A value may be atomic, or it may itself be another functional description, so the structure can be nested. Functional descriptions are combined by unification when they are compatible.

\ProjectDeepHeading{5}{Generative grammar}

% Source: robot/part/mind/cognition/part/object_level/process/parts/language/theory/generative_grammar/generative_grammar.tex

Generative grammar is a research tradition in linguistics that aims to explain the cognitive basis of language by formulating and testing explicit models of humans' subconscious grammatical knowledge \seeherefooter{https://en.wikipedia.org/wiki/Generative_grammar}.

\ProjectDeepHeading{5}{Interactional}

% Source: robot/part/mind/cognition/part/object_level/process/parts/language/theory/interactional/interactional.tex

Interactional linguistics (IL) is an interdisciplinary approach to grammar and interaction in the field of linguistics, that applies the methods of Conversation Analysis to the study of linguistic structures, including syntax, phonetics, morphology, and so on. Interactional linguistics is based on the principle that linguistic structures and uses are formed through interaction and it aims at understanding how languages are shaped through interaction. The approach focuses on temporality, activity implication and embodiment in interaction \seeherefooter{https://en.wikipedia.org/wiki/Interactional_linguistics}.
%    
%

\ProjectDeepHeading{5}{Phrase structure}

% Source: robot/part/mind/cognition/part/object_level/process/parts/language/theory/phrase_structure/phrase_structure.tex

\seeherefooter{https://en.wikipedia.org/wiki/Phrase_structure_grammar}

\ProjectDeepHeading{5}{Structuralism}

% Source: robot/part/mind/cognition/part/object_level/process/parts/language/theory/structuralism/structuralism.tex

Structural linguistics, or structuralism, in linguistics, denotes schools or theories in which language is conceived as a self-contained, self-regulating semiotic system whose elements are defined by their relationship to other elements within the syste \seeherefooter{https://en.wikipedia.org/wiki/Structural_linguistics}.

% Source: robot/part/mind/cognition/part/object_level/process/parts/language/theory/structuralism/computational.tex

\cite{mickus-2024-language-models-and-the-paradigmatic-axis} Mickus argues that language models and word embeddings can be interpreted in structuralist terms as models of the \textbf{paradigmatic axis}. Since language models estimate \(p(t \mid c)\), they model which lexical items can replace one another in a given context. The paper connects this idea to Saussurean associative series, Hjelmslevian paradigms, and Harrisian distributionally substitutable elements \cite{mickus-2024-language-models-and-the-paradigmatic-axis}.

\ProjectDeepHeading{6}{Glossematics}

% Source: robot/part/mind/cognition/part/object_level/process/parts/language/theory/structuralism/glossematics/glossematics.tex

Glossematics defines the glosseme as the smallest irreducible unit of both the content and expression planes of language; in the expression plane, the glosseme is nearly identical to the phoneme. In the content plane, it is the smallest unit of meaning which underlies a concept. A ewe, for example, consists of the taxemes sheep and female which may eventually be divided into even smaller units – glossemes – of meaning. The analysis is gradually expanded to the study of functions, more commonly known as dependencies, between elements on the level of discourse (which is called process), and between meaning and form in the linguistic system  \seeherefooter{https://en.wikipedia.org/wiki/Glossematics}.

\ProjectDeepHeading{6}{Saussure}

% Source: robot/part/mind/cognition/part/object_level/process/parts/language/theory/structuralism/saussure/saussure.tex

Saussure’s view \cite{saussure-1966-course-in-general-linguistics} is that language is not a list of names for things, but a system of signs, where each sign is formed by the relation between a signifier and a signified.
    
    The meaning or value of a sign does not come from itself alone, but from its differences and relations with other signs inside the whole language system.
    
    \begin{itemize}
        \item \textbf{Syntagmatic axis:} language elements combine in a linear sequence, for example words forming a sentence.
        \item \textbf{Paradigmatic axis:} language elements can replace each other in the same position, changing the meaning or value of the expression.
    \end{itemize}

\ProjectDeepHeading{5}{Transformational grammar}

% Source: robot/part/mind/cognition/part/object_level/process/parts/language/theory/transformational_grammar/transformational_grammar.tex

In linguistics, transformational grammar (TG) or transformational-generative grammar (TGG) was the earliest model of grammar proposed within the research tradition of generative grammar \seeherefooter{https://en.wikipedia.org/wiki/Transformational_grammar}.

\ProjectDeepHeading{5}{Unifacation base}

% Source: robot/part/mind/cognition/part/object_level/process/parts/language/theory/unifacation_base/unification_based.tex

FUG  = Functional Unification Grammar
    
    LFG  = Lexical Functional Grammar
    
    HPSG = Head-driven Phrase Structure Grammar
    
    PATR-II

\ProjectDeepHeading{5}{Universal grammar}

% Source: robot/part/mind/cognition/part/object_level/process/parts/language/theory/universal_grammar/universal_grammar.tex

The
    basic postulate of UG is that there are innate constraints on what the grammar of a possible human language could be. When linguistic stimuli are received in the course of language acquisition, children then adopt specific syntactic rules that conform to UG
    \seeherefooter{https://en.wikipedia.org/wiki/Universal_grammar}

\ProjectDeepHeading{6}{Part}

\ProjectDeepHeading{7}{Feature structure}

% Source: robot/part/mind/cognition/part/object_level/process/parts/language/theory/universal_grammar/part/feature_structure/feature_structure.tex

\seeherefooter{https://en.wikipedia.org/wiki/Feature_structure}
    \cite{kay-1984-functional-unification-grammar} defines ``Feature structure`` as a functional description of a linguistic object. It is built from features, where each feature consists of an attribute and an associated value. A value may be atomic, or it may itself be another functional description, so the structure can be nested. Functional descriptions are combined by unification when they are compatible.
    
    Feature structures \cite{kasper-1986-a-logical-semantics-for-feature-structures} represent linguistic objects as sets of attribute--value pairs. Kasper and Rounds formalize feature-structure descriptions as logical formulas interpreted by directed graphs.
    Their model explains unification and shows how disjunction and non-local path values increase computational complexity.

\ProjectDeepHeading{5}{Use based}

% Source: robot/part/mind/cognition/part/object_level/process/parts/language/theory/use_based/used_based.tex

Usage-based models of language are linguistic approaches within a broader functional or cognitive framework that emerged in the late 1980s. These models assume a profound relationship between linguistic structure and usage. \seeherefooter{https://en.wikipedia.org/wiki/Usage-based_models_of_language}

\ProjectDeepHeading{3}{Learning}

% Source: robot/part/mind/cognition/part/object_level/process/parts/learning/learning.tex

Changes the thinking system/ structure
    - makes model learned parameters
    - training config
    - architecture

\ProjectDeepHeading{4}{Motor}

% Source: robot/part/mind/cognition/part/object_level/process/parts/learning/motor/motor.tex

\url{https://en.wikipedia.org/wiki/Motor_learning}
    % 
    %

\ProjectDeepHeading{4}{Sensory motor coupling}

% Source: robot/part/mind/cognition/part/object_level/process/parts/learning/sensory_motor_coupling/sensory_motor_coupling.tex

Sensory-motor coupling is the coupling or integration of the sensory system and motor system. For a given stimulus, there is no one single motor command. \url{https://en.wikipedia.org/wiki/Sensory-motor_coupling}
    % 
    %

% Source: robot/part/mind/cognition/part/object_level/process/parts/learning/sensory_motor_coupling/by_teleportation_control.tex

\url{https://en.wikipedia.org/wiki/Sensory-motor_coupling}
    % 
    %

\ProjectDeepHeading{3}{Memory}

% Source: robot/part/mind/cognition/part/object_level/process/parts/memory/memory.tex

\ProjectDeepHeading{6}{Definition}
    Memory is the faculty of the mind by which data or information is encoded, stored, and retrieved when needed. It is the retention of information over time for the purpose of influencing future action. \seeherefooter{https://en.wikipedia.org/wiki/Memory}
    
    
    Most refrences such as~\cite[p.~3]{slotnick-2017-cognitive-neuroscience-of-memory}consider mental memory structure as a tree as follows
    \begin{figure}[H]
        \centering
        \ProjectSafeIncludeGraphics[width=0.9\textwidth]{robot/part/mind/cognition/part/object_level/process/parts/memory/slotnick-2017-cognitive-neuroscience-of-memory-figure-1-2.png}
        \caption{\cite{slotnick-2017-cognitive-neuroscience-of-memory} Fig. ~1.2}
    \end{figure}

\ProjectDeepHeading{4}{Part}

\ProjectDeepHeading{5}{Memorizing}

\ProjectDeepHeading{6}{Part}

\ProjectDeepHeading{7}{Explicity}

\ProjectDeepHeading{8}{Explicit}

\ProjectDeepHeading{9}{Long term}

\ProjectDeepHeading{10}{Part}

\ProjectDeepHeading{11}{Episodic}

% Source: robot/part/mind/cognition/part/object_level/process/parts/memory/part/memorizing/part/explicity/explicit/long_term/part/episodic/episodic.tex

\url{https://www.sciencedirect.com/science/article/abs/pii/S0028393209000645?via%3Dihub} says:
    There are essentially nine properties of episodic memory that collectively distinguish it from other types of memory. Other types of memory may exhibit a few of these properties, but only episodic memory has all nine:
    
    \begin{itemize}
        \item Contain summary records of sensory-perceptual-conceptual-affective processing.
        \item Retain patterns of activation/inhibition over long periods.
        \item Often represented in the form of (visual) images.
        \item They always have a perspective (field or observer).
        \item Represent short time slices of experience.
        \item They are represented on a temporal dimension roughly in order of occurrence.
        \item They are subject to rapid forgetting.
        \item They make autobiographical remembering specific.
        \item They are recollectively experienced when accessed.
    \end{itemize}
    
    \ProjectDeepHeading{14}{Definition of episodic and episode}
    Episodic Memory (EPMEM) automatically records snapshots of working memory in a temporal stream. Prior episodes can be retrieved into working memory through query. Once an episode has been retrieved, the next (or previous) episode can then be retrieved. An agent may employ EPMEM to sequentially play through episodes from its past (allowing it to predict the effects of actions), retrieve specific memories, or query for episodes possessing certain memory structures. \url{https://en.wikipedia.org/wiki/Soar_(cognitive_architecture)\#Episodic_memory}
    
    Episodic memory is the memory of events and experiences that can be recalled in relation to a specific time and in proper order \cite{tulving-2002-episodic-memory-from-mind-to-brain}
    
    Endel Tulving originally described episodic memory as a record of a person's experience that held temporally dated information and spatio-temporal relations \cite{tulving-1983-elements-of-episodic-memory}.
    
    A feature of episodic memory that Tulving later elaborates on is that it allows an agent to imagine traveling back in time \cite{tulving-2002-episodic-memory-from-mind-to-brain}.
    
    In the 1972 paper Tulving stated that episodic memory
    "stores information about temporally dated episodes or events and temporal-
    spatial relations among these events" (p. 385) \cite{tulving-1972-episodic-and-semantic-memory}
    
    
    \ProjectDeepHeading{14}{Episode}
    \cite{tulving-2002-episodic-memory-from-mind-to-brain} defines episodes as integrated representations of events situated in a spatiotemporal context.
    
    \ProjectDeepHeading{17}{Event} An event according to \cite{davidson-1969-the-individuation-of-events} Events are particular occurrences that can be described in terms of actions and their properties. Spatiotemporal is the association of location with time.
    
    \cite{morales-calva-2025-tell-me-why-the-missing-w-in-episodic-memorys-what-where-and-when} defines an eposide as it is shown in \ref{morales-calva-2025-tell-me-why-the-missing-w-in-episodic-memorys-what-where-and-when-figure-2}
        \begin{figure}[H]
            \centering
            \ProjectSafeIncludeGraphics[width=0.5\textwidth]{robot/part/mind/cognition/part/object_level/process/parts/memory/part/memorizing/part/explicity/explicit/long_term/part/episodic/morales-calva-2025-tell-me-why-the-missing-w-in-episodic-memorys-what-where-and-when-figure-2.png}
            \label{morales-calva-2025-tell-me-why-the-missing-w-in-episodic-memorys-what-where-and-when-figure-2}
            \caption{From \cite{morales-calva-2025-tell-me-why-the-missing-w-in-episodic-memorys-what-where-and-when} figure 2}
        \end{figure}

\ProjectDeepHeading{12}{Autobigraphical}

% Source: robot/part/mind/cognition/part/object_level/process/parts/memory/part/memorizing/part/explicity/explicit/long_term/part/episodic/autobigraphical/autobiographic.tex

According to \cite{brewer-1986-what-is-autobiographical-memory}, defined as memory for information related to the self additionally he includes the first three of the following as different forms of autobiographical memory
    \begin{enumerate}
        \item \textit{Personal memory}. I might experience a mental image that corresponds to a particular episode in my life, such as the time in California, while visiting Mt. Palomar, when I made a snowball and threw it down the trail at my sons, starting a snowball fight.
        \item \textit{Autobiographical fact}. I might simply recall the fact (with no accompanying imagery) that I drove to Mt. Palomar on another occasion when I was a senior in high school.
        \item \textit{Generic personal memory}. I might have a mental image of what, in general, it was like to drive north on Highway 1 in California’s Big Sur region: an image that does not appear to be of any specific moment during the drive but is a generic view from the driver’s seat of the car with the sharp jagged cliffs to my right and the Pacific stretching out to the horizon on my left.
        \item \textit{Semantic memory}. I might recall (with no accompanying imagery) that in California in 1978 the voters passed a limitation on property tax called Proposition 13.
        \item \textit{Generic perceptual memory}. I might have a visual image of the outline of the shape of the state of California.
    \end{enumerate}
    
    
    
    
    \url{https://en.wikipedia.org/wiki/Autobiographical_memory} says:
    Autobiographical memory (AM) is a memory system consisting of episodes recollected from an individual's life, based on a combination of episodic (personal experiences and specific objects, people and events experienced at particular time and place) and semantic (general knowledge and facts about the world) memory. It is thus a type of explicit memory.
    
    The autobiographical knowledge base contains knowledge of the self, used to provide information on what the self is, what the self was, and what the self can be. This information is categorized into three broad areas:
    \begin{itemize}
        \item lifetime periods
        \item general events
        \item event-specific knowledge
    \end{itemize}
    
    
    \ProjectDeepHeading{15}{Types}
        \begin{itemize}
            \item \textbf{Biographical or Personal:}
            These autobiographical memories contain personal biographical information, such as where a person was born or the names of their parents.
            
            \item \textbf{Copies vs.\ Reconstructions}
            \begin{itemize}
                \item \textbf{Copies:} Vivid autobiographical memories of an experience with a considerable amount of visual and sensory-perceptual detail. These memories may vary in authenticity.
                \item \textbf{Reconstructions:} Autobiographical memories that are not direct reflections of raw experiences, but are rebuilt by incorporating new information or interpretations formed in hindsight.
            \end{itemize}
            
            \item \textbf{Specific vs.\ Generic}
            \begin{itemize}
                \item \textbf{Specific:} Memories that contain detailed recollection of a particular event (event-specific knowledge).
                \item \textbf{Generic:} Memories that are vague and contain little detail beyond the type of event that occurred. Episodic autobiographical memories may also fall into this category when one remembered event represents a series of similar events.
            \end{itemize}
            
            \item \textbf{Field vs.\ Observer}
            \begin{itemize}
                \item \textbf{Field memories:} Memories recalled from the original perspective, experienced from a first-person point of view.
                \item \textbf{Observer memories:} Memories recalled from a perspective outside oneself, as if observing the event from a third-person viewpoint. Older memories are often recalled in this form and are more frequently reconstructions, while field memories tend to be more vivid and similar to copies.
            \end{itemize}
            
            \item \url{https://en.wikipedia.org/wiki/Autobiographical_memory\#Types}
        \end{itemize}

\ProjectDeepHeading{13}{Parts}

\ProjectDeepHeading{14}{Event specific knowledge}

% Source: robot/part/mind/cognition/part/object_level/process/parts/memory/part/memorizing/part/explicity/explicit/long_term/part/episodic/autobigraphical/parts/event_specific_knowledge/event_specific_knowledge.tex

\cite{pillemer-2001-momentous-events-and-the-life-story} describes Event-specific knowledge (ESK) as a vividly detailed information about individual events, often in the form of visual images and sensory-perceptual features. The high levels of detail in ESK fade very quickly, though certain memories for specific events tend to endure longer. Originating events (events that mark the beginning of a path towards long-term goals), turning points (events that re-direct plans from original goals), anchoring events (events that affirm an individual's beliefs and goals) and analogous events (past events that direct behaviour in the present) are all event specific memories that will resist memory decay.
    
    The sensory-perceptual details held in ESK, though short-lived, are a key component in distinguishing memory for experienced events from imagined events. In the majority of cases, it is found that the more ESK a memory contains, the more likely the recalled event has actually been experienced. Unlike lifetime periods and general events, ESK are not organized in their grouping or recall. Instead, they tend to simply 'pop' into the mind. ESK is also thought to be a summary of the content of episodic memories, which are contained in a separate memory system from the autobiographical knowledge base. This way of thinking could explain the rapid loss of event-specific detail, as the links between episodic memory and the autobiographical knowledge base are likewise quickly lost.
    \url{https://en.wikipedia.org/wiki/Autobiographical_memory?utm_source=chatgpt.com\#Autobiographical_knowledge_base}

\ProjectDeepHeading{15}{Remembering}

% Source: robot/part/mind/cognition/part/object_level/process/parts/memory/part/memorizing/part/explicity/explicit/long_term/part/episodic/autobigraphical/parts/event_specific_knowledge/remembering/remembering.tex

\cite{conway-2000-the-construction-of-autobiographical-memories-in-the-self-memory-system} says the process of remembering an event specific knowledge is as follows by order (This paper is the base paper for remebering event specific knowledge)
    \begin{enumerate}
        \item \textbf{WillingToRemember}\\
        The system intentionally initiates an attempt to recall a past experience.
        
        \item \textbf{RetrievalModeEntering}\\
        The mind enters a retrieval state in which memory construction becomes the current active task.
        
        \item \textbf{WorkingSelfGoalsActivating}\\
        Current self-related goals become active and constrain what kind of memory can be searched for and accepted.
        
        \item \textbf{CueElaborating}\\
        The initial cue is expanded and refined to make it more effective for accessing autobiographical knowledge.
        
        \item \textbf{VerificationCriteriaSetting}\\
        The system sets criteria for what will count as an acceptable retrieved memory.
        
        \item \textbf{RetrievalModelBuilding}\\
        A retrieval model is formed that combines cues, goals, and criteria to guide the search process.
        
        \item \textbf{KnowledgeBaseSearching}\\
        The autobiographical knowledge base is searched through spreading activation and controlled access.
        
        \item \textbf{GeneralEventAccessing}\\
        General event knowledge is typically accessed first, providing an initial structured entry point into memory.
        
        \item \textbf{StableMemoryPatternForming}\\
        A stable activation pattern is formed across relevant autobiographical knowledge structures.
        
        \item \textbf{EventSpecificKnowledgeActivating}\\
        Event-specific knowledge is activated, making detailed information about the target event available.
        
        \item \textbf{SensoryDetailsRetrieving}\\
        Sensory-perceptual details such as imagery, sounds, feelings, or other experiential fragments enter awareness.
    \end{enumerate}
    
    
    A probabale data flow from goal to planning
    
    \begin{verbatim}
SurvivalGoal
        ↓
        PlanningNeed
        ↓
        AutobiographicalMemoryRetrieval
        ↓
        RememberingMode
        ├── ConceptualRemembering
        └── PerceptualRemembering
        ↓
        Planning
    \end{verbatim}

% Source: robot/part/mind/cognition/part/object_level/process/parts/memory/part/memorizing/part/explicity/explicit/long_term/part/episodic/autobigraphical/parts/event_specific_knowledge/remembering/with_particle_filter.tex

\ProjectDeepHeading{18}{In human}
        \cite{socher-2009-a-bayesian-analysis-of-dynamics-in-free-recall} proposes LDA-TCM, a probabilistic model of free recall that combines the Temporal Context Model with latent Dirichlet allocation. It treats recall as inference over a drifting latent mental context, where remembered words are retrieved through both semantic and episodic paths. The authors use a particle filter for approximate posterior inference and show that the mixed semantic-episodic model predicts recall sequences better than purely semantic or purely episodic alternatives.
        
        \ProjectDeepHeading{19}{in human hippocampus}
            The hippocampus is mainly involved in forming and retrieving memories, especially episodic and associative memories, and in spatial navigation. In this paper, it is modeled as a Bayesian filtering system that combines memory retrieval and navigation in one mechanism.
            
            \cite{fox-2010-learning-in-a-unitary-coherent-hippocampus} extends the UCPF-HC(The name of their model) model, which treats the hippocampus as a unitary coherent particle filter combining associative memory and spatial navigation.
            The authors propose an online learning mechanism for CA3 connections, linking hippocampal learning to a temporal restricted Boltzmann machine and wake-sleep learning during theta cycles.
            The model suggests that CA3 can learn useful latent representations for denoising and state estimation, while the biological role of ADP and acetylcholine remains a hypothesis.
    
    
    \ProjectDeepHeading{18}{In robotics}
        \cite{ueda-2019-particle-filter-on-episode} proposes Particle Filter on Episode (PFoE), an algorithm that applies a particle filter to a recorded sequence of robot actions and sensor readings. The recorded sequence is called an episode, inspired by episodic memory, and is used to find situations similar to the robot’s current situation. Experiments with a small mobile robot show that PFoE can replay taught behaviors and recover from skids or interruptions without a separate learning phase.

\ProjectDeepHeading{12}{Context dependent}

% Source: robot/part/mind/cognition/part/object_level/process/parts/memory/part/memorizing/part/explicity/explicit/long_term/part/episodic/context_dependent/context_dependent.tex

context-dependent memory is the improved recall of specific episodes or information when the context present at encoding and retrieval are the same. \url{https://en.wikipedia.org/wiki/Context-dependent_memory}
    % 
    %

\ProjectDeepHeading{13}{Remembering}

% Source: robot/part/mind/cognition/part/object_level/process/parts/memory/part/memorizing/part/explicity/explicit/long_term/part/episodic/context_dependent/remembering/remebering.tex

Check the graph in neuro science memory
    % 
    %

\ProjectDeepHeading{14}{Recollection}

% Source: robot/part/mind/cognition/part/object_level/process/parts/memory/part/memorizing/part/explicity/explicit/long_term/part/episodic/context_dependent/remembering/recollection/recollection.tex

Check the graph in neuro science memory
    % 
    %

\ProjectDeepHeading{11}{Semantic}

% Source: robot/part/mind/cognition/part/object_level/process/parts/memory/part/memorizing/part/explicity/explicit/long_term/part/semantic/semantic.tex

context-dependent memory is the improved recall of specific episodes or information when the context present at encoding and retrieval are the same \seefootnotehere{https://en.wikipedia.org/wiki/Context-dependent_memory}

\ProjectDeepHeading{9}{Short term}

% Source: robot/part/mind/cognition/part/object_level/process/parts/memory/part/memorizing/part/explicity/explicit/short_term/short_term.tex

\seeherefooter{https://en.wikipedia.org/wiki/Short-term_memory} defines short memory as the capacity for holding a small amount of information in an active, readily available state for a short interval. For example, short-term memory holds a phone number that has just been recited. The duration of short-term memory (absent rehearsal or active maintenance) is estimated to be on the order of seconds. The commonly cited capacity of 7 items, found in Miller's law, has been superseded by 3-5 items.
    %
    %

\ProjectDeepHeading{5}{Remebering}

% Source: robot/part/mind/cognition/part/object_level/process/parts/memory/part/remebering/remebering.tex

\ProjectDeepHeading{8}{Relation with particle filter}
        Many papers including \cite{fox-2010-learning-in-a-unitary-coherent-hippocampus} consider hippocampus (the part in brain which is responsible for remebering) as a unitary coherent particle filter combining associative memory and spatial navigation.
    
    
    \ProjectDeepHeading{8}{Remebering autobiographical memory}
        For an overview of existing theories on how remembring autobiographical memory see \cite{sheldon-2019-a-neurocognitive-perspective-on-the-forms-and-functions-of-autobiographical-memory-retrieval}

\ProjectDeepHeading{5}{Shared}

\ProjectDeepHeading{6}{Engram}

% Source: robot/part/mind/cognition/part/object_level/process/parts/memory/part/shared/engram/engram.tex

An engram is a unit of cognitive information imprinted in a physical substance, theorized to be the means by which memories are stored[1] as biophysical or biochemical[2] changes in the brain or other biological tissue, in response to external stimuli \seeherefooter{https://en.wikipedia.org/wiki/Engram_(neuropsychology)}.
    
    \seeherefooter{https://pubmed.ncbi.nlm.nih.gov/31896692/}
    
    \seeherefooter{file:///home/donkarlo/Dropbox/repo/nd_shared_research/bibliography/corpora/1756-6606-5-32.pdf}
    
    \cite{sakaguchi-2012-catching-the-engram-strategies-to-examine-the-memory-trace}

\ProjectDeepHeading{3}{Perception}

% Source: robot/part/mind/cognition/part/object_level/process/parts/perception/perception.tex

\clickurlfooter{https://en.wikipedia.org/wiki/Perception} defines perception as organization, identification, and interpretation of sensory information, in order to represent and understand the presented information or environment
    %
    %

% Source: robot/part/mind/cognition/part/object_level/process/parts/perception/philosophy.tex

\clickurlfooter{https://en.wikipedia.org/wiki/Philosophy_of_perception} defines the philosophy of perception is concerned with the nature of perceptual experience and the status of perceptual data, in particular how they relate to beliefs about, or knowledge of, the world
    %
    %

\ProjectDeepHeading{4}{Motion}

% Source: robot/part/mind/cognition/part/object_level/process/parts/perception/motion/motion.tex

\url{https://en.wikipedia.org/wiki/Motion_perception}
    
    \url{https://en.wikipedia.org/wiki/Action-specific_perception?utm_source=chatgpt.com}
    %
    %

\ProjectDeepHeading{4}{State estimation}

\ProjectDeepHeading{5}{Kind}

\ProjectDeepHeading{6}{Bayesian}

\ProjectDeepHeading{7}{Kalman}

% Source: robot/part/mind/cognition/part/object_level/process/parts/perception/state_estimation/kind/bayesian/kalman/kalman.tex

The Kalman filter estimates the hidden state of a linear dynamical system from noisy measurements. It assumes that both the process noise and the measurement noise are Gaussian \cite{kalman-1960-a-new-approach-to-linear-filtering-and-prediction-problems}.
    
    
    \ProjectDeepHeading{10}{State-Space Model}
        
        The system is usually written as:
        
        \[
            x_k = A_k x_{k-1} + B_k u_k + w_k
        \]
        
        \[
            z_k = H_k x_k + v_k
        \]
        
        where:
        
        \[
            x_k
        \]
        
        is the hidden state of the system at time step \(k\).
        
        \[
            z_k
        \]
        
        is the measurement or observation at time step \(k\).
        
        \[
            u_k
        \]
        
        is the control input.
        
        \[
            A_k
        \]
        
        is the state transition matrix, which describes how the previous state moves to the current state.
        
        \[
            B_k
        \]
        
        is the control-input matrix.
        
        \[
            H_k
        \]
        
        is the observation matrix, which maps the hidden state into the measurement space.
        
        \[
            w_k
        \]
        
        is the process noise, usually assumed to be Gaussian:
        
        \[
            w_k \sim \mathcal{N}(0, Q_k)
        \]
        
        \[
            v_k
        \]
        
        is the measurement noise, also usually assumed to be Gaussian:
        
        \[
            v_k \sim \mathcal{N}(0, R_k)
        \]
        
        Here, \(Q_k\) is the process-noise covariance matrix, and \(R_k\) is the measurement-noise covariance matrix.
    
    
    \ProjectDeepHeading{10}{Prediction Step}
        
        In the prediction step, the filter predicts the next state before seeing the new measurement.
        
        The predicted state is:
        
        \[
            \hat{x}_{k|k-1} = A_k \hat{x}_{k-1|k-1} + B_k u_k
        \]
        
        The predicted covariance is:
        
        \[
            P_{k|k-1} = A_k P_{k-1|k-1} A_k^T + Q_k
        \]
        
        where:
        
        \[
            \hat{x}_{k|k-1}
        \]
        
        is the predicted state at time \(k\), based only on information available up to time \(k-1\).
        
        \[
            P_{k|k-1}
        \]
        
        is the predicted uncertainty of this state estimate.
        
        \[
            P_{k-1|k-1}
        \]
        
        is the previous corrected covariance.
    
    
    \ProjectDeepHeading{10}{Innovation}
        
        After the measurement \(z_k\) arrives, the filter computes the difference between the actual measurement and the predicted measurement.
        
        This difference is called the innovation or residual:
        
        \[
            y_k = z_k - H_k \hat{x}_{k|k-1}
        \]
        
        The innovation covariance is:
        
        \[
            S_k = H_k P_{k|k-1} H_k^T + R_k
        \]
        
        where \(y_k\) tells us how much the real measurement differs from what the model expected.
    
    
    \ProjectDeepHeading{10}{Kalman Gain}
        
        The Kalman gain decides how much the filter should trust the new measurement compared with the model prediction.
        
        \[
            K_k = P_{k|k-1} H_k^T S_k^{-1}
        \]
        
        or equivalently:
        
        \[
            K_k = P_{k|k-1} H_k^T
            \left(
                H_k P_{k|k-1} H_k^T + R_k
            \right)^{-1}
        \]
        
        If the measurement noise \(R_k\) is large, the Kalman gain becomes smaller, so the filter trusts the model more.
        
        If the predicted uncertainty \(P_{k|k-1}\) is large, the Kalman gain becomes larger, so the filter trusts the measurement more.
    
    
    \ProjectDeepHeading{10}{Correction Step}
        
        The corrected state estimate is:
        
        \[
            \hat{x}_{k|k} =
            \hat{x}_{k|k-1} + K_k y_k
        \]
        
        or:
        
        \[
            \hat{x}_{k|k} =
            \hat{x}_{k|k-1} +
            K_k
            \left(
                z_k - H_k \hat{x}_{k|k-1}
            \right)
        \]
        
        The corrected covariance is:
        
        \[
            P_{k|k} =
            \left(
                I - K_k H_k
            \right)
            P_{k|k-1}
        \]
        
        where:
        
        \[
            \hat{x}_{k|k}
        \]
        
        is the updated state estimate after using the measurement at time \(k\).
        
        \[
            P_{k|k}
        \]
        
        is the updated uncertainty after the correction step.
    
    
    \ProjectDeepHeading{10}{Compact Form}
        
        The Kalman filter can be summarized in two main phases.
        
        First, prediction:
        
        \[
            \hat{x}_{k|k-1} = A_k \hat{x}_{k-1|k-1} + B_k u_k
        \]
        
        \[
            P_{k|k-1} = A_k P_{k-1|k-1} A_k^T + Q_k
        \]
        
        Second, correction:
        
        \[
            K_k =
            P_{k|k-1} H_k^T
            \left(
                H_k P_{k|k-1} H_k^T + R_k
            \right)^{-1}
        \]
        
        \[
            \hat{x}_{k|k} =
            \hat{x}_{k|k-1}
            +
            K_k
            \left(
                z_k - H_k \hat{x}_{k|k-1}
            \right)
        \]
        
        \[
            P_{k|k} =
            \left(
                I - K_k H_k
            \right)
            P_{k|k-1}
        \]
    
    
    \ProjectDeepHeading{10}{Interpretation}
        
        The Kalman filter is a recursive estimator.
        It does not need the full history of measurements.
        At each time step, it only needs the previous estimate, the previous covariance, the current control input, and the current measurement.
        
        The prediction step answers:
        
        \[
            \text{What do I expect the state to be before seeing the new measurement?}
        \]
        
        The correction step answers:
        
        \[
            \text{How should I modify that expectation after seeing the new measurement?}
        \]
        
        The covariance matrix \(P\) represents uncertainty about the state estimate.
        The matrices \(Q\) and \(R\) control how uncertainty enters the system: \(Q\) describes uncertainty in the system dynamics, while \(R\) describes uncertainty in the sensor measurements.

% Source: robot/part/mind/cognition/part/object_level/process/parts/perception/state_estimation/kind/bayesian/kalman/biological.tex

\cite{wolpert-2007-probabilistic-models-in-human-sensorimotor-control} argues that human sensorimotor control can be understood through Bayesian decision theory, because the brain must estimate body and world states under sensory and motor uncertainty. The paper reviews evidence that humans combine sensory cues according to their reliability, use prior beliefs in estimation, and may use Kalman-filter-like processes for time-varying state estimation. It also connects estimation to action selection, suggesting that movement planning depends on uncertainty, loss functions, and signal-dependent motor noise.

\ProjectDeepHeading{7}{Particle}

% Source: robot/part/mind/cognition/part/object_level/process/parts/perception/state_estimation/kind/bayesian/particle/particle.tex

\cite{gordon-1993-novel-approach-to-nonlinear-non-gaussian-bayesian-state-estimation} introduced particle filter as a sequential Monte Carlo method used to estimate the hidden state of a dynamic system from noisy observations.
    
    Assume that the hidden state at time $t$ is $x_t$ and the observation is $y_t$.
    
    \begin{align}
        x_t &\sim p(x_t \mid x_{t-1}) \\
        y_t &\sim p(y_t \mid x_t)
    \end{align}
    
    The first equation describes how the state changes over time.
    The second equation describes how the observation is generated from the hidden state.
    
    Instead of representing the posterior distribution exactly,
    
    \begin{equation}
        p(x_t \mid y_{1:t})
    \end{equation}
    
    a particle filter approximates it using many weighted samples, called particles:
    
    \begin{equation}
        p(x_t \mid y_{1:t}) \approx \sum_{i=1}^{N} w_t^{(i)} \ensuremath{\delta}(x_t - x_t^{(i)})
    \end{equation}
    
    where $x_t^{(i)}$ is the $i$-th particle, $w_t^{(i)}$ is its weight, and $N$ is the number of particles.
    
    At each time step, the algorithm has three main steps.
    
    
    \ProjectDeepHeading{10}{Prediction}
        
        Each particle is moved forward using the state transition model:
        
        \begin{equation}
            x_t^{(i)} \sim p(x_t \mid x_{t-1}^{(i)})
        \end{equation}
    
    
    \ProjectDeepHeading{10}{Update}
        
        Each particle receives a weight according to how well it explains the new observation:
        
        \begin{equation}
            w_t^{(i)} \propto p(y_t \mid x_t^{(i)})
        \end{equation}
        
        Then the weights are normalized:
        
        \begin{equation}
            w_t^{(i)} =
            \frac{w_t^{(i)}}{\sum_{j=1}^{N} w_t^{(j)}}
        \end{equation}
    
    
    \ProjectDeepHeading{10}{Resampling}
        
        Particles with large weights are copied more often, and particles with small weights are removed.
        
        \begin{equation}
            \{x_t^{(i)}, w_t^{(i)}\}_{i=1}^{N}
            \rightarrow
            \{\tilde{x}_t^{(i)}, \frac{1}{N}\}_{i=1}^{N}
        \end{equation}
        
        After resampling, all particles have equal weight again.
    
    
    \ProjectDeepHeading{10}{Simple Interpretation}
        
        A particle filter represents belief about the hidden state by many guesses.
        Each guess is tested against the new observation.
        Good guesses survive, and bad guesses disappear.
        Over time, the cloud of particles follows the probable state of the system.

\ProjectDeepHeading{8}{Discrete}

% Source: robot/part/mind/cognition/part/object_level/process/parts/perception/state_estimation/kind/bayesian/particle/discrete/discrete.tex

\cite{chopin-2001-sequential-inference-and-state-number-determination-for-discrete-state-space-models-through-particle-filtering} applies particle filtering to discrete state-space models, especially hidden Markov models with unobserved regime/state sequences. The paper proposes a Monte Carlo HMM filter that marginalizes hidden states to reduce particle-weight volatility and improve sequential Bayesian inference. It also introduces sequential state-number determination, estimating how many distinct hidden regimes have appeared over time.
    
    
    \ProjectDeepHeading{11}{Discrete State-Space Model}
        
        In a discrete state-space model, the observed signal \(y_t\) is assumed to be generated by an unobserved discrete state \(z_t\). If the hidden state at time \(t\) is \(k\), the observation model can be written as
        
        \[
            y_t
            \mid
            z_t = k
            \sim
            f_{\ensuremath{\xi}_k}(y_t),
        \]
        
        where \(f_{\ensuremath{\xi}_k}\) is the observation distribution associated with discrete state \(k\).
        
        The hidden state evolves in a discrete state space according to
        
        \[
            \Pr
            \left(
                z_t = l
                \mid
                z_{t-1} = k
            \right)
            =
            p_{kl},
        \]
        
        where \(p_{kl}\) is the transition probability from state \(k\) to state \(l\).
    
    
    \ProjectDeepHeading{11}{Particle Approximation}
        
        Particle filtering approximates the posterior distribution over the unknown parameters and hidden states by a set of weighted particles:
        
        \[
            \ensuremath{\pi}_t
            =
            \ensuremath{\pi}
            \left(
                \ensuremath{\theta},
            z_{1:t}
                \mid
                y_{1:t}
            \right)
            \approx
            \left\{
                \left(
                    \ensuremath{\theta}^{(j)},
                z_{1:t}^{(j)},
                w_t^{(j)}
                \right)
            \right\}_{j=1}^{H}.
        \]
        
        Here, \(\ensuremath{\theta}\) denotes the unknown model parameters, \(z_{1:t}\) is the hidden state sequence, \(y_{1:t}\) is the observed sequence, \(w_t^{(j)}\) is the weight of particle \(j\), and \(H\) is the number of particles.
        
        When a new observation \(y_t\) arrives, each particle is updated according to how well it explains the new observation:
        
        \[
            w_t^{(j)}
            \propto
            w_{t-1}^{(j)}
            p
            \left(
                y_t
                \mid
                y_{1:t-1},
                \ensuremath{\theta}^{(j)}
            \right).
        \]
    
    
    \ProjectDeepHeading{11}{Marginalization of Discrete Hidden States}
        
        In the Monte Carlo HMM filter, the discrete hidden states are marginalized out. Therefore, the likelihood of the new observation is computed by summing over all possible discrete states:
        
        \[
            p
            \left(
                y_t
                \mid
                y_{1:t-1},
                \ensuremath{\theta}^{(j)}
            \right)
            =
            \sum_{k=1}^{K}
            \Pr
            \left(
                z_t = k
                \mid
                y_{1:t-1},
                \ensuremath{\theta}^{(j)}
            \right)
            f_{\ensuremath{\xi}_k^{(j)}}(y_t).
        \]
        
        This means that each particle does not need to carry one sampled hidden-state trajectory. Instead, it carries a probability distribution over the discrete states.
        
        The filtering probability of each state can then be estimated by the weighted average
        
        \[
            \widehat{\Pr}
            \left(
                z_t = k
                \mid
                y_{1:t}
            \right)
            =
            \frac{
                \sum_{j=1}^{H}
                w_t^{(j)}
                \Pr
                \left(
                    z_t = k
                    \mid
                    y_{1:t},
                    \ensuremath{\theta}^{(j)}
                \right)
            }{
                \sum_{j=1}^{H}
                w_t^{(j)}
            }.
        \]
    
    
    \ProjectDeepHeading{11}{General Interpretation}
        
        The method combines three elements:
        
        \[
            \text{state transition probabilities}
            \quad
            p_{kl},
        \]
        
        \[
            \text{observation likelihoods}
            \quad
            f_{\ensuremath{\xi}_k}(y_t),
        \]
        
        and
        
        \[
            \text{particle weights}
            \quad
            w_t^{(j)}.
        \]
        
        Thus, particle filtering in a discrete state-space model estimates the current hidden regime by combining how states evolve, how well each state explains the observation, and how probable each parameter particle is.

\ProjectDeepHeading{5}{Odometry}

% Source: robot/part/mind/cognition/part/object_level/process/parts/perception/state_estimation/odometry/odometry.tex

Odometry is the use of data from motion sensors to estimate change in position over time. \url{https://en.wikipedia.org/wiki/Odometry}
    % 
    %

\ProjectDeepHeading{5}{Slam}

% Source: robot/part/mind/cognition/part/object_level/process/parts/perception/state_estimation/slam/slam.tex

SLAM is mapping+localization +state estimation

\ProjectDeepHeading{6}{Kind}

\ProjectDeepHeading{7}{Graph based}

% Source: robot/part/mind/cognition/part/object_level/process/parts/perception/state_estimation/slam/kind/graph_based/graph_based.tex

Represent SLAM as a graph:
    \begin{itemize}
        \item Nodes = robot poses
        \item Edges = constraints (odometry, loop closure)
    \end{itemize}

\ProjectDeepHeading{7}{Learning based}

\ProjectDeepHeading{8}{Neural}

% Source: robot/part/mind/cognition/part/object_level/process/parts/perception/state_estimation/slam/kind/learning_based/neural/neural.tex

Emerging direction.
    Examples:
    \begin{itemize}
    
        \item DeepVO
        \item Neural Radiance Fields (NeRF-SLAM)
        \item Learned depth + pose estimation
    \end{itemize}
    
    Still research-heavy.

\ProjectDeepHeading{7}{Particle filter}

% Source: robot/part/mind/cognition/part/object_level/process/parts/perception/state_estimation/slam/kind/particle_filter/particle_filter.tex

use Particle filter for robot pose and EKF per landmark

\ProjectDeepHeading{6}{Localization}

% Source: robot/part/mind/cognition/part/object_level/process/parts/perception/state_estimation/slam/localization/localization.tex

Localization is a special kind of state estimation in which the agent tries to only find its location and not further derivatives such as velocity, acceleration or jerk

\ProjectDeepHeading{4}{Temporal sequence}

\ProjectDeepHeading{5}{Kind}

\ProjectDeepHeading{6}{Sequential}

% Source: robot/part/mind/cognition/part/object_level/process/parts/perception/temporal_sequence/kind/sequential/sequential.tex

In cognitive psychology, sequence learning is inherent to human ability because it is an integrated part of conscious and nonconscious learning as well as activities. Sequences of information or sequences of actions are used in various everyday tasks: "from sequencing sounds in speech, to sequencing movements in typing or playing instruments, to sequencing actions in driving an automobile \seeherefooter{https://en.wikipedia.org/wiki/Sequence_learning}

\ProjectDeepHeading{6}{Temporal}

% Source: robot/part/mind/cognition/part/object_level/process/parts/perception/temporal_sequence/kind/temporal/temporal.tex

In psychology and neuroscience, time perception or chronoception is the subjective experience, or sense, of time, which is measured by someone's own perception of the duration of the indefinite and unfolding of events \seeherefooter{https://en.wikipedia.org/wiki/Time_perception}

% Source: robot/part/mind/cognition/part/object_level/process/parts/thinking/thinking.tex

None architectural
    (work on existing learned
    structures)

\ProjectDeepHeading{4}{Concept learning}

% Source: robot/part/mind/cognition/part/object_level/process/parts/thinking/concept_learning/concept_learning.tex

Concept learning is defined by Bruner, Goodnow, \& Austin (1956) as ``the search for and testing of attributes that can be used to distinguish exemplars from non exemplars of various categories'' \seeherefooter{https://en.wikipedia.org/wiki/Concept_learning}

\ProjectDeepHeading{4}{Decision making}

% Source: robot/part/mind/cognition/part/object_level/process/parts/thinking/decision_making/decision_making.tex

decision-making is regarded as the cognitive process resulting in the selection of a belief or a course of action among several possible alternative options. It could be either rational or irrational.
    \seeherefooter{https://en.wikipedia.org/wiki/Decision-making}
    
    \ProjectDeepHeading{7}{Automated decision making}
        Automated decision-making (ADM) is the use of data, machines and algorithms to make decisions
        \seeherefooter{https://en.wikipedia.org/wiki/Automated_decision-making}
    %
    %

% Source: robot/part/mind/cognition/part/object_level/process/parts/thinking/planning/mental_time_travel.tex

mental time travel is the capacity to mentally reconstruct personal events from the past (episodic memory) as well as to imagine possible scenarios in the future (episodic foresight/episodic future thinking). \seeherefooter{https://en.wikipedia.org/wiki/Mental_time_travel}
    %
    %

% Source: robot/part/mind/cognition/part/object_level/process/parts/thinking/planning/perception_action_cycle.tex

\cite{khilkevich-2024-brain-wide-dynamics-linking-sensation-to-action-during-decision-making}  shows that, after learning, sensory evidence is integrated across many brain regions in mice rather than in a single specialized area. The integrated evidence drives movement-preparatory activity, and this preparatory activity is distinct from the neural dynamics that execute the movement. Overall, the study argues that the transformation from sensation to action is brain-wide, parallel, and shaped by learning.

\ProjectDeepHeading{5}{Motion}

% Source: robot/part/mind/cognition/part/object_level/process/parts/thinking/planning/motion/motion.tex

Motion planning, also path planning (also known as the navigation problem or the piano mover's problem) is a computational problem to find a sequence of valid configurations that moves the object from the source to destination.
    \seeherefooter{https://en.wikipedia.org/wiki/Motion_planning}
    %
    %

% Source: robot/part/mind/cognition/part/object_level/process/parts/thinking/planning/motion/sensorimotor_transformation.tex

\url{https://en.wikipedia.org/wiki/Common_coding_theory?utm_source=chatgpt.com}
    
    \url{https://en.wikipedia.org/wiki/Sensory-motor_coupling?utm_source=chatgpt.com}
    %
    %

\ProjectDeepHeading{6}{Navigation}

% Source: robot/part/mind/cognition/part/object_level/process/parts/thinking/planning/motion/navigation/navigation.tex

%
    %

\ProjectDeepHeading{7}{Path planning}

% Source: robot/part/mind/cognition/part/object_level/process/parts/thinking/planning/motion/navigation/path_planning/path_planning.tex

Computes a geometric path from a start position to a goal position.
    
    The output is usually a sequence of waypoints or a continuous curve in space.
    
    Common algorithms: A*, D*, Dijkstra, RRT*, PRM.

\ProjectDeepHeading{7}{Trajectory planning}

% Source: robot/part/mind/cognition/part/object_level/process/parts/thinking/planning/motion/navigation/trajectory_planning/trajectory_planning.tex

Converts the geometric path into a time-parameterized trajectory while respecting dynamic and kinematic constraints (velocity, acceleration, jerk limits).
    
    Output: a trajectory that can be executed by the controller.

\ProjectDeepHeading{3}{Transition}

% Source: robot/part/mind/cognition/part/object_level/process/parts/transition/transition.tex

\cite{smith-2018-role-of-the-default-mode-network-in-cognitive-transitions} argues that the default mode network is involved not only in internally directed thought, but also in major cognitive transitions during externally focused tasks.
    Using fMRI task-switching experiments, the authors found that core and medial temporal DMN regions become more active during large task switches, brief rest periods, and task restarts after rest.
    The main conclusion is that the DMN likely represents scene, episode, or context, and that this context representation becomes prominent again when cognition shifts from one task state to another.

\ProjectDeepHeading{1}{Theory}

% Source: robot/part/mind/cognition/part/object_level/theory/piaget_cognitive_development.tex

\url{https://en.wikipedia.org/wiki/Piaget%27s_theory_of_cognitive_development}
    \\
    \url{}
    % 
    %

\ProjectDeepHeading{2}{Embodied}

% Source: robot/part/mind/cognition/part/object_level/theory/embodied/embodied.tex

Embodied cognition represents a diverse group of theories which investigate how cognition is shaped by the bodily state and capacities of the organism. Proponents of the embodied cognition thesis emphasize the active and significant role the body plays in the shaping of cognition and in the understanding of an agent's mind and cognitive capacities.
    \url{https://en.wikipedia.org/wiki/Embodied_cognition}
    
    %
    %

\paragraph{Theories}

\subparagraph{Connectionism}

% Source: robot/part/mind/cognition/theories/connectionism/connectionism.tex

Connectionism is an approach to the study of human mental processes and cognition that utilizes mathematical models known as connectionist networks or artificial neural networks \url{https://en.wikipedia.org/wiki/Connectionism}
    % 
    %

\subsubsection{Epistomology}

% Source: robot/part/mind/epistomology/epistomology.tex

Epistomology is the source of knowledge.
    
    To discover how knowledge arises, they investigate sources of justification, such as perception, introspection, memory, reason, and testimony.\url{https://en.wikipedia.org/wiki/Epistemology}

% Source: robot/part/mind/epistomology/knowedge_hierarchy.tex

\url{https://en.wikipedia.org/wiki/DIKW_pyramid}
    % 
    %

\paragraph{Theory}

% Source: robot/part/mind/epistomology/theory/genetic.tex

\url{https://en.wikipedia.org/wiki/Genetic_epistemology}
    % 
    %

\subsubsection{Event}

% Source: robot/part/mind/event/event.tex

\url{https://en.wikipedia.org/wiki/Mental_event}
    % 
    %

\subsubsection{Goal}

% Source: robot/part/mind/goal/goal.tex

A goal or objective is an idea of the future or desired result that a person or a group of people envision, plan, and commit to achieve. \seehere{https://en.wikipedia.org/wiki/Goal}
    
    \cite{trapp-2017-commentary-on-the-joys-of-perceiving-affect-as-feedback-for-perceptual-predictions} discusses predictive coding and the idea that the brain minimizes surprise using internal models. It addresses the “dark room problem” by introducing affect as a motivational signal. Positive affect is proposed to arise from successful predictions, especially under uncertainty. This framework links perception with reinforcement learning and the exploration–exploitation dilemma. It suggests that novelty is intrinsically rewarding because it can increase positive affect.
    
    \cite{sternson-2013-hypothalamic-survival-circuits-blueprints-for-purposive-behaviors} reviews how hypothalamic neural circuits generate goal-directed behaviors essential for survival, such as feeding, aggression, and reproduction.
    It shows that small, molecularly defined neuron populations (e.g., AGRP neurons) can rapidly trigger complex motivated behaviors when activated.
    It proposes that survival needs provide entry points to map neural circuits linking internal states to behavioral intent and intensity.
    
    
    \ProjectDeepHeading{1}{Goal seeking}
        \cite{kerr-2017-biological-goal-seeking} presents a biologically inspired approach to goal-seeking behavior in autonomous agents. It models how agents can move toward goals while avoiding obstacles using simple, reactive mechanisms. The method combines elements like potential fields and finite-state control to generate adaptive behavior.
        It demonstrates how complex navigation can emerge from relatively simple rules without global planning. The work is evaluated in scenarios such as dynamic environments and predator–prey interactions.
        
    \ProjectDeepHeading{1}{Motor goal}
        A motor goal is a neurally planned motor outcome that is used to organize motor control.
        \seeherefooter{https://en.wikipedia.org/wiki/Motor_goal}

\subsubsection{Knowledge}

% Source: robot/part/mind/knowledge/knowledge.tex

\url{https://en.wikipedia.org/wiki/Procedural_knowledge}
    % 
    %

\subsubsection{Ontology}

% Source: robot/part/mind/ontology/ontology.tex

Ontology is the phylosophical study of being(Wikipedia)
        
        Luciano Floridi says:
        Information is the flow of well-formed, meaningful data between ontological entities.
        
        % 
        %

\subsubsection{Phylosophy}

% Source: robot/part/mind/phylosophy/phylosophy.tex

\url{https://en.wikipedia.org/wiki/Philosophy_of_mind}
    
    % 
    %

% Source: robot/part/mind/plan/action_plan.tex

An action plan is a detailed plan outlining actions needed to reach one or more goals. \seeherefooter{https://en.wikipedia.org/wiki/Action_plan}
    %
    %

\subsubsection{Psyche}

% Source: robot/part/mind/psyche/psyche.tex

Talks about the quality or the condition under which a part of mind is activated

\paragraph{Kind}

\subparagraph{Consciousness}

% Source: robot/part/mind/psyche/kind/consciousness/consciousness.tex

It is a quality or the condition under which a part of mind is activated
    \url{https://en.wikipedia.org/wiki/Consciousness}

\ProjectDeepHeading{1}{Artificial}

% Source: robot/part/mind/psyche/kind/consciousness/artificial/artificial.tex

\url{https://en.wikipedia.org/wiki/Artificial_consciousness}

\ProjectDeepHeading{1}{Awareness}

% Source: robot/part/mind/psyche/kind/consciousness/awareness/awareness.tex

\url{https://en.wikipedia.org/wiki/Awareness}
    % 
    %

% Source: robot/part/mind/psyche/kind/consciousness/awareness/awarness_vs_conciousness.tex

\ProjectDeepHeading{4}{Terminology and Scope}
        
        In everyday English, \emph{awareness} and \emph{consciousness} are often used interchangeably, but in philosophy of mind, cognitive science, and neuroscience they typically refer to different levels of description. This section fixes working definitions and clarifies how the terms relate.
    
    
    \ProjectDeepHeading{4}{Awareness}
        
        \emph{Awareness} refers to being aware \emph{of something}. It is usually object-directed: one is aware of a sound, a pain, a visual object, a thought, or a change in the environment. In many scientific contexts, awareness is treated as a capacity for detection, monitoring, or access to information, and it can be discussed without committing to any specific theory about subjective experience.
    
    
    \ProjectDeepHeading{4}{Consciousness}
        
        \emph{Consciousness} is the broader notion. At a minimum, it can mean being in a conscious state (as opposed to being under deep anesthesia or in a coma). In many theoretical discussions, consciousness also includes the presence of \emph{subjective experience} (often described as there being something it is like to be in that state). Some usages further distinguish forms such as phenomenal consciousness (experience), access consciousness (availability of information for report and control), and self-consciousness (awareness of oneself as oneself).
    
    
    \ProjectDeepHeading{4}{Relationship Between the Terms}
        
        A common working relationship is:
        \begin{itemize}
            \item Awareness is typically treated as a component or manifestation within a conscious state.
            \item Consciousness is not reducible to awareness alone, because it can include the global condition of being conscious and, in many accounts, the qualitative character of experience.
        \end{itemize}
        
        Illustrative cases:
        \begin{itemize}
            \item Deep anesthesia: neither consciousness nor awareness is present.
            \item Dreaming sleep: consciousness (as experience) may be present, while awareness of the external environment is reduced or absent.
            \item Normal wakefulness: both consciousness and multiple forms of awareness (of environment, body, thoughts) can be present.
        \end{itemize}
    
    
    \ProjectDeepHeading{4}{A Compact Distinction}
        
        A concise way to separate the ideas is:
        \begin{itemize}
            \item Awareness: access or sensitivity to particular contents (being aware \emph{of} X).
            \item Consciousness: the overall condition in which experience occurs (and, depending on the theory, the subjective character of that experience).
        \end{itemize}
    % 
    %

\ProjectDeepHeading{1}{Kind}

% Source: robot/part/mind/psyche/kind/consciousness/kind/autonoetic.tex

Autonoetic consciousness is the human ability to mentally place oneself in the past and future (i.e. mental time travel) or in counterfactual situations (i.e. alternative outcomes), and to thus be able to examine one's own thoughts.
    \url{https://en.wikipedia.org/wiki/Autonoetic_consciousness}
    % 
    %

\ProjectDeepHeading{1}{Thoery}

% Source: robot/part/mind/psyche/kind/consciousness/thoery/theory.tex

\url{https://en.wikipedia.org/wiki/Models_of_consciousness} 
    \\ 
    and
    \\
    \url{https://en.wikipedia.org/wiki/Neural_correlates_of_consciousness}

\ProjectDeepHeading{2}{Kind}

% Source: robot/part/mind/psyche/kind/consciousness/thoery/kind/grobal_work_space.tex

\url{https://en.wikipedia.org/wiki/Global_workspace_theory}

% Source: robot/part/mind/psyche/kind/consciousness/thoery/kind/higher_order.tex

\url{https://en.wikipedia.org/wiki/Higher-order_theories_of_consciousness}

% Source: robot/part/mind/psyche/kind/consciousness/thoery/kind/integrated_information_theory.tex

\url{https://en.wikipedia.org/wiki/Integrated_information_theory}

\ProjectDeepHeading{3}{Damasio}

% Source: robot/part/mind/psyche/kind/consciousness/thoery/kind/damasio/damasio.tex

in \cite{damasio-1999-the-feeling-of-what-happens} proposes that consciousness arises from the interactions between the brain, the body, and the environment. According to this theory, consciousness is not a unitary experience, but rather emerges from the dynamic interplay between different brain regions and their corresponding bodily states.

\ProjectDeepHeading{4}{Proto}

% Source: robot/part/mind/psyche/kind/consciousness/thoery/kind/damasio/proto/proto.tex

According to Damasio's theory of consciousness, the protoself is the first stage in the hierarchical process of consciousness generation. Shared by many species, the protoself is the most basic representation of the organism, and it arises from the brain's constant interaction with the body. The protoself is an unconscious process that creates a "map" of the body's physiological state, which is then used by the brain to generate conscious experience. This "map" is constantly updated as the brain receives new stimuli from the body, and it forms the foundation for the development of more complex forms of consciousness.
    \url{https://en.wikipedia.org/wiki/Damasio%27s_theory_of_consciousness#Protoself}
    
    % 
    %

% Source: robot/part/mind/psyche/kind/consciousness/thoery/kind/damasio/proto/autobiographical_self.tex

\begin{verbatim}
"""
Auto biographical self
"""
\end{verbatim}

% Source: robot/part/mind/psyche/kind/consciousness/thoery/kind/damasio/proto/core.tex

"""core self"""

% Source: robot/part/mind/psyche/kind/consciousness/thoery/kind/damasio/proto/hierarchy.tex

\begin{verbatim}
class Hierarchy:
    def __init__(self):
        hierarchy = ["proto self", "core self", "autobiographical"]
\end{verbatim}

\ProjectDeepHeading{3}{Global workspace thoery}

% Source: robot/part/mind/psyche/kind/consciousness/thoery/kind/global_workspace_thoery/global_workspace_theory.tex

\clickurlfooter{https://en.wikipedia.org/wiki/Global_workspace_theory} argues that the brain contains many specialized processes or modules that operate in parallel, much of which is unconscious. The global workspace is a functional hub of broadcast and integration that allows information to be disseminated across modules. As such GWT can be classified as a functionalist theory of consciousness. When sensory input, memories, or internal representations receive attention, they enter the global workspace and become accessible to various cognitive processes. As elements compete for attention, those that succeed gain entry to the global workspace, allowing their information to be distributed and coordinated throughout the whole cognitive system. GWT resembles the concept of working memory and is proposed to correspond to a 'momentarily active, subjectively experienced' event in working memory. It facilitates top-down control of attention, working memory, planning, and problem-solving through this information sharing.
    %
    %

\ProjectDeepHeading{3}{Higher order theories of consciousness}

% Source: robot/part/mind/psyche/kind/consciousness/thoery/kind/higher-order-theories-of-consciousness/higher-order-theories-of-consciousness.tex

\url{https://www.researchgate.net/publication/288652264_Higher-Order_Theories_of_Consciousness}
    %
    %

\ProjectDeepHeading{3}{Recurrent processing theory}

% Source: robot/part/mind/psyche/kind/consciousness/thoery/kind/recurrent_processing_theory/recurrent_processing_theory.tex

\url{https://selfawarepatterns.com/2020/01/25/recurrent-processing-theory-and-the-function-of-consciousness/}
    %
    %

\subparagraph{Unconsciousness}

% Source: robot/part/mind/psyche/kind/unconsciousness/unconsciousness.tex

\url{https://en.wikipedia.org/wiki/Unconscious_mind}
    
    for unconcious mind see
    \url{https://en.wikipedia.org/wiki/Unconscious_mind}
    % 
    %

\subsubsection{State}

% Source: robot/part/mind/state/state.tex

A mental state, or a mental property, is a state of mind of a person. Mental states comprise a diverse class, including perception, pain/pleasure experience, belief, desire, intention, emotion, and memory.
    \url{https://en.wikipedia.org/wiki/Mental_state}
    
    
    %
    %

\paragraph{Parts}

\subparagraph{Belief}

% Source: robot/part/mind/state/parts/belief/belief.tex

%
    %

\subparagraph{Desire}

% Source: robot/part/mind/state/parts/desire/desire.tex

%
    %

\subparagraph{Emotion}

% Source: robot/part/mind/state/parts/emotion/emotion.tex

\seeherefooter{https://en.wikipedia.org/wiki/Emotion}
    %
    %

\subparagraph{Intension}

% Source: robot/part/mind/state/parts/intension/intention.tex

An intention is a mental state in which a person commits themselves to a course of action. \seeherefooter{https://en.wikipedia.org/wiki/Intention}
    
    \ProjectDeepHeading{3}{Theories}
    %
    %

% Source: robot/part/mind/state/parts/intension/theories.tex

\seeherefooter{https://en.wikipedia.org/wiki/Intention\#Theories_of_intention}
    %
    %

\ProjectDeepHeading{1}{Intentionality}

% Source: robot/part/mind/state/parts/intension/intentionality/intentionality.tex

\url{https://en.wikipedia.org/wiki/Intentionality}
    % 
    %

\subparagraph{Pain pleasure}

% Source: robot/part/mind/state/parts/pain_pleasure/pain_pleasure.tex

\ProjectDeepHeading{3}{Pleasure}
        Pleasure is an experience that feels good, that involves the enjoyment of something.
        \seeherefooter{https://en.wikipedia.org/wiki/Pleasure}
    
    
    \ProjectDeepHeading{3}{Pain}
        Pain is a distressing sensation often caused by intense or damaging stimuli.
        \seeherefooter{https://en.wikipedia.org/wiki/Pain}
    %
    %

\subsubsection{Theory}

% Source: robot/part/mind/theory/theories.tex

\ProjectDeepHeading{1}{Major Thinkers Beyond Claude Shannon in Information Theory}
    
    \ProjectDeepHeading{2}{1. Norbert Wiener (Cybernetics)}
        \textbf{Viewpoint:} Information as the basis of order and control.
        In \textit{Cybernetics} (1948), Wiener wrote that
        \emph{"Information is the measure of organization."}
        He treated information not merely as a mathematical quantity, but as a means to maintain stability in both living and mechanical systems.
        Wiener linked information to physical entropy: information flow corresponds to the preservation of order (negative entropy).
    
    \ProjectDeepHeading{2}{2. Rolf Landauer (Physical Information)}
        \textbf{Viewpoint:} Information is physical and has energetic cost.
        Landauer formulated the principle that \emph{"Information is physical"} and proved in 1961 that erasing one bit of information dissipates at least $kT \ln 2$ energy.
        This connected information theory to thermodynamics and the foundations of physics.
    
    \ProjectDeepHeading{2}{3. Edward Jaynes (Maximum Entropy Principle)}
        \textbf{Viewpoint:} Information as a rule for rational inference.
        Jaynes introduced the \textit{Maximum Entropy Principle}, stating that when only partial information is known,
        the most unbiased probability distribution is the one with maximum entropy.
        Information thus defines the rational boundary between what is known and unknown.
    
    \ProjectDeepHeading{2}{4. Kolmogorov, Chaitin, and Solomonoff (Algorithmic Information Theory)}
        \textbf{Viewpoint:} Information as the length of the shortest possible description.
        They developed the concept of \textit{Kolmogorov Complexity},
        where the information content of a string is defined as the length of the shortest program that can generate it.
        Hence, regular patterns have low information; random data have high information.
    
    \ProjectDeepHeading{2}{5. Luciano Floridi (Philosophy of Information)}
        \textbf{Viewpoint:} Information as a fundamental constituent of reality.
        Floridi, the founder of the \textit{Philosophy of Information}, argues that information is not just epistemic but ontological —
        a basic form of being in the digital age.
        He distinguishes four levels:
        \begin{itemize}
            \item Data
            \item Semantic Information
            \item Knowledge
            \item Wisdom
        \end{itemize}
    
    \ProjectDeepHeading{2}{6. Karl Friston (Free Energy Principle)}
        \textbf{Viewpoint:} Information as prediction and uncertainty reduction in biological systems.
        Friston’s \textit{Free Energy Principle} builds on Shannon and statistical physics:
        the brain minimizes \emph{informational surprise} by continuously updating its internal model of the world.
        In this view, life itself can be described as the perpetual minimization of informational entropy.
    
    \ProjectDeepHeading{2}{7. Other Notable Figures}
        \begin{itemize}
            \item \textbf{Thomas Cover} and \textbf{Joy Thomas}: Authors of the canonical text \textit{Elements of Information Theory}.
            \item \textbf{Claude Berrou}: Developer of turbo codes and the concept of ``soft information.''
            \item \textbf{John Archibald Wheeler}: Proposed the philosophical idea \textit{``It from bit''} — that physical reality arises from informational distinctions.
        \end{itemize}
    
    \ProjectDeepHeading{2}{Summary Table}
        
        \begin{center}
            \begin{tabular}{|l|l|p{8cm}|}
                \hline
                \textbf{Name}        & \textbf{Domain}      & \textbf{View of Information}                                   \\
                \hline
                Claude Shannon       & Communication theory & Reduction of uncertainty about the message source              \\
                Norbert Wiener       & Cybernetics          & Medium of control and order in systems                         \\
                Rolf Landauer        & Physics              & Physical entity with energetic cost                            \\
                Edward Jaynes        & Statistics           & Principle of rational inference and uncertainty quantification \\
                Kolmogorov / Chaitin & Computation theory   & Shortest description length (algorithmic complexity)           \\
                Luciano Floridi      & Philosophy           & Foundational structure of reality and knowledge                \\
                Karl Friston         & Neuroscience         & Minimization of informational surprise (predictive brain)      \\
                \hline
            \end{tabular}
        \end{center}
        
        \bigskip
        \textit{In summary:}
        Shannon defined information as the \emph{reduction of uncertainty about the source of a message},
        while subsequent thinkers expanded it into physical, statistical, computational, philosophical, and cognitive domains.

\paragraph{Computational theory of mind}

% Source: robot/part/mind/theory/computational_theory_of_mind/computational_theory_of_mind.tex

\url{https://en.wikipedia.org/wiki/Computational_theory_of_mind}

\paragraph{Connectionism}

% Source: robot/part/mind/theory/connectionism/connectionism.tex

Connectionism is an approach to the study of human mental processes and cognition that utilizes mathematical models known as connectionist networks or artificial neural networks \seeherefooter{https://en.wikipedia.org/wiki/Connectionism}

\paragraph{Functionalism}

% Source: robot/part/mind/theory/functionalism/functionalism.tex

\url{https://en.wikipedia.org/wiki/Functionalism_(philosophy_of_mind)}

\paragraph{Information processing}

% Source: robot/part/mind/theory/information_processing/information_processing.tex

information processing is an approach to the goal of understanding human thinking that treats cognition as essentially computational in nature, with the mind being the software and the brain being the hardware\url{https://en.wikipedia.org/wiki/Information_processing_(psychology)}
    % 
    %

\paragraph{Information processing thoery}

% Source: robot/part/mind/theory/information_processing_thoery/information_processing_theory.tex

Information processing theory is the approach to the study of cognitive development evolved out of the American experimental tradition in psychology.
    % 
    %

\paragraph{Knowledge representation}

% Source: robot/part/mind/theory/knowledge_representation/knowledge_representation.tex

%
    %

\paragraph{Physicalism}

% Source: robot/part/mind/theory/physicalism/physicalism.tex

\url{https://en.wikipedia.org/wiki/Type_physicalism}
    % 
    %

\chapter{Society}

% Source: society/society.tex

\url{https://en.wikipedia.org/wiki/Multi-agent_system}
    
    my idea to test:\\
    Each robot has a ``Markov blanket`` if multiple robots can keep their ``Markobv Blankets`` overlaped then we have a multi robot system.
    % 
    % 




    \subsubsection{Multi agent systems}
        A multi-agent system (MAS) or "self-organized system" is a computational system composed of multiple interacting intelligent agents
        \seeherefooter{https://en.wikipedia.org/wiki/Multi-agent_system}

% Source: society/blackboard.tex

\url{https://en.wikipedia.org/wiki/Blackboard_system}
    % 
    %

% Source: society/social_norm.tex

\clickurlfooter{https://en.wikipedia.org/wiki/Social_norm} norm is a shared standard of acceptable behavior by a group.
    %
    %

% Source: society/sociality.tex

Sociality is the degree to which individuals in an animal population tend to associate in social groups (for which, the desire or inclination is known as gregariousness) and form cooperative societies. \seeherefooter{https://en.wikipedia.org/wiki/Sociality}
    %
    %

\section{Collective}

% Source: society/collective/behavior.tex

Collective behavior constitutes social processes and events which do not reflect existing social structure (laws, conventions, and institutions), but which emerge in a "spontaneous" way. \seeherefooter{https://en.wikipedia.org/wiki/Collective_behavior}
    
    
    \paragraph{Colective animal behaviour}
        Collective animal behaviour is a form of social behavior involving the coordinated behavior of large groups of similar animals as well as emergent properties of these groups.
        \seeherefooter{https://en.wikipedia.org/wiki/Collective_animal_behavior}
    %
    %

\section{Communication}

% Source: society/communication/communication.tex

\url{https://en.wikipedia.org/wiki/A_Mathematical_Theory_of_Communication}
    % 
    %

\subsection{Business}

\subsubsection{Comapny}

\paragraph{Kinds}

% Source: society/communication/business/comapny/kinds/gmbh.tex

% 
    %

% Source: society/communication/business/comapny/kinds/startup.tex

% 
    %

\section{Part}

\subsection{Mind}

% Source: society/part/mind/collective_self_awareness.tex

\url{https://en.wikipedia.org/wiki/Self-awareness\#Collective_self-awareness}
    % 
    %

\section{Self}

% Source: society/self/self.tex

\url{https://en.wikipedia.org/wiki/Neural_basis_of_self}
    
    also \cite{friston-2013-life-as-we-know-it} mentiones the idea of using ``Markov blankets`` as a model to distinguish between self and the environment. The same idea kan be applied to make a distinguish between self and other robots.

\chapter{System}

% Source: system/system.tex

A system is a group of interacting or interrelated elements that act according to a set of rules to form a unified whole. \seeherefooter{https://en.wikipedia.org/wiki/System}
    \subsubsection{System theory}
    \seeherefooter{https://en.wikipedia.org/wiki/Systems_theory}
    
    %
    %

\section{Dynamical systems}

% Source: system/dynamical_systems/dynamic_systems.tex

a dynamical system is the description of how a system evolves in time.
    \seeherefooter{https://en.wikipedia.org/wiki/Dynamical_system}
    
    %
    %

\bibliographystyle{plainnat}
    \bibliography{/home/donkarlo/Dropbox/repo/nd_shared_research/bibliography/bibliography.bib}

\end{document}
